\chapter{Core Concepts}

% === Source file: concepts/agent-loop.md ===
\section{Agent Loop (OpenClaw)}


An agentic loop is the full “real” run of an agent: intake → context assembly → model inference →
tool execution → streaming replies → persistence. It’s the authoritative path that turns a message
into actions and a final reply, while keeping session state consistent.

In OpenClaw, a loop is a single, serialized run per session that emits lifecycle and stream events
as the model thinks, calls tools, and streams output. This doc explains how that authentic loop is
wired end-to-end.

\subsection{Entry points}


\begin{itemize}
  \item Gateway RPC: \texttt{agent} and \texttt{agent.wait}.
  \item CLI: \texttt{agent} command.
\end{itemize}


\subsection{How it works (high-level)}


\begin{enumerate}
  \item \texttt{agent} RPC validates params, resolves session (sessionKey/sessionId), persists session metadata, returns \texttt{\{ runId, acceptedAt \}} immediately.
  \item \texttt{agentCommand} runs the agent:
\end{enumerate}

\begin{itemize}
  \item resolves model + thinking/verbose defaults
  \item loads skills snapshot
  \item calls \texttt{runEmbeddedPiAgent} (pi-agent-core runtime)
  \item emits \textbf{lifecycle end/error} if the embedded loop does not emit one
\end{itemize}

\begin{enumerate}
  \item \texttt{runEmbeddedPiAgent}:
\end{enumerate}

\begin{itemize}
  \item serializes runs via per-session + global queues
  \item resolves model + auth profile and builds the pi session
  \item subscribes to pi events and streams assistant/tool deltas
  \item enforces timeout -> aborts run if exceeded
  \item returns payloads + usage metadata
\end{itemize}

\begin{enumerate}
  \item \texttt{subscribeEmbeddedPiSession} bridges pi-agent-core events to OpenClaw \texttt{agent} stream:
\end{enumerate}

\begin{itemize}
  \item tool events => \texttt{stream: "tool"}
  \item assistant deltas => \texttt{stream: "assistant"}
  \item lifecycle events => \texttt{stream: "lifecycle"} (\texttt{phase: "start" | "end" | "error"})
\end{itemize}

\begin{enumerate}
  \item \texttt{agent.wait} uses \texttt{waitForAgentJob}:
\end{enumerate}

\begin{itemize}
  \item waits for \textbf{lifecycle end/error} for \texttt{runId}
  \item returns \texttt{\{ status: ok|error|timeout, startedAt, endedAt, error? \}}
\end{itemize}


\subsection{Queueing + concurrency}


\begin{itemize}
  \item Runs are serialized per session key (session lane) and optionally through a global lane.
  \item This prevents tool/session races and keeps session history consistent.
  \item Messaging channels can choose queue modes (collect/steer/followup) that feed this lane system.
\end{itemize}

See \href{/concepts/queue}{Command Queue}.

\subsection{Session + workspace preparation}


\begin{itemize}
  \item Workspace is resolved and created; sandboxed runs may redirect to a sandbox workspace root.
  \item Skills are loaded (or reused from a snapshot) and injected into env and prompt.
  \item Bootstrap/context files are resolved and injected into the system prompt report.
  \item A session write lock is acquired; \texttt{SessionManager} is opened and prepared before streaming.
\end{itemize}


\subsection{Prompt assembly + system prompt}


\begin{itemize}
  \item System prompt is built from OpenClaw’s base prompt, skills prompt, bootstrap context, and per-run overrides.
  \item Model-specific limits and compaction reserve tokens are enforced.
  \item See \href{/concepts/system-prompt}{System prompt} for what the model sees.
\end{itemize}


\subsection{Hook points (where you can intercept)}


OpenClaw has two hook systems:

\begin{itemize}
  \item \textbf{Internal hooks} (Gateway hooks): event-driven scripts for commands and lifecycle events.
  \item \textbf{Plugin hooks}: extension points inside the agent/tool lifecycle and gateway pipeline.
\end{itemize}


\subsubsection{Internal hooks (Gateway hooks)}


\begin{itemize}
  \item \textbf{\texttt{agent:bootstrap}}: runs while building bootstrap files before the system prompt is finalized.
\end{itemize}

Use this to add/remove bootstrap context files.
\begin{itemize}
  \item \textbf{Command hooks}: \texttt{/new}, \texttt{/reset}, \texttt{/stop}, and other command events (see Hooks doc).
\end{itemize}


See \href{/automation/hooks}{Hooks} for setup and examples.

\subsubsection{Plugin hooks (agent + gateway lifecycle)}


These run inside the agent loop or gateway pipeline:

\begin{itemize}
  \item \textbf{\texttt{before\_model\_resolve}}: runs pre-session (no \texttt{messages}) to deterministically override provider/model before model resolution.
  \item \textbf{\texttt{before\_prompt\_build}}: runs after session load (with \texttt{messages}) to inject \texttt{prependContext}/\texttt{systemPrompt} before prompt submission.
  \item \textbf{\texttt{before\_agent\_start}}: legacy compatibility hook that may run in either phase; prefer the explicit hooks above.
  \item \textbf{\texttt{agent\_end}}: inspect the final message list and run metadata after completion.
  \item \textbf{\texttt{before\_compaction} / \texttt{after\_compaction}}: observe or annotate compaction cycles.
  \item \textbf{\texttt{before\_tool\_call} / \texttt{after\_tool\_call}}: intercept tool params/results.
  \item \textbf{\texttt{tool\_result\_persist}}: synchronously transform tool results before they are written to the session transcript.
  \item \textbf{\texttt{message\_received} / \texttt{message\_sending} / \texttt{message\_sent}}: inbound + outbound message hooks.
  \item \textbf{\texttt{session\_start} / \texttt{session\_end}}: session lifecycle boundaries.
  \item \textbf{\texttt{gateway\_start} / \texttt{gateway\_stop}}: gateway lifecycle events.
\end{itemize}


See \href{/tools/plugin\#plugin-hooks}{Plugins} for the hook API and registration details.

\subsection{Streaming + partial replies}


\begin{itemize}
  \item Assistant deltas are streamed from pi-agent-core and emitted as \texttt{assistant} events.
  \item Block streaming can emit partial replies either on \texttt{text\_end} or \texttt{message\_end}.
  \item Reasoning streaming can be emitted as a separate stream or as block replies.
  \item See \href{/concepts/streaming}{Streaming} for chunking and block reply behavior.
\end{itemize}


\subsection{Tool execution + messaging tools}


\begin{itemize}
  \item Tool start/update/end events are emitted on the \texttt{tool} stream.
  \item Tool results are sanitized for size and image payloads before logging/emitting.
  \item Messaging tool sends are tracked to suppress duplicate assistant confirmations.
\end{itemize}


\subsection{Reply shaping + suppression}


\begin{itemize}
  \item Final payloads are assembled from:
  \item assistant text (and optional reasoning)
  \item inline tool summaries (when verbose + allowed)
  \item assistant error text when the model errors
  \item \texttt{NO\_REPLY} is treated as a silent token and filtered from outgoing payloads.
  \item Messaging tool duplicates are removed from the final payload list.
  \item If no renderable payloads remain and a tool errored, a fallback tool error reply is emitted
\end{itemize}

(unless a messaging tool already sent a user-visible reply).

\subsection{Compaction + retries}


\begin{itemize}
  \item Auto-compaction emits \texttt{compaction} stream events and can trigger a retry.
  \item On retry, in-memory buffers and tool summaries are reset to avoid duplicate output.
  \item See \href{/concepts/compaction}{Compaction} for the compaction pipeline.
\end{itemize}


\subsection{Event streams (today)}


\begin{itemize}
  \item \texttt{lifecycle}: emitted by \texttt{subscribeEmbeddedPiSession} (and as a fallback by \texttt{agentCommand})
  \item \texttt{assistant}: streamed deltas from pi-agent-core
  \item \texttt{tool}: streamed tool events from pi-agent-core
\end{itemize}


\subsection{Chat channel handling}


\begin{itemize}
  \item Assistant deltas are buffered into chat \texttt{delta} messages.
  \item A chat \texttt{final} is emitted on \textbf{lifecycle end/error}.
\end{itemize}


\subsection{Timeouts}


\begin{itemize}
  \item \texttt{agent.wait} default: 30s (just the wait). \texttt{timeoutMs} param overrides.
  \item Agent runtime: \texttt{agents.defaults.timeoutSeconds} default 600s; enforced in \texttt{runEmbeddedPiAgent} abort timer.
\end{itemize}


\subsection{Where things can end early}


\begin{itemize}
  \item Agent timeout (abort)
  \item AbortSignal (cancel)
  \item Gateway disconnect or RPC timeout
  \item \texttt{agent.wait} timeout (wait-only, does not stop agent)
\end{itemize}



\clearpage

% === Source file: concepts/agent-workspace.md ===
\section{Agent workspace}


The workspace is the agent's home. It is the only working directory used for
file tools and for workspace context. Keep it private and treat it as memory.

This is separate from \texttt{\textasciitilde{}/.openclaw/}, which stores config, credentials, and
sessions.

\textbf{Important:} the workspace is the \textbf{default cwd}, not a hard sandbox. Tools
resolve relative paths against the workspace, but absolute paths can still reach
elsewhere on the host unless sandboxing is enabled. If you need isolation, use
\href{/gateway/sandboxing}{\texttt{agents.defaults.sandbox}} (and/or per‑agent sandbox config).
When sandboxing is enabled and \texttt{workspaceAccess} is not \texttt{"rw"}, tools operate
inside a sandbox workspace under \texttt{\textasciitilde{}/.openclaw/sandboxes}, not your host workspace.

\subsection{Default location}


\begin{itemize}
  \item Default: \texttt{\textasciitilde{}/.openclaw/workspace}
  \item If \texttt{OPENCLAW\_PROFILE} is set and not \texttt{"default"}, the default becomes
\end{itemize}

\texttt{\textasciitilde{}/.openclaw/workspace-}.
\begin{itemize}
  \item Override in \texttt{\textasciitilde{}/.openclaw/openclaw.json}:
\end{itemize}


\begin{lstlisting}[language=json]
{
  agent: {
    workspace: "~/.openclaw/workspace",
  },
}
\end{lstlisting}


\texttt{openclaw onboard}, \texttt{openclaw configure}, or \texttt{openclaw setup} will create the
workspace and seed the bootstrap files if they are missing.
Sandbox seed copies only accept regular in-workspace files; symlink/hardlink
aliases that resolve outside the source workspace are ignored.

If you already manage the workspace files yourself, you can disable bootstrap
file creation:

\begin{lstlisting}[language=json]
{ agent: { skipBootstrap: true } }
\end{lstlisting}


\subsection{Extra workspace folders}


Older installs may have created \texttt{\textasciitilde{}/openclaw}. Keeping multiple workspace
directories around can cause confusing auth or state drift, because only one
workspace is active at a time.

\textbf{Recommendation:} keep a single active workspace. If you no longer use the
extra folders, archive or move them to Trash (for example \texttt{trash \textasciitilde{}/openclaw}).
If you intentionally keep multiple workspaces, make sure
\texttt{agents.defaults.workspace} points to the active one.

\texttt{openclaw doctor} warns when it detects extra workspace directories.

\subsection{Workspace file map (what each file means)}


These are the standard files OpenClaw expects inside the workspace:

\begin{itemize}
  \item \texttt{AGENTS.md}
  \item Operating instructions for the agent and how it should use memory.
  \item Loaded at the start of every session.
  \item Good place for rules, priorities, and "how to behave" details.
\end{itemize}


\begin{itemize}
  \item \texttt{SOUL.md}
  \item Persona, tone, and boundaries.
  \item Loaded every session.
\end{itemize}


\begin{itemize}
  \item \texttt{USER.md}
  \item Who the user is and how to address them.
  \item Loaded every session.
\end{itemize}


\begin{itemize}
  \item \texttt{IDENTITY.md}
  \item The agent's name, vibe, and emoji.
  \item Created/updated during the bootstrap ritual.
\end{itemize}


\begin{itemize}
  \item \texttt{TOOLS.md}
  \item Notes about your local tools and conventions.
  \item Does not control tool availability; it is only guidance.
\end{itemize}


\begin{itemize}
  \item \texttt{HEARTBEAT.md}
  \item Optional tiny checklist for heartbeat runs.
  \item Keep it short to avoid token burn.
\end{itemize}


\begin{itemize}
  \item \texttt{BOOT.md}
  \item Optional startup checklist executed on gateway restart when internal hooks are enabled.
  \item Keep it short; use the message tool for outbound sends.
\end{itemize}


\begin{itemize}
  \item \texttt{BOOTSTRAP.md}
  \item One-time first-run ritual.
  \item Only created for a brand-new workspace.
  \item Delete it after the ritual is complete.
\end{itemize}


\begin{itemize}
  \item \texttt{memory/YYYY-MM-DD.md}
  \item Daily memory log (one file per day).
  \item Recommended to read today + yesterday on session start.
\end{itemize}


\begin{itemize}
  \item \texttt{MEMORY.md} (optional)
  \item Curated long-term memory.
  \item Only load in the main, private session (not shared/group contexts).
\end{itemize}


See \href{/concepts/memory}{Memory} for the workflow and automatic memory flush.

\begin{itemize}
  \item \texttt{skills/} (optional)
  \item Workspace-specific skills.
  \item Overrides managed/bundled skills when names collide.
\end{itemize}


\begin{itemize}
  \item \texttt{canvas/} (optional)
  \item Canvas UI files for node displays (for example \texttt{canvas/index.html}).
\end{itemize}


If any bootstrap file is missing, OpenClaw injects a "missing file" marker into
the session and continues. Large bootstrap files are truncated when injected;
adjust limits with \texttt{agents.defaults.bootstrapMaxChars} (default: 20000) and
\texttt{agents.defaults.bootstrapTotalMaxChars} (default: 150000).
\texttt{openclaw setup} can recreate missing defaults without overwriting existing
files.

\subsection{What is NOT in the workspace}


These live under \texttt{\textasciitilde{}/.openclaw/} and should NOT be committed to the workspace repo:

\begin{itemize}
  \item \texttt{\textasciitilde{}/.openclaw/openclaw.json} (config)
  \item \texttt{\textasciitilde{}/.openclaw/credentials/} (OAuth tokens, API keys)
  \item \texttt{\textasciitilde{}/.openclaw/agents//sessions/} (session transcripts + metadata)
  \item \texttt{\textasciitilde{}/.openclaw/skills/} (managed skills)
\end{itemize}


If you need to migrate sessions or config, copy them separately and keep them
out of version control.

\subsection{Git backup (recommended, private)}


Treat the workspace as private memory. Put it in a \textbf{private} git repo so it is
backed up and recoverable.

Run these steps on the machine where the Gateway runs (that is where the
workspace lives).

\subsubsection{1) Initialize the repo}


If git is installed, brand-new workspaces are initialized automatically. If this
workspace is not already a repo, run:

\begin{lstlisting}[language=bash]
cd ~/.openclaw/workspace
git init
git add AGENTS.md SOUL.md TOOLS.md IDENTITY.md USER.md HEARTBEAT.md memory/
git commit -m "Add agent workspace"
\end{lstlisting}


\subsubsection{2) Add a private remote (beginner-friendly options)}


Option A: GitHub web UI

\begin{enumerate}
  \item Create a new \textbf{private} repository on GitHub.
  \item Do not initialize with a README (avoids merge conflicts).
  \item Copy the HTTPS remote URL.
  \item Add the remote and push:
\end{enumerate}


\begin{lstlisting}[language=bash]
git branch -M main
git remote add origin <https-url>
git push -u origin main
\end{lstlisting}


Option B: GitHub CLI (\texttt{gh})

\begin{lstlisting}[language=bash]
gh auth login
gh repo create openclaw-workspace --private --source . --remote origin --push
\end{lstlisting}


Option C: GitLab web UI

\begin{enumerate}
  \item Create a new \textbf{private} repository on GitLab.
  \item Do not initialize with a README (avoids merge conflicts).
  \item Copy the HTTPS remote URL.
  \item Add the remote and push:
\end{enumerate}


\begin{lstlisting}[language=bash]
git branch -M main
git remote add origin <https-url>
git push -u origin main
\end{lstlisting}


\subsubsection{3) Ongoing updates}


\begin{lstlisting}[language=bash]
git status
git add .
git commit -m "Update memory"
git push
\end{lstlisting}


\subsection{Do not commit secrets}


Even in a private repo, avoid storing secrets in the workspace:

\begin{itemize}
  \item API keys, OAuth tokens, passwords, or private credentials.
  \item Anything under \texttt{\textasciitilde{}/.openclaw/}.
  \item Raw dumps of chats or sensitive attachments.
\end{itemize}


If you must store sensitive references, use placeholders and keep the real
secret elsewhere (password manager, environment variables, or \texttt{\textasciitilde{}/.openclaw/}).

Suggested \texttt{.gitignore} starter:

\begin{lstlisting}
.DS_Store
.env
**/*.key
**/*.pem
**/secrets*
\end{lstlisting}


\subsection{Moving the workspace to a new machine}


\begin{enumerate}
  \item Clone the repo to the desired path (default \texttt{\textasciitilde{}/.openclaw/workspace}).
  \item Set \texttt{agents.defaults.workspace} to that path in \texttt{\textasciitilde{}/.openclaw/openclaw.json}.
  \item Run \texttt{openclaw setup --workspace } to seed any missing files.
  \item If you need sessions, copy \texttt{\textasciitilde{}/.openclaw/agents//sessions/} from the
\end{enumerate}

old machine separately.

\subsection{Advanced notes}


\begin{itemize}
  \item Multi-agent routing can use different workspaces per agent. See
\end{itemize}

\href{/channels/channel-routing}{Channel routing} for routing configuration.
\begin{itemize}
  \item If \texttt{agents.defaults.sandbox} is enabled, non-main sessions can use per-session sandbox
\end{itemize}

workspaces under \texttt{agents.defaults.sandbox.workspaceRoot}.


\clearpage

% === Source file: concepts/agent.md ===
\section{Agent Runtime}


OpenClaw runs a single embedded agent runtime derived from \textbf{pi-mono}.

\subsection{Workspace (required)}


OpenClaw uses a single agent workspace directory (\texttt{agents.defaults.workspace}) as the agent’s \textbf{only} working directory (\texttt{cwd}) for tools and context.

Recommended: use \texttt{openclaw setup} to create \texttt{\textasciitilde{}/.openclaw/openclaw.json} if missing and initialize the workspace files.

Full workspace layout + backup guide: \href{/concepts/agent-workspace}{Agent workspace}

If \texttt{agents.defaults.sandbox} is enabled, non-main sessions can override this with
per-session workspaces under \texttt{agents.defaults.sandbox.workspaceRoot} (see
\href{/gateway/configuration}{Gateway configuration}).

\subsection{Bootstrap files (injected)}


Inside \texttt{agents.defaults.workspace}, OpenClaw expects these user-editable files:

\begin{itemize}
  \item \texttt{AGENTS.md} — operating instructions + “memory”
  \item \texttt{SOUL.md} — persona, boundaries, tone
  \item \texttt{TOOLS.md} — user-maintained tool notes (e.g. \texttt{imsg}, \texttt{sag}, conventions)
  \item \texttt{BOOTSTRAP.md} — one-time first-run ritual (deleted after completion)
  \item \texttt{IDENTITY.md} — agent name/vibe/emoji
  \item \texttt{USER.md} — user profile + preferred address
\end{itemize}


On the first turn of a new session, OpenClaw injects the contents of these files directly into the agent context.

Blank files are skipped. Large files are trimmed and truncated with a marker so prompts stay lean (read the file for full content).

If a file is missing, OpenClaw injects a single “missing file” marker line (and \texttt{openclaw setup} will create a safe default template).

\texttt{BOOTSTRAP.md} is only created for a \textbf{brand new workspace} (no other bootstrap files present). If you delete it after completing the ritual, it should not be recreated on later restarts.

To disable bootstrap file creation entirely (for pre-seeded workspaces), set:

\begin{lstlisting}[language=json]
{ agent: { skipBootstrap: true } }
\end{lstlisting}


\subsection{Built-in tools}


Core tools (read/exec/edit/write and related system tools) are always available,
subject to tool policy. \texttt{apply\_patch} is optional and gated by
\texttt{tools.exec.applyPatch}. \texttt{TOOLS.md} does \textbf{not} control which tools exist; it’s
guidance for how \_you\_ want them used.

\subsection{Skills}


OpenClaw loads skills from three locations (workspace wins on name conflict):

\begin{itemize}
  \item Bundled (shipped with the install)
  \item Managed/local: \texttt{\textasciitilde{}/.openclaw/skills}
  \item Workspace: \texttt{/skills}
\end{itemize}


Skills can be gated by config/env (see \texttt{skills} in \href{/gateway/configuration}{Gateway configuration}).

\subsection{pi-mono integration}


OpenClaw reuses pieces of the pi-mono codebase (models/tools), but \textbf{session management, discovery, and tool wiring are OpenClaw-owned}.

\begin{itemize}
  \item No pi-coding agent runtime.
  \item No \texttt{\textasciitilde{}/.pi/agent} or \texttt{/.pi} settings are consulted.
\end{itemize}


\subsection{Sessions}


Session transcripts are stored as JSONL at:

\begin{itemize}
  \item \texttt{\textasciitilde{}/.openclaw/agents//sessions/.jsonl}
\end{itemize}


The session ID is stable and chosen by OpenClaw.
Legacy Pi/Tau session folders are \textbf{not} read.

\subsection{Steering while streaming}


When queue mode is \texttt{steer}, inbound messages are injected into the current run.
The queue is checked \textbf{after each tool call}; if a queued message is present,
remaining tool calls from the current assistant message are skipped (error tool
results with "Skipped due to queued user message."), then the queued user
message is injected before the next assistant response.

When queue mode is \texttt{followup} or \texttt{collect}, inbound messages are held until the
current turn ends, then a new agent turn starts with the queued payloads. See
\href{/concepts/queue}{Queue} for mode + debounce/cap behavior.

Block streaming sends completed assistant blocks as soon as they finish; it is
\textbf{off by default} (\texttt{agents.defaults.blockStreamingDefault: "off"}).
Tune the boundary via \texttt{agents.defaults.blockStreamingBreak} (\texttt{text\_end} vs \texttt{message\_end}; defaults to text\_end).
Control soft block chunking with \texttt{agents.defaults.blockStreamingChunk} (defaults to
800–1200 chars; prefers paragraph breaks, then newlines; sentences last).
Coalesce streamed chunks with \texttt{agents.defaults.blockStreamingCoalesce} to reduce
single-line spam (idle-based merging before send). Non-Telegram channels require
explicit \texttt{*.blockStreaming: true} to enable block replies.
Verbose tool summaries are emitted at tool start (no debounce); Control UI
streams tool output via agent events when available.
More details: \href{/concepts/streaming}{Streaming + chunking}.

\subsection{Model refs}


Model refs in config (for example \texttt{agents.defaults.model} and \texttt{agents.defaults.models}) are parsed by splitting on the \textbf{first} \texttt{/}.

\begin{itemize}
  \item Use \texttt{provider/model} when configuring models.
  \item If the model ID itself contains \texttt{/} (OpenRouter-style), include the provider prefix (example: \texttt{openrouter/moonshotai/kimi-k2}).
  \item If you omit the provider, OpenClaw treats the input as an alias or a model for the \textbf{default provider} (only works when there is no \texttt{/} in the model ID).
\end{itemize}


\subsection{Configuration (minimal)}


At minimum, set:

\begin{itemize}
  \item \texttt{agents.defaults.workspace}
  \item \texttt{channels.whatsapp.allowFrom} (strongly recommended)
\end{itemize}


\hrulefill


\_Next: \href{/channels/group-messages}{Group Chats}\_ 🦞


\clearpage

% === Source file: concepts/architecture.md ===
\section{Gateway architecture}


Last updated: 2026-01-22

\subsection{Overview}


\begin{itemize}
  \item A single long‑lived \textbf{Gateway} owns all messaging surfaces (WhatsApp via
\end{itemize}

Baileys, Telegram via grammY, Slack, Discord, Signal, iMessage, WebChat).
\begin{itemize}
  \item Control-plane clients (macOS app, CLI, web UI, automations) connect to the
\end{itemize}

Gateway over \textbf{WebSocket} on the configured bind host (default
\texttt{127.0.0.1:18789}).
\begin{itemize}
  \item \textbf{Nodes} (macOS/iOS/Android/headless) also connect over \textbf{WebSocket}, but
\end{itemize}

declare \texttt{role: node} with explicit caps/commands.
\begin{itemize}
  \item One Gateway per host; it is the only place that opens a WhatsApp session.
  \item The \textbf{canvas host} is served by the Gateway HTTP server under:
  \item \texttt{/\_\_openclaw\_\_/canvas/} (agent-editable HTML/CSS/JS)
  \item \texttt{/\_\_openclaw\_\_/a2ui/} (A2UI host)
\end{itemize}

It uses the same port as the Gateway (default \texttt{18789}).

\subsection{Components and flows}


\subsubsection{Gateway (daemon)}


\begin{itemize}
  \item Maintains provider connections.
  \item Exposes a typed WS API (requests, responses, server‑push events).
  \item Validates inbound frames against JSON Schema.
  \item Emits events like \texttt{agent}, \texttt{chat}, \texttt{presence}, \texttt{health}, \texttt{heartbeat}, \texttt{cron}.
\end{itemize}


\subsubsection{Clients (mac app / CLI / web admin)}


\begin{itemize}
  \item One WS connection per client.
  \item Send requests (\texttt{health}, \texttt{status}, \texttt{send}, \texttt{agent}, \texttt{system-presence}).
  \item Subscribe to events (\texttt{tick}, \texttt{agent}, \texttt{presence}, \texttt{shutdown}).
\end{itemize}


\subsubsection{Nodes (macOS / iOS / Android / headless)}


\begin{itemize}
  \item Connect to the \textbf{same WS server} with \texttt{role: node}.
  \item Provide a device identity in \texttt{connect}; pairing is \textbf{device‑based} (role \texttt{node}) and
\end{itemize}

approval lives in the device pairing store.
\begin{itemize}
  \item Expose commands like \texttt{canvas.\textit{}, \texttt{camera.}}, \texttt{screen.record}, \texttt{location.get}.
\end{itemize}


Protocol details:

\begin{itemize}
  \item \href{/gateway/protocol}{Gateway protocol}
\end{itemize}


\subsubsection{WebChat}


\begin{itemize}
  \item Static UI that uses the Gateway WS API for chat history and sends.
  \item In remote setups, connects through the same SSH/Tailscale tunnel as other
\end{itemize}

clients.

\subsection{Connection lifecycle (single client)}


\begin{verbatim}
% Mermaid diagram
sequenceDiagram
    participant Client
    participant Gateway

    Client->>Gateway: req:connect
    Gateway-->>Client: res (ok)
    Note right of Gateway: or res error + close
    Note left of Client: payload=hello-ok<br>snapshot: presence + health

    Gateway-->>Client: event:presence
    Gateway-->>Client: event:tick

    Client->>Gateway: req:agent
    Gateway-->>Client: res:agent<br>ack {runId, status:"accepted"}
    Gateway-->>Client: event:agent<br>(streaming)
    Gateway-->>Client: res:agent<br>final {runId, status, summary}
\end{verbatim}


\subsection{Wire protocol (summary)}


\begin{itemize}
  \item Transport: WebSocket, text frames with JSON payloads.
  \item First frame \textbf{must} be \texttt{connect}.
  \item After handshake:
  \item Requests: \texttt{\{type:"req", id, method, params\}} → \texttt{\{type:"res", id, ok, payload|error\}}
  \item Events: \texttt{\{type:"event", event, payload, seq?, stateVersion?\}}
  \item If \texttt{OPENCLAW\_GATEWAY\_TOKEN} (or \texttt{--token}) is set, \texttt{connect.params.auth.token}
\end{itemize}

must match or the socket closes.
\begin{itemize}
  \item Idempotency keys are required for side‑effecting methods (\texttt{send}, \texttt{agent}) to
\end{itemize}

safely retry; the server keeps a short‑lived dedupe cache.
\begin{itemize}
  \item Nodes must include \texttt{role: "node"} plus caps/commands/permissions in \texttt{connect}.
\end{itemize}


\subsection{Pairing + local trust}


\begin{itemize}
  \item All WS clients (operators + nodes) include a \textbf{device identity} on \texttt{connect}.
  \item New device IDs require pairing approval; the Gateway issues a \textbf{device token}
\end{itemize}

for subsequent connects.
\begin{itemize}
  \item \textbf{Local} connects (loopback or the gateway host’s own tailnet address) can be
\end{itemize}

auto‑approved to keep same‑host UX smooth.
\begin{itemize}
  \item All connects must sign the \texttt{connect.challenge} nonce.
  \item Signature payload \texttt{v3} also binds \texttt{platform} + \texttt{deviceFamily}; the gateway
\end{itemize}

pins paired metadata on reconnect and requires repair pairing for metadata
changes.
\begin{itemize}
  \item \textbf{Non‑local} connects still require explicit approval.
  \item Gateway auth (\texttt{gateway.auth.*}) still applies to \textbf{all} connections, local or
\end{itemize}

remote.

Details: \href{/gateway/protocol}{Gateway protocol}, \href{/channels/pairing}{Pairing},
\href{/gateway/security}{Security}.

\subsection{Protocol typing and codegen}


\begin{itemize}
  \item TypeBox schemas define the protocol.
  \item JSON Schema is generated from those schemas.
  \item Swift models are generated from the JSON Schema.
\end{itemize}


\subsection{Remote access}


\begin{itemize}
  \item Preferred: Tailscale or VPN.
  \item Alternative: SSH tunnel
\end{itemize}


\begin{lstlisting}[language=bash]
  ssh -N -L 18789:127.0.0.1:18789 user@host
\end{lstlisting}


\begin{itemize}
  \item The same handshake + auth token apply over the tunnel.
  \item TLS + optional pinning can be enabled for WS in remote setups.
\end{itemize}


\subsection{Operations snapshot}


\begin{itemize}
  \item Start: \texttt{openclaw gateway} (foreground, logs to stdout).
  \item Health: \texttt{health} over WS (also included in \texttt{hello-ok}).
  \item Supervision: launchd/systemd for auto‑restart.
\end{itemize}


\subsection{Invariants}


\begin{itemize}
  \item Exactly one Gateway controls a single Baileys session per host.
  \item Handshake is mandatory; any non‑JSON or non‑connect first frame is a hard close.
  \item Events are not replayed; clients must refresh on gaps.
\end{itemize}



\clearpage

% === Source file: concepts/compaction.md ===
\section{Context Window \& Compaction}


Every model has a \textbf{context window} (max tokens it can see). Long-running chats accumulate messages and tool results; once the window is tight, OpenClaw \textbf{compacts} older history to stay within limits.

\subsection{What compaction is}


Compaction \textbf{summarizes older conversation} into a compact summary entry and keeps recent messages intact. The summary is stored in the session history, so future requests use:

\begin{itemize}
  \item The compaction summary
  \item Recent messages after the compaction point
\end{itemize}


Compaction \textbf{persists} in the session’s JSONL history.

\subsection{Configuration}


Use the \texttt{agents.defaults.compaction} setting in your \texttt{openclaw.json} to configure compaction behavior (mode, target tokens, etc.).
Compaction summarization preserves opaque identifiers by default (\texttt{identifierPolicy: "strict"}). You can override this with \texttt{identifierPolicy: "off"} or provide custom text with \texttt{identifierPolicy: "custom"} and \texttt{identifierInstructions}.

\subsection{Auto-compaction (default on)}


When a session nears or exceeds the model’s context window, OpenClaw triggers auto-compaction and may retry the original request using the compacted context.

You’ll see:

\begin{itemize}
  \item \texttt{🧹 Auto-compaction complete} in verbose mode
  \item \texttt{/status} showing \texttt{🧹 Compactions: }
\end{itemize}


Before compaction, OpenClaw can run a \textbf{silent memory flush} turn to store
durable notes to disk. See \href{/concepts/memory}{Memory} for details and config.

\subsection{Manual compaction}


Use \texttt{/compact} (optionally with instructions) to force a compaction pass:

\begin{lstlisting}
/compact Focus on decisions and open questions
\end{lstlisting}


\subsection{Context window source}


Context window is model-specific. OpenClaw uses the model definition from the configured provider catalog to determine limits.

\subsection{Compaction vs pruning}


\begin{itemize}
  \item \textbf{Compaction}: summarises and \textbf{persists} in JSONL.
  \item \textbf{Session pruning}: trims old \textbf{tool results} only, \textbf{in-memory}, per request.
\end{itemize}


See \href{/concepts/session-pruning}{/concepts/session-pruning} for pruning details.

\subsection{OpenAI server-side compaction}


OpenClaw also supports OpenAI Responses server-side compaction hints for
compatible direct OpenAI models. This is separate from local OpenClaw
compaction and can run alongside it.

\begin{itemize}
  \item Local compaction: OpenClaw summarizes and persists into session JSONL.
  \item Server-side compaction: OpenAI compacts context on the provider side when
\end{itemize}

\texttt{store} + \texttt{context\_management} are enabled.

See \href{/providers/openai}{OpenAI provider} for model params and overrides.

\subsection{Tips}


\begin{itemize}
  \item Use \texttt{/compact} when sessions feel stale or context is bloated.
  \item Large tool outputs are already truncated; pruning can further reduce tool-result buildup.
  \item If you need a fresh slate, \texttt{/new} or \texttt{/reset} starts a new session id.
\end{itemize}



\clearpage

% === Source file: concepts/context.md ===
\section{Context}


“Context” is \textbf{everything OpenClaw sends to the model for a run}. It is bounded by the model’s \textbf{context window} (token limit).

Beginner mental model:

\begin{itemize}
  \item \textbf{System prompt} (OpenClaw-built): rules, tools, skills list, time/runtime, and injected workspace files.
  \item \textbf{Conversation history}: your messages + the assistant’s messages for this session.
  \item \textbf{Tool calls/results + attachments}: command output, file reads, images/audio, etc.
\end{itemize}


Context is \_not the same thing\_ as “memory”: memory can be stored on disk and reloaded later; context is what’s inside the model’s current window.

\subsection{Quick start (inspect context)}


\begin{itemize}
  \item \texttt{/status} → quick “how full is my window?” view + session settings.
  \item \texttt{/context list} → what’s injected + rough sizes (per file + totals).
  \item \texttt{/context detail} → deeper breakdown: per-file, per-tool schema sizes, per-skill entry sizes, and system prompt size.
  \item \texttt{/usage tokens} → append per-reply usage footer to normal replies.
  \item \texttt{/compact} → summarize older history into a compact entry to free window space.
\end{itemize}


See also: \href{/tools/slash-commands}{Slash commands}, \href{/reference/token-use}{Token use \& costs}, \href{/concepts/compaction}{Compaction}.

\subsection{Example output}


Values vary by model, provider, tool policy, and what’s in your workspace.

\subsubsection{/context list}


\begin{lstlisting}
🧠 Context breakdown
Workspace: <workspaceDir>
Bootstrap max/file: 20,000 chars
Sandbox: mode=non-main sandboxed=false
System prompt (run): 38,412 chars (~9,603 tok) (Project Context 23,901 chars (~5,976 tok))

Injected workspace files:
- AGENTS.md: OK | raw 1,742 chars (~436 tok) | injected 1,742 chars (~436 tok)
- SOUL.md: OK | raw 912 chars (~228 tok) | injected 912 chars (~228 tok)
- TOOLS.md: TRUNCATED | raw 54,210 chars (~13,553 tok) | injected 20,962 chars (~5,241 tok)
- IDENTITY.md: OK | raw 211 chars (~53 tok) | injected 211 chars (~53 tok)
- USER.md: OK | raw 388 chars (~97 tok) | injected 388 chars (~97 tok)
- HEARTBEAT.md: MISSING | raw 0 | injected 0
- BOOTSTRAP.md: OK | raw 0 chars (~0 tok) | injected 0 chars (~0 tok)

Skills list (system prompt text): 2,184 chars (~546 tok) (12 skills)
Tools: read, edit, write, exec, process, browser, message, sessions_send, …
Tool list (system prompt text): 1,032 chars (~258 tok)
Tool schemas (JSON): 31,988 chars (~7,997 tok) (counts toward context; not shown as text)
Tools: (same as above)

Session tokens (cached): 14,250 total / ctx=32,000
\end{lstlisting}


\subsubsection{/context detail}


\begin{lstlisting}
🧠 Context breakdown (detailed)
…
Top skills (prompt entry size):
- frontend-design: 412 chars (~103 tok)
- oracle: 401 chars (~101 tok)
… (+10 more skills)

Top tools (schema size):
- browser: 9,812 chars (~2,453 tok)
- exec: 6,240 chars (~1,560 tok)
… (+N more tools)
\end{lstlisting}


\subsection{What counts toward the context window}


Everything the model receives counts, including:

\begin{itemize}
  \item System prompt (all sections).
  \item Conversation history.
  \item Tool calls + tool results.
  \item Attachments/transcripts (images/audio/files).
  \item Compaction summaries and pruning artifacts.
  \item Provider “wrappers” or hidden headers (not visible, still counted).
\end{itemize}


\subsection{How OpenClaw builds the system prompt}


The system prompt is \textbf{OpenClaw-owned} and rebuilt each run. It includes:

\begin{itemize}
  \item Tool list + short descriptions.
  \item Skills list (metadata only; see below).
  \item Workspace location.
  \item Time (UTC + converted user time if configured).
  \item Runtime metadata (host/OS/model/thinking).
  \item Injected workspace bootstrap files under \textbf{Project Context}.
\end{itemize}


Full breakdown: \href{/concepts/system-prompt}{System Prompt}.

\subsection{Injected workspace files (Project Context)}


By default, OpenClaw injects a fixed set of workspace files (if present):

\begin{itemize}
  \item \texttt{AGENTS.md}
  \item \texttt{SOUL.md}
  \item \texttt{TOOLS.md}
  \item \texttt{IDENTITY.md}
  \item \texttt{USER.md}
  \item \texttt{HEARTBEAT.md}
  \item \texttt{BOOTSTRAP.md} (first-run only)
\end{itemize}


Large files are truncated per-file using \texttt{agents.defaults.bootstrapMaxChars} (default \texttt{20000} chars). OpenClaw also enforces a total bootstrap injection cap across files with \texttt{agents.defaults.bootstrapTotalMaxChars} (default \texttt{150000} chars). \texttt{/context} shows \textbf{raw vs injected} sizes and whether truncation happened.

\subsection{Skills: what’s injected vs loaded on-demand}


The system prompt includes a compact \textbf{skills list} (name + description + location). This list has real overhead.

Skill instructions are \_not\_ included by default. The model is expected to \texttt{read} the skill’s \texttt{SKILL.md} \textbf{only when needed}.

\subsection{Tools: there are two costs}


Tools affect context in two ways:

\begin{enumerate}
  \item \textbf{Tool list text} in the system prompt (what you see as “Tooling”).
  \item \textbf{Tool schemas} (JSON). These are sent to the model so it can call tools. They count toward context even though you don’t see them as plain text.
\end{enumerate}


\texttt{/context detail} breaks down the biggest tool schemas so you can see what dominates.

\subsection{Commands, directives, and “inline shortcuts”}


Slash commands are handled by the Gateway. There are a few different behaviors:

\begin{itemize}
  \item \textbf{Standalone commands}: a message that is only \texttt{/...} runs as a command.
  \item \textbf{Directives}: \texttt{/think}, \texttt{/verbose}, \texttt{/reasoning}, \texttt{/elevated}, \texttt{/model}, \texttt{/queue} are stripped before the model sees the message.
  \item Directive-only messages persist session settings.
  \item Inline directives in a normal message act as per-message hints.
  \item \textbf{Inline shortcuts} (allowlisted senders only): certain \texttt{/...} tokens inside a normal message can run immediately (example: “hey /status”), and are stripped before the model sees the remaining text.
\end{itemize}


Details: \href{/tools/slash-commands}{Slash commands}.

\subsection{Sessions, compaction, and pruning (what persists)}


What persists across messages depends on the mechanism:

\begin{itemize}
  \item \textbf{Normal history} persists in the session transcript until compacted/pruned by policy.
  \item \textbf{Compaction} persists a summary into the transcript and keeps recent messages intact.
  \item \textbf{Pruning} removes old tool results from the \_in-memory\_ prompt for a run, but does not rewrite the transcript.
\end{itemize}


Docs: \href{/concepts/session}{Session}, \href{/concepts/compaction}{Compaction}, \href{/concepts/session-pruning}{Session pruning}.

\subsection{What /context actually reports}


\texttt{/context} prefers the latest \textbf{run-built} system prompt report when available:

\begin{itemize}
  \item \texttt{System prompt (run)} = captured from the last embedded (tool-capable) run and persisted in the session store.
  \item \texttt{System prompt (estimate)} = computed on the fly when no run report exists (or when running via a CLI backend that doesn’t generate the report).
\end{itemize}


Either way, it reports sizes and top contributors; it does \textbf{not} dump the full system prompt or tool schemas.


\clearpage

\section{Features}
% Source file: concepts/features.md

% === Source file: concepts/features.md ===
\subsection{Highlights}


WhatsApp, Telegram, Discord, and iMessage with a single Gateway.
Add Mattermost and more with extensions.
Multi-agent routing with isolated sessions.
Images, audio, and documents in and out.
Web Control UI and macOS companion app.
iOS and Android nodes with pairing, voice/chat, and rich device commands.

\subsection{Full list}


\begin{itemize}
  \item WhatsApp integration via WhatsApp Web (Baileys)
  \item Telegram bot support (grammY)
  \item Discord bot support (channels.discord.js)
  \item Mattermost bot support (plugin)
  \item iMessage integration via local imsg CLI (macOS)
  \item Agent bridge for Pi in RPC mode with tool streaming
  \item Streaming and chunking for long responses
  \item Multi-agent routing for isolated sessions per workspace or sender
  \item Subscription auth for Anthropic and OpenAI via OAuth
  \item Sessions: direct chats collapse into shared \texttt{main}; groups are isolated
  \item Group chat support with mention based activation
  \item Media support for images, audio, and documents
  \item Optional voice note transcription hook
  \item WebChat and macOS menu bar app
  \item iOS node with pairing, Canvas, camera, screen recording, location, and voice features
  \item Android node with pairing, Connect tab, chat sessions, voice tab, Canvas/camera/screen, plus device, notifications, contacts/calendar, motion, photos, SMS, and app update commands
\end{itemize}


Legacy Claude, Codex, Gemini, and Opencode paths have been removed. Pi is the only
coding agent path.


\clearpage

% === Source file: concepts/markdown-formatting.md ===
\section{Markdown formatting}


OpenClaw formats outbound Markdown by converting it into a shared intermediate
representation (IR) before rendering channel-specific output. The IR keeps the
source text intact while carrying style/link spans so chunking and rendering can
stay consistent across channels.

\subsection{Goals}


\begin{itemize}
  \item \textbf{Consistency:} one parse step, multiple renderers.
  \item \textbf{Safe chunking:} split text before rendering so inline formatting never
\end{itemize}

breaks across chunks.
\begin{itemize}
  \item \textbf{Channel fit:} map the same IR to Slack mrkdwn, Telegram HTML, and Signal
\end{itemize}

style ranges without re-parsing Markdown.

\subsection{Pipeline}


\begin{enumerate}
  \item \textbf{Parse Markdown -> IR}
\end{enumerate}

\begin{itemize}
  \item IR is plain text plus style spans (bold/italic/strike/code/spoiler) and link spans.
  \item Offsets are UTF-16 code units so Signal style ranges align with its API.
  \item Tables are parsed only when a channel opts into table conversion.
\end{itemize}

\begin{enumerate}
  \item \textbf{Chunk IR (format-first)}
\end{enumerate}

\begin{itemize}
  \item Chunking happens on the IR text before rendering.
  \item Inline formatting does not split across chunks; spans are sliced per chunk.
\end{itemize}

\begin{enumerate}
  \item \textbf{Render per channel}
\end{enumerate}

\begin{itemize}
  \item \textbf{Slack:} mrkdwn tokens (bold/italic/strike/code), links as ``.
  \item \textbf{Telegram:} HTML tags (`\texttt{, }\texttt{, }\texttt{, }\texttt{, }\texttt{, }`).
  \item \textbf{Signal:} plain text + \texttt{text-style} ranges; links become \texttt{label (url)} when label differs.
\end{itemize}


\subsection{IR example}


Input Markdown:

\begin{lstlisting}
Hello **world** — see [docs](https://docs.openclaw.ai).
\end{lstlisting}


IR (schematic):

\begin{lstlisting}[language=json]
{
  "text": "Hello world — see docs.",
  "styles": [{ "start": 6, "end": 11, "style": "bold" }],
  "links": [{ "start": 19, "end": 23, "href": "https://docs.openclaw.ai" }]
}
\end{lstlisting}


\subsection{Where it is used}


\begin{itemize}
  \item Slack, Telegram, and Signal outbound adapters render from the IR.
  \item Other channels (WhatsApp, iMessage, MS Teams, Discord) still use plain text or
\end{itemize}

their own formatting rules, with Markdown table conversion applied before
chunking when enabled.

\subsection{Table handling}


Markdown tables are not consistently supported across chat clients. Use
\texttt{markdown.tables} to control conversion per channel (and per account).

\begin{itemize}
  \item \texttt{code}: render tables as code blocks (default for most channels).
  \item \texttt{bullets}: convert each row into bullet points (default for Signal + WhatsApp).
  \item \texttt{off}: disable table parsing and conversion; raw table text passes through.
\end{itemize}


Config keys:

\begin{lstlisting}[language=yaml]
channels:
  discord:
    markdown:
      tables: code
    accounts:
      work:
        markdown:
          tables: off
\end{lstlisting}


\subsection{Chunking rules}


\begin{itemize}
  \item Chunk limits come from channel adapters/config and are applied to the IR text.
  \item Code fences are preserved as a single block with a trailing newline so channels
\end{itemize}

render them correctly.
\begin{itemize}
  \item List prefixes and blockquote prefixes are part of the IR text, so chunking
\end{itemize}

does not split mid-prefix.
\begin{itemize}
  \item Inline styles (bold/italic/strike/inline-code/spoiler) are never split across
\end{itemize}

chunks; the renderer reopens styles inside each chunk.

If you need more on chunking behavior across channels, see
\href{/concepts/streaming}{Streaming + chunking}.

\subsection{Link policy}


\begin{itemize}
  \item \textbf{Slack:} \texttt{\href{url}{label}} -> ``; bare URLs remain bare. Autolink
\end{itemize}

is disabled during parse to avoid double-linking.
\begin{itemize}
  \item \textbf{Telegram:} \texttt{\href{url}{label}} -> \texttt{label} (HTML parse mode).
  \item \textbf{Signal:} \texttt{\href{url}{label}} -> \texttt{label (url)} unless label matches the URL.
\end{itemize}


\subsection{Spoilers}


Spoiler markers (\texttt{||spoiler||}) are parsed only for Signal, where they map to
SPOILER style ranges. Other channels treat them as plain text.

\subsection{How to add or update a channel formatter}


\begin{enumerate}
  \item \textbf{Parse once:} use the shared \texttt{markdownToIR(...)} helper with channel-appropriate
\end{enumerate}

options (autolink, heading style, blockquote prefix).
\begin{enumerate}
  \item \textbf{Render:} implement a renderer with \texttt{renderMarkdownWithMarkers(...)} and a
\end{enumerate}

style marker map (or Signal style ranges).
\begin{enumerate}
  \item \textbf{Chunk:} call \texttt{chunkMarkdownIR(...)} before rendering; render each chunk.
  \item \textbf{Wire adapter:} update the channel outbound adapter to use the new chunker
\end{enumerate}

and renderer.
\begin{enumerate}
  \item \textbf{Test:} add or update format tests and an outbound delivery test if the
\end{enumerate}

channel uses chunking.

\subsection{Common gotchas}


\begin{itemize}
  \item Slack angle-bracket tokens (\texttt{<@U123>}, \texttt{<\#C123>}, ``) must be
\end{itemize}

preserved; escape raw HTML safely.
\begin{itemize}
  \item Telegram HTML requires escaping text outside tags to avoid broken markup.
  \item Signal style ranges depend on UTF-16 offsets; do not use code point offsets.
  \item Preserve trailing newlines for fenced code blocks so closing markers land on
\end{itemize}

their own line.


\clearpage

% === Source file: concepts/memory.md ===
\section{Memory}


OpenClaw memory is \textbf{plain Markdown in the agent workspace}. The files are the
source of truth; the model only "remembers" what gets written to disk.

Memory search tools are provided by the active memory plugin (default:
\texttt{memory-core}). Disable memory plugins with \texttt{plugins.slots.memory = "none"}.

\subsection{Memory files (Markdown)}


The default workspace layout uses two memory layers:

\begin{itemize}
  \item \texttt{memory/YYYY-MM-DD.md}
  \item Daily log (append-only).
  \item Read today + yesterday at session start.
  \item \texttt{MEMORY.md} (optional)
  \item Curated long-term memory.
  \item \textbf{Only load in the main, private session} (never in group contexts).
\end{itemize}


These files live under the workspace (\texttt{agents.defaults.workspace}, default
\texttt{\textasciitilde{}/.openclaw/workspace}). See \href{/concepts/agent-workspace}{Agent workspace} for the full layout.

\subsection{Memory tools}


OpenClaw exposes two agent-facing tools for these Markdown files:

\begin{itemize}
  \item \texttt{memory\_search} — semantic recall over indexed snippets.
  \item \texttt{memory\_get} — targeted read of a specific Markdown file/line range.
\end{itemize}


\texttt{memory\_get} now \textbf{degrades gracefully when a file doesn't exist} (for example,
today's daily log before the first write). Both the builtin manager and the QMD
backend return \texttt{\{ text: "", path \}} instead of throwing \texttt{ENOENT}, so agents can
handle "nothing recorded yet" and continue their workflow without wrapping the
tool call in try/catch logic.

\subsection{When to write memory}


\begin{itemize}
  \item Decisions, preferences, and durable facts go to \texttt{MEMORY.md}.
  \item Day-to-day notes and running context go to \texttt{memory/YYYY-MM-DD.md}.
  \item If someone says "remember this," write it down (do not keep it in RAM).
  \item This area is still evolving. It helps to remind the model to store memories; it will know what to do.
  \item If you want something to stick, \textbf{ask the bot to write it} into memory.
\end{itemize}


\subsection{Automatic memory flush (pre-compaction ping)}


When a session is \textbf{close to auto-compaction}, OpenClaw triggers a **silent,
agentic turn\textbf{ that reminds the model to write durable memory }before** the
context is compacted. The default prompts explicitly say the model \_may reply\_,
but usually \texttt{NO\_REPLY} is the correct response so the user never sees this turn.

This is controlled by \texttt{agents.defaults.compaction.memoryFlush}:

\begin{lstlisting}[language=json]
{
  agents: {
    defaults: {
      compaction: {
        reserveTokensFloor: 20000,
        memoryFlush: {
          enabled: true,
          softThresholdTokens: 4000,
          systemPrompt: "Session nearing compaction. Store durable memories now.",
          prompt: "Write any lasting notes to memory/YYYY-MM-DD.md; reply with NO_REPLY if nothing to store.",
        },
      },
    },
  },
}
\end{lstlisting}


Details:

\begin{itemize}
  \item \textbf{Soft threshold}: flush triggers when the session token estimate crosses
\end{itemize}

\texttt{contextWindow - reserveTokensFloor - softThresholdTokens}.
\begin{itemize}
  \item \textbf{Silent} by default: prompts include \texttt{NO\_REPLY} so nothing is delivered.
  \item \textbf{Two prompts}: a user prompt plus a system prompt append the reminder.
  \item \textbf{One flush per compaction cycle} (tracked in \texttt{sessions.json}).
  \item \textbf{Workspace must be writable}: if the session runs sandboxed with
\end{itemize}

\texttt{workspaceAccess: "ro"} or \texttt{"none"}, the flush is skipped.

For the full compaction lifecycle, see
\href{/reference/session-management-compaction}{Session management + compaction}.

\subsection{Vector memory search}


OpenClaw can build a small vector index over \texttt{MEMORY.md} and \texttt{memory/*.md} so
semantic queries can find related notes even when wording differs.

Defaults:

\begin{itemize}
  \item Enabled by default.
  \item Watches memory files for changes (debounced).
  \item Configure memory search under \texttt{agents.defaults.memorySearch} (not top-level
\end{itemize}

\texttt{memorySearch}).
\begin{itemize}
  \item Uses remote embeddings by default. If \texttt{memorySearch.provider} is not set, OpenClaw auto-selects:
\end{itemize}

\begin{enumerate}
  \item \texttt{local} if a \texttt{memorySearch.local.modelPath} is configured and the file exists.
  \item \texttt{openai} if an OpenAI key can be resolved.
  \item \texttt{gemini} if a Gemini key can be resolved.
  \item \texttt{voyage} if a Voyage key can be resolved.
  \item \texttt{mistral} if a Mistral key can be resolved.
  \item Otherwise memory search stays disabled until configured.
\end{enumerate}

\begin{itemize}
  \item Local mode uses node-llama-cpp and may require \texttt{pnpm approve-builds}.
  \item Uses sqlite-vec (when available) to accelerate vector search inside SQLite.
  \item \texttt{memorySearch.provider = "ollama"} is also supported for local/self-hosted
\end{itemize}

Ollama embeddings (\texttt{/api/embeddings}), but it is not auto-selected.

Remote embeddings \textbf{require} an API key for the embedding provider. OpenClaw
resolves keys from auth profiles, \texttt{models.providers.*.apiKey}, or environment
variables. Codex OAuth only covers chat/completions and does \textbf{not} satisfy
embeddings for memory search. For Gemini, use \texttt{GEMINI\_API\_KEY} or
\texttt{models.providers.google.apiKey}. For Voyage, use \texttt{VOYAGE\_API\_KEY} or
\texttt{models.providers.voyage.apiKey}. For Mistral, use \texttt{MISTRAL\_API\_KEY} or
\texttt{models.providers.mistral.apiKey}. Ollama typically does not require a real API
key (a placeholder like \texttt{OLLAMA\_API\_KEY=ollama-local} is enough when needed by
local policy).
When using a custom OpenAI-compatible endpoint,
set \texttt{memorySearch.remote.apiKey} (and optional \texttt{memorySearch.remote.headers}).

\subsubsection{QMD backend (experimental)}


Set \texttt{memory.backend = "qmd"} to swap the built-in SQLite indexer for
\href{https://github.com/tobi/qmd}{QMD}: a local-first search sidecar that combines
BM25 + vectors + reranking. Markdown stays the source of truth; OpenClaw shells
out to QMD for retrieval. Key points:

\textbf{Prereqs}

\begin{itemize}
  \item Disabled by default. Opt in per-config (\texttt{memory.backend = "qmd"}).
  \item Install the QMD CLI separately (\texttt{bun install -g https://github.com/tobi/qmd} or grab
\end{itemize}

a release) and make sure the \texttt{qmd} binary is on the gateway’s \texttt{PATH}.
\begin{itemize}
  \item QMD needs an SQLite build that allows extensions (\texttt{brew install sqlite} on
\end{itemize}

macOS).
\begin{itemize}
  \item QMD runs fully locally via Bun + \texttt{node-llama-cpp} and auto-downloads GGUF
\end{itemize}

models from HuggingFace on first use (no separate Ollama daemon required).
\begin{itemize}
  \item The gateway runs QMD in a self-contained XDG home under
\end{itemize}

\texttt{\textasciitilde{}/.openclaw/agents//qmd/} by setting \texttt{XDG\_CONFIG\_HOME} and
\texttt{XDG\_CACHE\_HOME}.
\begin{itemize}
  \item OS support: macOS and Linux work out of the box once Bun + SQLite are
\end{itemize}

installed. Windows is best supported via WSL2.

\textbf{How the sidecar runs}

\begin{itemize}
  \item The gateway writes a self-contained QMD home under
\end{itemize}

\texttt{\textasciitilde{}/.openclaw/agents//qmd/} (config + cache + sqlite DB).
\begin{itemize}
  \item Collections are created via \texttt{qmd collection add} from \texttt{memory.qmd.paths}
\end{itemize}

(plus default workspace memory files), then \texttt{qmd update} + \texttt{qmd embed} run
on boot and on a configurable interval (\texttt{memory.qmd.update.interval},
default 5 m).
\begin{itemize}
  \item The gateway now initializes the QMD manager on startup, so periodic update
\end{itemize}

timers are armed even before the first \texttt{memory\_search} call.
\begin{itemize}
  \item Boot refresh now runs in the background by default so chat startup is not
\end{itemize}

blocked; set \texttt{memory.qmd.update.waitForBootSync = true} to keep the previous
blocking behavior.
\begin{itemize}
  \item Searches run via \texttt{memory.qmd.searchMode} (default \texttt{qmd search --json}; also
\end{itemize}

supports \texttt{vsearch} and \texttt{query}). If the selected mode rejects flags on your
QMD build, OpenClaw retries with \texttt{qmd query}. If QMD fails or the binary is
missing, OpenClaw automatically falls back to the builtin SQLite manager so
memory tools keep working.
\begin{itemize}
  \item OpenClaw does not expose QMD embed batch-size tuning today; batch behavior is
\end{itemize}

controlled by QMD itself.
\begin{itemize}
  \item \textbf{First search may be slow}: QMD may download local GGUF models (reranker/query
\end{itemize}

expansion) on the first \texttt{qmd query} run.
\begin{itemize}
  \item OpenClaw sets \texttt{XDG\_CONFIG\_HOME}/\texttt{XDG\_CACHE\_HOME} automatically when it runs QMD.
  \item If you want to pre-download models manually (and warm the same index OpenClaw
\end{itemize}

uses), run a one-off query with the agent’s XDG dirs.

OpenClaw’s QMD state lives under your \textbf{state dir} (defaults to \texttt{\textasciitilde{}/.openclaw}).
You can point \texttt{qmd} at the exact same index by exporting the same XDG vars
OpenClaw uses:

\begin{lstlisting}[language=bash]
    # Pick the same state dir OpenClaw uses
    STATE_DIR="${OPENCLAW_STATE_DIR:-$HOME/.openclaw}"

    export XDG_CONFIG_HOME="$STATE_DIR/agents/main/qmd/xdg-config"
    export XDG_CACHE_HOME="$STATE_DIR/agents/main/qmd/xdg-cache"

    # (Optional) force an index refresh + embeddings
    qmd update
    qmd embed

    # Warm up / trigger first-time model downloads
    qmd query "test" -c memory-root --json >/dev/null 2>&1
\end{lstlisting}


**Config surface (\texttt{memory.qmd.*})**

\begin{itemize}
  \item \texttt{command} (default \texttt{qmd}): override the executable path.
  \item \texttt{searchMode} (default \texttt{search}): pick which QMD command backs
\end{itemize}

\texttt{memory\_search} (\texttt{search}, \texttt{vsearch}, \texttt{query}).
\begin{itemize}
  \item \texttt{includeDefaultMemory} (default \texttt{true}): auto-index \texttt{MEMORY.md} + \texttt{memory/**/*.md}.
  \item \texttt{paths[]}: add extra directories/files (\texttt{path}, optional \texttt{pattern}, optional
\end{itemize}

stable \texttt{name}).
\begin{itemize}
  \item \texttt{sessions}: opt into session JSONL indexing (\texttt{enabled}, \texttt{retentionDays},
\end{itemize}

\texttt{exportDir}).
\begin{itemize}
  \item \texttt{update}: controls refresh cadence and maintenance execution:
\end{itemize}

(\texttt{interval}, \texttt{debounceMs}, \texttt{onBoot}, \texttt{waitForBootSync}, \texttt{embedInterval},
\texttt{commandTimeoutMs}, \texttt{updateTimeoutMs}, \texttt{embedTimeoutMs}).
\begin{itemize}
  \item \texttt{limits}: clamp recall payload (\texttt{maxResults}, \texttt{maxSnippetChars},
\end{itemize}

\texttt{maxInjectedChars}, \texttt{timeoutMs}).
\begin{itemize}
  \item \texttt{scope}: same schema as \href{/gateway/configuration\#session}{\texttt{session.sendPolicy}}.
\end{itemize}

Default is DM-only (\texttt{deny} all, \texttt{allow} direct chats); loosen it to surface QMD
hits in groups/channels.
\begin{itemize}
  \item \texttt{match.keyPrefix} matches the \textbf{normalized} session key (lowercased, with any
\end{itemize}

leading \texttt{agent::} stripped). Example: \texttt{discord:channel:}.
\begin{itemize}
  \item \texttt{match.rawKeyPrefix} matches the \textbf{raw} session key (lowercased), including
\end{itemize}

\texttt{agent::}. Example: \texttt{agent:main:discord:}.
\begin{itemize}
  \item Legacy: \texttt{match.keyPrefix: "agent:..."} is still treated as a raw-key prefix,
\end{itemize}

but prefer \texttt{rawKeyPrefix} for clarity.
\begin{itemize}
  \item When \texttt{scope} denies a search, OpenClaw logs a warning with the derived
\end{itemize}

\texttt{channel}/\texttt{chatType} so empty results are easier to debug.
\begin{itemize}
  \item Snippets sourced outside the workspace show up as
\end{itemize}

\texttt{qmd//} in \texttt{memory\_search} results; \texttt{memory\_get}
understands that prefix and reads from the configured QMD collection root.
\begin{itemize}
  \item When \texttt{memory.qmd.sessions.enabled = true}, OpenClaw exports sanitized session
\end{itemize}

transcripts (User/Assistant turns) into a dedicated QMD collection under
\texttt{\textasciitilde{}/.openclaw/agents//qmd/sessions/}, so \texttt{memory\_search} can recall recent
conversations without touching the builtin SQLite index.
\begin{itemize}
  \item \texttt{memory\_search} snippets now include a \texttt{Source: } footer when
\end{itemize}

\texttt{memory.citations} is \texttt{auto}/\texttt{on}; set \texttt{memory.citations = "off"} to keep
the path metadata internal (the agent still receives the path for
\texttt{memory\_get}, but the snippet text omits the footer and the system prompt
warns the agent not to cite it).

\textbf{Example}

\begin{lstlisting}[language=json]
memory: {
  backend: "qmd",
  citations: "auto",
  qmd: {
    includeDefaultMemory: true,
    update: { interval: "5m", debounceMs: 15000 },
    limits: { maxResults: 6, timeoutMs: 4000 },
    scope: {
      default: "deny",
      rules: [
        { action: "allow", match: { chatType: "direct" } },
        // Normalized session-key prefix (strips `agent:<id>:`).
        { action: "deny", match: { keyPrefix: "discord:channel:" } },
        // Raw session-key prefix (includes `agent:<id>:`).
        { action: "deny", match: { rawKeyPrefix: "agent:main:discord:" } },
      ]
    },
    paths: [
      { name: "docs", path: "~/notes", pattern: "**/*.md" }
    ]
  }
}
\end{lstlisting}


\textbf{Citations \& fallback}

\begin{itemize}
  \item \texttt{memory.citations} applies regardless of backend (\texttt{auto}/\texttt{on}/\texttt{off}).
  \item When \texttt{qmd} runs, we tag \texttt{status().backend = "qmd"} so diagnostics show which
\end{itemize}

engine served the results. If the QMD subprocess exits or JSON output can’t be
parsed, the search manager logs a warning and returns the builtin provider
(existing Markdown embeddings) until QMD recovers.

\subsubsection{Additional memory paths}


If you want to index Markdown files outside the default workspace layout, add
explicit paths:

\begin{lstlisting}[language=json]
agents: {
  defaults: {
    memorySearch: {
      extraPaths: ["../team-docs", "/srv/shared-notes/overview.md"]
    }
  }
}
\end{lstlisting}


Notes:

\begin{itemize}
  \item Paths can be absolute or workspace-relative.
  \item Directories are scanned recursively for \texttt{.md} files.
  \item Only Markdown files are indexed.
  \item Symlinks are ignored (files or directories).
\end{itemize}


\subsubsection{Gemini embeddings (native)}


Set the provider to \texttt{gemini} to use the Gemini embeddings API directly:

\begin{lstlisting}[language=json]
agents: {
  defaults: {
    memorySearch: {
      provider: "gemini",
      model: "gemini-embedding-001",
      remote: {
        apiKey: "YOUR_GEMINI_API_KEY"
      }
    }
  }
}
\end{lstlisting}


Notes:

\begin{itemize}
  \item \texttt{remote.baseUrl} is optional (defaults to the Gemini API base URL).
  \item \texttt{remote.headers} lets you add extra headers if needed.
  \item Default model: \texttt{gemini-embedding-001}.
\end{itemize}


If you want to use a \textbf{custom OpenAI-compatible endpoint} (OpenRouter, vLLM, or a proxy),
you can use the \texttt{remote} configuration with the OpenAI provider:

\begin{lstlisting}[language=json]
agents: {
  defaults: {
    memorySearch: {
      provider: "openai",
      model: "text-embedding-3-small",
      remote: {
        baseUrl: "https://api.example.com/v1/",
        apiKey: "YOUR_OPENAI_COMPAT_API_KEY",
        headers: { "X-Custom-Header": "value" }
      }
    }
  }
}
\end{lstlisting}


If you don't want to set an API key, use \texttt{memorySearch.provider = "local"} or set
\texttt{memorySearch.fallback = "none"}.

Fallbacks:

\begin{itemize}
  \item \texttt{memorySearch.fallback} can be \texttt{openai}, \texttt{gemini}, \texttt{voyage}, \texttt{mistral}, \texttt{ollama}, \texttt{local}, or \texttt{none}.
  \item The fallback provider is only used when the primary embedding provider fails.
\end{itemize}


Batch indexing (OpenAI + Gemini + Voyage):

\begin{itemize}
  \item Disabled by default. Set \texttt{agents.defaults.memorySearch.remote.batch.enabled = true} to enable for large-corpus indexing (OpenAI, Gemini, and Voyage).
  \item Default behavior waits for batch completion; tune \texttt{remote.batch.wait}, \texttt{remote.batch.pollIntervalMs}, and \texttt{remote.batch.timeoutMinutes} if needed.
  \item Set \texttt{remote.batch.concurrency} to control how many batch jobs we submit in parallel (default: 2).
  \item Batch mode applies when \texttt{memorySearch.provider = "openai"} or \texttt{"gemini"} and uses the corresponding API key.
  \item Gemini batch jobs use the async embeddings batch endpoint and require Gemini Batch API availability.
\end{itemize}


Why OpenAI batch is fast + cheap:

\begin{itemize}
  \item For large backfills, OpenAI is typically the fastest option we support because we can submit many embedding requests in a single batch job and let OpenAI process them asynchronously.
  \item OpenAI offers discounted pricing for Batch API workloads, so large indexing runs are usually cheaper than sending the same requests synchronously.
  \item See the OpenAI Batch API docs and pricing for details:
  \item \href{https://platform.openai.com/docs/api-reference/batch}{https://platform.openai.com/docs/api-reference/batch}
  \item \href{https://platform.openai.com/pricing}{https://platform.openai.com/pricing}
\end{itemize}


Config example:

\begin{lstlisting}[language=json]
agents: {
  defaults: {
    memorySearch: {
      provider: "openai",
      model: "text-embedding-3-small",
      fallback: "openai",
      remote: {
        batch: { enabled: true, concurrency: 2 }
      },
      sync: { watch: true }
    }
  }
}
\end{lstlisting}


Tools:

\begin{itemize}
  \item \texttt{memory\_search} — returns snippets with file + line ranges.
  \item \texttt{memory\_get} — read memory file content by path.
\end{itemize}


Local mode:

\begin{itemize}
  \item Set \texttt{agents.defaults.memorySearch.provider = "local"}.
  \item Provide \texttt{agents.defaults.memorySearch.local.modelPath} (GGUF or \texttt{hf:} URI).
  \item Optional: set \texttt{agents.defaults.memorySearch.fallback = "none"} to avoid remote fallback.
\end{itemize}


\subsubsection{How the memory tools work}


\begin{itemize}
  \item \texttt{memory\_search} semantically searches Markdown chunks (\textasciitilde{}400 token target, 80-token overlap) from \texttt{MEMORY.md} + \texttt{memory/**/*.md}. It returns snippet text (capped \textasciitilde{}700 chars), file path, line range, score, provider/model, and whether we fell back from local → remote embeddings. No full file payload is returned.
  \item \texttt{memory\_get} reads a specific memory Markdown file (workspace-relative), optionally from a starting line and for N lines. Paths outside \texttt{MEMORY.md} / \texttt{memory/} are rejected.
  \item Both tools are enabled only when \texttt{memorySearch.enabled} resolves true for the agent.
\end{itemize}


\subsubsection{What gets indexed (and when)}


\begin{itemize}
  \item File type: Markdown only (\texttt{MEMORY.md}, \texttt{memory/**/*.md}).
  \item Index storage: per-agent SQLite at \texttt{\textasciitilde{}/.openclaw/memory/.sqlite} (configurable via \texttt{agents.defaults.memorySearch.store.path}, supports \texttt{\{agentId\}} token).
  \item Freshness: watcher on \texttt{MEMORY.md} + \texttt{memory/} marks the index dirty (debounce 1.5s). Sync is scheduled on session start, on search, or on an interval and runs asynchronously. Session transcripts use delta thresholds to trigger background sync.
  \item Reindex triggers: the index stores the embedding \textbf{provider/model + endpoint fingerprint + chunking params}. If any of those change, OpenClaw automatically resets and reindexes the entire store.
\end{itemize}


\subsubsection{Hybrid search (BM25 + vector)}


When enabled, OpenClaw combines:

\begin{itemize}
  \item \textbf{Vector similarity} (semantic match, wording can differ)
  \item \textbf{BM25 keyword relevance} (exact tokens like IDs, env vars, code symbols)
\end{itemize}


If full-text search is unavailable on your platform, OpenClaw falls back to vector-only search.

\paragraph{Why hybrid?}


Vector search is great at “this means the same thing”:

\begin{itemize}
  \item “Mac Studio gateway host” vs “the machine running the gateway”
  \item “debounce file updates” vs “avoid indexing on every write”
\end{itemize}


But it can be weak at exact, high-signal tokens:

\begin{itemize}
  \item IDs (\texttt{a828e60}, \texttt{b3b9895a…})
  \item code symbols (\texttt{memorySearch.query.hybrid})
  \item error strings ("sqlite-vec unavailable")
\end{itemize}


BM25 (full-text) is the opposite: strong at exact tokens, weaker at paraphrases.
Hybrid search is the pragmatic middle ground: \textbf{use both retrieval signals} so you get
good results for both "natural language" queries and "needle in a haystack" queries.

\paragraph{How we merge results (the current design)}


Implementation sketch:

\begin{enumerate}
  \item Retrieve a candidate pool from both sides:
\end{enumerate}


\begin{itemize}
  \item \textbf{Vector}: top \texttt{maxResults * candidateMultiplier} by cosine similarity.
  \item \textbf{BM25}: top \texttt{maxResults * candidateMultiplier} by FTS5 BM25 rank (lower is better).
\end{itemize}


\begin{enumerate}
  \item Convert BM25 rank into a 0..1-ish score:
\end{enumerate}


\begin{itemize}
  \item \texttt{textScore = 1 / (1 + max(0, bm25Rank))}
\end{itemize}


\begin{enumerate}
  \item Union candidates by chunk id and compute a weighted score:
\end{enumerate}


\begin{itemize}
  \item \texttt{finalScore = vectorWeight \textit{ vectorScore + textWeight } textScore}
\end{itemize}


Notes:

\begin{itemize}
  \item \texttt{vectorWeight} + \texttt{textWeight} is normalized to 1.0 in config resolution, so weights behave as percentages.
  \item If embeddings are unavailable (or the provider returns a zero-vector), we still run BM25 and return keyword matches.
  \item If FTS5 can't be created, we keep vector-only search (no hard failure).
\end{itemize}


This isn't "IR-theory perfect", but it's simple, fast, and tends to improve recall/precision on real notes.
If we want to get fancier later, common next steps are Reciprocal Rank Fusion (RRF) or score normalization
(min/max or z-score) before mixing.

\paragraph{Post-processing pipeline}


After merging vector and keyword scores, two optional post-processing stages
refine the result list before it reaches the agent:

\begin{lstlisting}
Vector + Keyword → Weighted Merge → Temporal Decay → Sort → MMR → Top-K Results
\end{lstlisting}


Both stages are \textbf{off by default} and can be enabled independently.

\paragraph{MMR re-ranking (diversity)}


When hybrid search returns results, multiple chunks may contain similar or overlapping content.
For example, searching for "home network setup" might return five nearly identical snippets
from different daily notes that all mention the same router configuration.

\textbf{MMR (Maximal Marginal Relevance)} re-ranks the results to balance relevance with diversity,
ensuring the top results cover different aspects of the query instead of repeating the same information.

How it works:

\begin{enumerate}
  \item Results are scored by their original relevance (vector + BM25 weighted score).
  \item MMR iteratively selects results that maximize: \texttt{λ × relevance − (1−λ) × max\_similarity\_to\_selected}.
  \item Similarity between results is measured using Jaccard text similarity on tokenized content.
\end{enumerate}


The \texttt{lambda} parameter controls the trade-off:

\begin{itemize}
  \item \texttt{lambda = 1.0} → pure relevance (no diversity penalty)
  \item \texttt{lambda = 0.0} → maximum diversity (ignores relevance)
  \item Default: \texttt{0.7} (balanced, slight relevance bias)
\end{itemize}


\textbf{Example — query: "home network setup"}

Given these memory files:

\begin{lstlisting}
memory/2026-02-10.md  → "Configured Omada router, set VLAN 10 for IoT devices"
memory/2026-02-08.md  → "Configured Omada router, moved IoT to VLAN 10"
memory/2026-02-05.md  → "Set up AdGuard DNS on 192.168.10.2"
memory/network.md     → "Router: Omada ER605, AdGuard: 192.168.10.2, VLAN 10: IoT"
\end{lstlisting}


Without MMR — top 3 results:

\begin{lstlisting}
1. memory/2026-02-10.md  (score: 0.92)  ← router + VLAN
2. memory/2026-02-08.md  (score: 0.89)  ← router + VLAN (near-duplicate!)
3. memory/network.md     (score: 0.85)  ← reference doc
\end{lstlisting}


With MMR (λ=0.7) — top 3 results:

\begin{lstlisting}
1. memory/2026-02-10.md  (score: 0.92)  ← router + VLAN
2. memory/network.md     (score: 0.85)  ← reference doc (diverse!)
3. memory/2026-02-05.md  (score: 0.78)  ← AdGuard DNS (diverse!)
\end{lstlisting}


The near-duplicate from Feb 8 drops out, and the agent gets three distinct pieces of information.

\textbf{When to enable:} If you notice \texttt{memory\_search} returning redundant or near-duplicate snippets,
especially with daily notes that often repeat similar information across days.

\paragraph{Temporal decay (recency boost)}


Agents with daily notes accumulate hundreds of dated files over time. Without decay,
a well-worded note from six months ago can outrank yesterday's update on the same topic.

\textbf{Temporal decay} applies an exponential multiplier to scores based on the age of each result,
so recent memories naturally rank higher while old ones fade:

\begin{lstlisting}
decayedScore = score × e^(-λ × ageInDays)
\end{lstlisting}


where \texttt{λ = ln(2) / halfLifeDays}.

With the default half-life of 30 days:

\begin{itemize}
  \item Today's notes: \textbf{100\%} of original score
  \item 7 days ago: \textbf{\textasciitilde{}84\%}
  \item 30 days ago: \textbf{50\%}
  \item 90 days ago: \textbf{12.5\%}
  \item 180 days ago: \textbf{\textasciitilde{}1.6\%}
\end{itemize}


\textbf{Evergreen files are never decayed:}

\begin{itemize}
  \item \texttt{MEMORY.md} (root memory file)
  \item Non-dated files in \texttt{memory/} (e.g., \texttt{memory/projects.md}, \texttt{memory/network.md})
  \item These contain durable reference information that should always rank normally.
\end{itemize}


\textbf{Dated daily files} (\texttt{memory/YYYY-MM-DD.md}) use the date extracted from the filename.
Other sources (e.g., session transcripts) fall back to file modification time (\texttt{mtime}).

\textbf{Example — query: "what's Rod's work schedule?"}

Given these memory files (today is Feb 10):

\begin{lstlisting}
memory/2025-09-15.md  → "Rod works Mon-Fri, standup at 10am, pairing at 2pm"  (148 days old)
memory/2026-02-10.md  → "Rod has standup at 14:15, 1:1 with Zeb at 14:45"    (today)
memory/2026-02-03.md  → "Rod started new team, standup moved to 14:15"        (7 days old)
\end{lstlisting}


Without decay:

\begin{lstlisting}
1. memory/2025-09-15.md  (score: 0.91)  ← best semantic match, but stale!
2. memory/2026-02-10.md  (score: 0.82)
3. memory/2026-02-03.md  (score: 0.80)
\end{lstlisting}


With decay (halfLife=30):

\begin{lstlisting}
1. memory/2026-02-10.md  (score: 0.82 × 1.00 = 0.82)  ← today, no decay
2. memory/2026-02-03.md  (score: 0.80 × 0.85 = 0.68)  ← 7 days, mild decay
3. memory/2025-09-15.md  (score: 0.91 × 0.03 = 0.03)  ← 148 days, nearly gone
\end{lstlisting}


The stale September note drops to the bottom despite having the best raw semantic match.

\textbf{When to enable:} If your agent has months of daily notes and you find that old,
stale information outranks recent context. A half-life of 30 days works well for
daily-note-heavy workflows; increase it (e.g., 90 days) if you reference older notes frequently.

\paragraph{Configuration}


Both features are configured under \texttt{memorySearch.query.hybrid}:

\begin{lstlisting}[language=json]
agents: {
  defaults: {
    memorySearch: {
      query: {
        hybrid: {
          enabled: true,
          vectorWeight: 0.7,
          textWeight: 0.3,
          candidateMultiplier: 4,
          // Diversity: reduce redundant results
          mmr: {
            enabled: true,    // default: false
            lambda: 0.7       // 0 = max diversity, 1 = max relevance
          },
          // Recency: boost newer memories
          temporalDecay: {
            enabled: true,    // default: false
            halfLifeDays: 30  // score halves every 30 days
          }
        }
      }
    }
  }
}
\end{lstlisting}


You can enable either feature independently:

\begin{itemize}
  \item \textbf{MMR only} — useful when you have many similar notes but age doesn't matter.
  \item \textbf{Temporal decay only} — useful when recency matters but your results are already diverse.
  \item \textbf{Both} — recommended for agents with large, long-running daily note histories.
\end{itemize}


\subsubsection{Embedding cache}


OpenClaw can cache \textbf{chunk embeddings} in SQLite so reindexing and frequent updates (especially session transcripts) don't re-embed unchanged text.

Config:

\begin{lstlisting}[language=json]
agents: {
  defaults: {
    memorySearch: {
      cache: {
        enabled: true,
        maxEntries: 50000
      }
    }
  }
}
\end{lstlisting}


\subsubsection{Session memory search (experimental)}


You can optionally index \textbf{session transcripts} and surface them via \texttt{memory\_search}.
This is gated behind an experimental flag.

\begin{lstlisting}[language=json]
agents: {
  defaults: {
    memorySearch: {
      experimental: { sessionMemory: true },
      sources: ["memory", "sessions"]
    }
  }
}
\end{lstlisting}


Notes:

\begin{itemize}
  \item Session indexing is \textbf{opt-in} (off by default).
  \item Session updates are debounced and \textbf{indexed asynchronously} once they cross delta thresholds (best-effort).
  \item \texttt{memory\_search} never blocks on indexing; results can be slightly stale until background sync finishes.
  \item Results still include snippets only; \texttt{memory\_get} remains limited to memory files.
  \item Session indexing is isolated per agent (only that agent’s session logs are indexed).
  \item Session logs live on disk (\texttt{\textasciitilde{}/.openclaw/agents//sessions/*.jsonl}). Any process/user with filesystem access can read them, so treat disk access as the trust boundary. For stricter isolation, run agents under separate OS users or hosts.
\end{itemize}


Delta thresholds (defaults shown):

\begin{lstlisting}[language=json]
agents: {
  defaults: {
    memorySearch: {
      sync: {
        sessions: {
          deltaBytes: 100000,   // ~100 KB
          deltaMessages: 50     // JSONL lines
        }
      }
    }
  }
}
\end{lstlisting}


\subsubsection{SQLite vector acceleration (sqlite-vec)}


When the sqlite-vec extension is available, OpenClaw stores embeddings in a
SQLite virtual table (\texttt{vec0}) and performs vector distance queries in the
database. This keeps search fast without loading every embedding into JS.

Configuration (optional):

\begin{lstlisting}[language=json]
agents: {
  defaults: {
    memorySearch: {
      store: {
        vector: {
          enabled: true,
          extensionPath: "/path/to/sqlite-vec"
        }
      }
    }
  }
}
\end{lstlisting}


Notes:

\begin{itemize}
  \item \texttt{enabled} defaults to true; when disabled, search falls back to in-process
\end{itemize}

cosine similarity over stored embeddings.
\begin{itemize}
  \item If the sqlite-vec extension is missing or fails to load, OpenClaw logs the
\end{itemize}

error and continues with the JS fallback (no vector table).
\begin{itemize}
  \item \texttt{extensionPath} overrides the bundled sqlite-vec path (useful for custom builds
\end{itemize}

or non-standard install locations).

\subsubsection{Local embedding auto-download}


\begin{itemize}
  \item Default local embedding model: \texttt{hf:ggml-org/embeddinggemma-300m-qat-q8\_0-GGUF/embeddinggemma-300m-qat-Q8\_0.gguf} (\textasciitilde{}0.6 GB).
  \item When \texttt{memorySearch.provider = "local"}, \texttt{node-llama-cpp} resolves \texttt{modelPath}; if the GGUF is missing it \textbf{auto-downloads} to the cache (or \texttt{local.modelCacheDir} if set), then loads it. Downloads resume on retry.
  \item Native build requirement: run \texttt{pnpm approve-builds}, pick \texttt{node-llama-cpp}, then \texttt{pnpm rebuild node-llama-cpp}.
  \item Fallback: if local setup fails and \texttt{memorySearch.fallback = "openai"}, we automatically switch to remote embeddings (\texttt{openai/text-embedding-3-small} unless overridden) and record the reason.
\end{itemize}


\subsubsection{Custom OpenAI-compatible endpoint example}


\begin{lstlisting}[language=json]
agents: {
  defaults: {
    memorySearch: {
      provider: "openai",
      model: "text-embedding-3-small",
      remote: {
        baseUrl: "https://api.example.com/v1/",
        apiKey: "YOUR_REMOTE_API_KEY",
        headers: {
          "X-Organization": "org-id",
          "X-Project": "project-id"
        }
      }
    }
  }
}
\end{lstlisting}


Notes:

\begin{itemize}
  \item \texttt{remote.\textit{} takes precedence over \texttt{models.providers.openai.}}.
  \item \texttt{remote.headers} merge with OpenAI headers; remote wins on key conflicts. Omit \texttt{remote.headers} to use the OpenAI defaults.
\end{itemize}



\clearpage

% === Source file: concepts/messages.md ===
\section{Messages}


This page ties together how OpenClaw handles inbound messages, sessions, queueing,
streaming, and reasoning visibility.

\subsection{Message flow (high level)}


\begin{lstlisting}
Inbound message
  -> routing/bindings -> session key
  -> queue (if a run is active)
  -> agent run (streaming + tools)
  -> outbound replies (channel limits + chunking)
\end{lstlisting}


Key knobs live in configuration:

\begin{itemize}
  \item \texttt{messages.*} for prefixes, queueing, and group behavior.
  \item \texttt{agents.defaults.*} for block streaming and chunking defaults.
  \item Channel overrides (\texttt{channels.whatsapp.\textit{}, \texttt{channels.telegram.}}, etc.) for caps and streaming toggles.
\end{itemize}


See \href{/gateway/configuration}{Configuration} for full schema.

\subsection{Inbound dedupe}


Channels can redeliver the same message after reconnects. OpenClaw keeps a
short-lived cache keyed by channel/account/peer/session/message id so duplicate
deliveries do not trigger another agent run.

\subsection{Inbound debouncing}


Rapid consecutive messages from the \textbf{same sender} can be batched into a single
agent turn via \texttt{messages.inbound}. Debouncing is scoped per channel + conversation
and uses the most recent message for reply threading/IDs.

Config (global default + per-channel overrides):

\begin{lstlisting}[language=json]
{
  messages: {
    inbound: {
      debounceMs: 2000,
      byChannel: {
        whatsapp: 5000,
        slack: 1500,
        discord: 1500,
      },
    },
  },
}
\end{lstlisting}


Notes:

\begin{itemize}
  \item Debounce applies to \textbf{text-only} messages; media/attachments flush immediately.
  \item Control commands bypass debouncing so they remain standalone.
\end{itemize}


\subsection{Sessions and devices}


Sessions are owned by the gateway, not by clients.

\begin{itemize}
  \item Direct chats collapse into the agent main session key.
  \item Groups/channels get their own session keys.
  \item The session store and transcripts live on the gateway host.
\end{itemize}


Multiple devices/channels can map to the same session, but history is not fully
synced back to every client. Recommendation: use one primary device for long
conversations to avoid divergent context. The Control UI and TUI always show the
gateway-backed session transcript, so they are the source of truth.

Details: \href{/concepts/session}{Session management}.

\subsection{Inbound bodies and history context}


OpenClaw separates the \textbf{prompt body} from the \textbf{command body}:

\begin{itemize}
  \item \texttt{Body}: prompt text sent to the agent. This may include channel envelopes and
\end{itemize}

optional history wrappers.
\begin{itemize}
  \item \texttt{CommandBody}: raw user text for directive/command parsing.
  \item \texttt{RawBody}: legacy alias for \texttt{CommandBody} (kept for compatibility).
\end{itemize}


When a channel supplies history, it uses a shared wrapper:

\begin{itemize}
  \item \texttt{[Chat messages since your last reply - for context]}
  \item \texttt{[Current message - respond to this]}
\end{itemize}


For \textbf{non-direct chats} (groups/channels/rooms), the \textbf{current message body} is prefixed with the
sender label (same style used for history entries). This keeps real-time and queued/history
messages consistent in the agent prompt.

History buffers are \textbf{pending-only}: they include group messages that did \_not\_
trigger a run (for example, mention-gated messages) and \textbf{exclude} messages
already in the session transcript.

Directive stripping only applies to the \textbf{current message} section so history
remains intact. Channels that wrap history should set \texttt{CommandBody} (or
\texttt{RawBody}) to the original message text and keep \texttt{Body} as the combined prompt.
History buffers are configurable via \texttt{messages.groupChat.historyLimit} (global
default) and per-channel overrides like \texttt{channels.slack.historyLimit} or
\texttt{channels.telegram.accounts..historyLimit} (set \texttt{0} to disable).

\subsection{Queueing and followups}


If a run is already active, inbound messages can be queued, steered into the
current run, or collected for a followup turn.

\begin{itemize}
  \item Configure via \texttt{messages.queue} (and \texttt{messages.queue.byChannel}).
  \item Modes: \texttt{interrupt}, \texttt{steer}, \texttt{followup}, \texttt{collect}, plus backlog variants.
\end{itemize}


Details: \href{/concepts/queue}{Queueing}.

\subsection{Streaming, chunking, and batching}


Block streaming sends partial replies as the model produces text blocks.
Chunking respects channel text limits and avoids splitting fenced code.

Key settings:

\begin{itemize}
  \item \texttt{agents.defaults.blockStreamingDefault} (\texttt{on|off}, default off)
  \item \texttt{agents.defaults.blockStreamingBreak} (\texttt{text\_end|message\_end})
  \item \texttt{agents.defaults.blockStreamingChunk} (\texttt{minChars|maxChars|breakPreference})
  \item \texttt{agents.defaults.blockStreamingCoalesce} (idle-based batching)
  \item \texttt{agents.defaults.humanDelay} (human-like pause between block replies)
  \item Channel overrides: \texttt{\textit{.blockStreaming} and \texttt{}.blockStreamingCoalesce} (non-Telegram channels require explicit \texttt{*.blockStreaming: true})
\end{itemize}


Details: \href{/concepts/streaming}{Streaming + chunking}.

\subsection{Reasoning visibility and tokens}


OpenClaw can expose or hide model reasoning:

\begin{itemize}
  \item \texttt{/reasoning on|off|stream} controls visibility.
  \item Reasoning content still counts toward token usage when produced by the model.
  \item Telegram supports reasoning stream into the draft bubble.
\end{itemize}


Details: \href{/tools/thinking}{Thinking + reasoning directives} and \href{/reference/token-use}{Token use}.

\subsection{Prefixes, threading, and replies}


Outbound message formatting is centralized in \texttt{messages}:

\begin{itemize}
  \item \texttt{messages.responsePrefix}, \texttt{channels..responsePrefix}, and \texttt{channels..accounts..responsePrefix} (outbound prefix cascade), plus \texttt{channels.whatsapp.messagePrefix} (WhatsApp inbound prefix)
  \item Reply threading via \texttt{replyToMode} and per-channel defaults
\end{itemize}


Details: \href{/gateway/configuration\#messages}{Configuration} and channel docs.


\clearpage

% === Source file: concepts/model-failover.md ===
\section{Model failover}


OpenClaw handles failures in two stages:

\begin{enumerate}
  \item \textbf{Auth profile rotation} within the current provider.
  \item \textbf{Model fallback} to the next model in \texttt{agents.defaults.model.fallbacks}.
\end{enumerate}


This doc explains the runtime rules and the data that backs them.

\subsection{Auth storage (keys + OAuth)}


OpenClaw uses \textbf{auth profiles} for both API keys and OAuth tokens.

\begin{itemize}
  \item Secrets live in \texttt{\textasciitilde{}/.openclaw/agents//agent/auth-profiles.json} (legacy: \texttt{\textasciitilde{}/.openclaw/agent/auth-profiles.json}).
  \item Config \texttt{auth.profiles} / \texttt{auth.order} are \textbf{metadata + routing only} (no secrets).
  \item Legacy import-only OAuth file: \texttt{\textasciitilde{}/.openclaw/credentials/oauth.json} (imported into \texttt{auth-profiles.json} on first use).
\end{itemize}


More detail: \href{/concepts/oauth}{/concepts/oauth}

Credential types:

\begin{itemize}
  \item \texttt{type: "api\_key"} → \texttt{\{ provider, key \}}
  \item \texttt{type: "oauth"} → \texttt{\{ provider, access, refresh, expires, email? \}} (+ \texttt{projectId}/\texttt{enterpriseUrl} for some providers)
\end{itemize}


\subsection{Profile IDs}


OAuth logins create distinct profiles so multiple accounts can coexist.

\begin{itemize}
  \item Default: \texttt{provider:default} when no email is available.
  \item OAuth with email: \texttt{provider:} (for example \texttt{google-antigravity:user@gmail.com}).
\end{itemize}


Profiles live in \texttt{\textasciitilde{}/.openclaw/agents//agent/auth-profiles.json} under \texttt{profiles}.

\subsection{Rotation order}


When a provider has multiple profiles, OpenClaw chooses an order like this:

\begin{enumerate}
  \item \textbf{Explicit config}: \texttt{auth.order[provider]} (if set).
  \item \textbf{Configured profiles}: \texttt{auth.profiles} filtered by provider.
  \item \textbf{Stored profiles}: entries in \texttt{auth-profiles.json} for the provider.
\end{enumerate}


If no explicit order is configured, OpenClaw uses a round‑robin order:

\begin{itemize}
  \item \textbf{Primary key:} profile type (\textbf{OAuth before API keys}).
  \item \textbf{Secondary key:} \texttt{usageStats.lastUsed} (oldest first, within each type).
  \item \textbf{Cooldown/disabled profiles} are moved to the end, ordered by soonest expiry.
\end{itemize}


\subsubsection{Session stickiness (cache-friendly)}


OpenClaw \textbf{pins the chosen auth profile per session} to keep provider caches warm.
It does \textbf{not} rotate on every request. The pinned profile is reused until:

\begin{itemize}
  \item the session is reset (\texttt{/new} / \texttt{/reset})
  \item a compaction completes (compaction count increments)
  \item the profile is in cooldown/disabled
\end{itemize}


Manual selection via \texttt{/model …@} sets a \textbf{user override} for that session
and is not auto‑rotated until a new session starts.

Auto‑pinned profiles (selected by the session router) are treated as a \textbf{preference}:
they are tried first, but OpenClaw may rotate to another profile on rate limits/timeouts.
User‑pinned profiles stay locked to that profile; if it fails and model fallbacks
are configured, OpenClaw moves to the next model instead of switching profiles.

\subsubsection{Why OAuth can “look lost”}


If you have both an OAuth profile and an API key profile for the same provider, round‑robin can switch between them across messages unless pinned. To force a single profile:

\begin{itemize}
  \item Pin with \texttt{auth.order[provider] = ["provider:profileId"]}, or
  \item Use a per-session override via \texttt{/model …} with a profile override (when supported by your UI/chat surface).
\end{itemize}


\subsection{Cooldowns}


When a profile fails due to auth/rate‑limit errors (or a timeout that looks
like rate limiting), OpenClaw marks it in cooldown and moves to the next profile.
Format/invalid‑request errors (for example Cloud Code Assist tool call ID
validation failures) are treated as failover‑worthy and use the same cooldowns.
OpenAI-compatible stop-reason errors such as \texttt{Unhandled stop reason: error},
\texttt{stop reason: error}, and \texttt{reason: error} are classified as timeout/failover
signals.

Cooldowns use exponential backoff:

\begin{itemize}
  \item 1 minute
  \item 5 minutes
  \item 25 minutes
  \item 1 hour (cap)
\end{itemize}


State is stored in \texttt{auth-profiles.json} under \texttt{usageStats}:

\begin{lstlisting}[language=json]
{
  "usageStats": {
    "provider:profile": {
      "lastUsed": 1736160000000,
      "cooldownUntil": 1736160600000,
      "errorCount": 2
    }
  }
}
\end{lstlisting}


\subsection{Billing disables}


Billing/credit failures (for example “insufficient credits” / “credit balance too low”) are treated as failover‑worthy, but they’re usually not transient. Instead of a short cooldown, OpenClaw marks the profile as \textbf{disabled} (with a longer backoff) and rotates to the next profile/provider.

State is stored in \texttt{auth-profiles.json}:

\begin{lstlisting}[language=json]
{
  "usageStats": {
    "provider:profile": {
      "disabledUntil": 1736178000000,
      "disabledReason": "billing"
    }
  }
}
\end{lstlisting}


Defaults:

\begin{itemize}
  \item Billing backoff starts at \textbf{5 hours}, doubles per billing failure, and caps at \textbf{24 hours}.
  \item Backoff counters reset if the profile hasn’t failed for \textbf{24 hours} (configurable).
\end{itemize}


\subsection{Model fallback}


If all profiles for a provider fail, OpenClaw moves to the next model in
\texttt{agents.defaults.model.fallbacks}. This applies to auth failures, rate limits, and
timeouts that exhausted profile rotation (other errors do not advance fallback).

When a run starts with a model override (hooks or CLI), fallbacks still end at
\texttt{agents.defaults.model.primary} after trying any configured fallbacks.

\subsection{Related config}


See \href{/gateway/configuration}{Gateway configuration} for:

\begin{itemize}
  \item \texttt{auth.profiles} / \texttt{auth.order}
  \item \texttt{auth.cooldowns.billingBackoffHours} / \texttt{auth.cooldowns.billingBackoffHoursByProvider}
  \item \texttt{auth.cooldowns.billingMaxHours} / \texttt{auth.cooldowns.failureWindowHours}
  \item \texttt{agents.defaults.model.primary} / \texttt{agents.defaults.model.fallbacks}
  \item \texttt{agents.defaults.imageModel} routing
\end{itemize}


See \href{/concepts/models}{Models} for the broader model selection and fallback overview.


\clearpage

% === Source file: concepts/model-providers.md ===
\section{Model providers}


This page covers \textbf{LLM/model providers} (not chat channels like WhatsApp/Telegram).
For model selection rules, see \href{/concepts/models}{/concepts/models}.

\subsection{Quick rules}


\begin{itemize}
  \item Model refs use \texttt{provider/model} (example: \texttt{opencode/claude-opus-4-6}).
  \item If you set \texttt{agents.defaults.models}, it becomes the allowlist.
  \item CLI helpers: \texttt{openclaw onboard}, \texttt{openclaw models list}, \texttt{openclaw models set }.
\end{itemize}


\subsection{API key rotation}


\begin{itemize}
  \item Supports generic provider rotation for selected providers.
  \item Configure multiple keys via:
  \item \texttt{OPENCLAW\_LIVE\_\_KEY} (single live override, highest priority)
  \item \texttt{\_API\_KEYS} (comma or semicolon list)
  \item \texttt{\_API\_KEY} (primary key)
  \item \texttt{\_API\_KEY\_*} (numbered list, e.g. \texttt{\_API\_KEY\_1})
  \item For Google providers, \texttt{GOOGLE\_API\_KEY} is also included as fallback.
  \item Key selection order preserves priority and deduplicates values.
  \item Requests are retried with the next key only on rate-limit responses (for example \texttt{429}, \texttt{rate\_limit}, \texttt{quota}, \texttt{resource exhausted}).
  \item Non-rate-limit failures fail immediately; no key rotation is attempted.
  \item When all candidate keys fail, the final error is returned from the last attempt.
\end{itemize}


\subsection{Built-in providers (pi-ai catalog)}


OpenClaw ships with the pi‑ai catalog. These providers require \textbf{no}
\texttt{models.providers} config; just set auth + pick a model.

\subsubsection{OpenAI}


\begin{itemize}
  \item Provider: \texttt{openai}
  \item Auth: \texttt{OPENAI\_API\_KEY}
  \item Optional rotation: \texttt{OPENAI\_API\_KEYS}, \texttt{OPENAI\_API\_KEY\_1}, \texttt{OPENAI\_API\_KEY\_2}, plus \texttt{OPENCLAW\_LIVE\_OPENAI\_KEY} (single override)
  \item Example model: \texttt{openai/gpt-5.1-codex}
  \item CLI: \texttt{openclaw onboard --auth-choice openai-api-key}
  \item Default transport is \texttt{auto} (WebSocket-first, SSE fallback)
  \item Override per model via \texttt{agents.defaults.models["openai/"].params.transport} (\texttt{"sse"}, \texttt{"websocket"}, or \texttt{"auto"})
  \item OpenAI Responses WebSocket warm-up defaults to enabled via \texttt{params.openaiWsWarmup} (\texttt{true}/\texttt{false})
\end{itemize}


\begin{lstlisting}[language=json]
{
  agents: { defaults: { model: { primary: "openai/gpt-5.1-codex" } } },
}
\end{lstlisting}


\subsubsection{Anthropic}


\begin{itemize}
  \item Provider: \texttt{anthropic}
  \item Auth: \texttt{ANTHROPIC\_API\_KEY} or \texttt{claude setup-token}
  \item Optional rotation: \texttt{ANTHROPIC\_API\_KEYS}, \texttt{ANTHROPIC\_API\_KEY\_1}, \texttt{ANTHROPIC\_API\_KEY\_2}, plus \texttt{OPENCLAW\_LIVE\_ANTHROPIC\_KEY} (single override)
  \item Example model: \texttt{anthropic/claude-opus-4-6}
  \item CLI: \texttt{openclaw onboard --auth-choice token} (paste setup-token) or \texttt{openclaw models auth paste-token --provider anthropic}
  \item Policy note: setup-token support is technical compatibility; Anthropic has blocked some subscription usage outside Claude Code in the past. Verify current Anthropic terms and decide based on your risk tolerance.
  \item Recommendation: Anthropic API key auth is the safer, recommended path over subscription setup-token auth.
\end{itemize}


\begin{lstlisting}[language=json]
{
  agents: { defaults: { model: { primary: "anthropic/claude-opus-4-6" } } },
}
\end{lstlisting}


\subsubsection{OpenAI Code (Codex)}


\begin{itemize}
  \item Provider: \texttt{openai-codex}
  \item Auth: OAuth (ChatGPT)
  \item Example model: \texttt{openai-codex/gpt-5.3-codex}
  \item CLI: \texttt{openclaw onboard --auth-choice openai-codex} or \texttt{openclaw models auth login --provider openai-codex}
  \item Default transport is \texttt{auto} (WebSocket-first, SSE fallback)
  \item Override per model via \texttt{agents.defaults.models["openai-codex/"].params.transport} (\texttt{"sse"}, \texttt{"websocket"}, or \texttt{"auto"})
  \item Policy note: OpenAI Codex OAuth is explicitly supported for external tools/workflows like OpenClaw.
\end{itemize}


\begin{lstlisting}[language=json]
{
  agents: { defaults: { model: { primary: "openai-codex/gpt-5.3-codex" } } },
}
\end{lstlisting}


\subsubsection{OpenCode Zen}


\begin{itemize}
  \item Provider: \texttt{opencode}
  \item Auth: \texttt{OPENCODE\_API\_KEY} (or \texttt{OPENCODE\_ZEN\_API\_KEY})
  \item Example model: \texttt{opencode/claude-opus-4-6}
  \item CLI: \texttt{openclaw onboard --auth-choice opencode-zen}
\end{itemize}


\begin{lstlisting}[language=json]
{
  agents: { defaults: { model: { primary: "opencode/claude-opus-4-6" } } },
}
\end{lstlisting}


\subsubsection{Google Gemini (API key)}


\begin{itemize}
  \item Provider: \texttt{google}
  \item Auth: \texttt{GEMINI\_API\_KEY}
  \item Optional rotation: \texttt{GEMINI\_API\_KEYS}, \texttt{GEMINI\_API\_KEY\_1}, \texttt{GEMINI\_API\_KEY\_2}, \texttt{GOOGLE\_API\_KEY} fallback, and \texttt{OPENCLAW\_LIVE\_GEMINI\_KEY} (single override)
  \item Example model: \texttt{google/gemini-3-pro-preview}
  \item CLI: \texttt{openclaw onboard --auth-choice gemini-api-key}
\end{itemize}


\subsubsection{Google Vertex, Antigravity, and Gemini CLI}


\begin{itemize}
  \item Providers: \texttt{google-vertex}, \texttt{google-antigravity}, \texttt{google-gemini-cli}
  \item Auth: Vertex uses gcloud ADC; Antigravity/Gemini CLI use their respective auth flows
  \item Caution: Antigravity and Gemini CLI OAuth in OpenClaw are unofficial integrations. Some users have reported Google account restrictions after using third-party clients. Review Google terms and use a non-critical account if you choose to proceed.
  \item Antigravity OAuth is shipped as a bundled plugin (\texttt{google-antigravity-auth}, disabled by default).
  \item Enable: \texttt{openclaw plugins enable google-antigravity-auth}
  \item Login: \texttt{openclaw models auth login --provider google-antigravity --set-default}
  \item Gemini CLI OAuth is shipped as a bundled plugin (\texttt{google-gemini-cli-auth}, disabled by default).
  \item Enable: \texttt{openclaw plugins enable google-gemini-cli-auth}
  \item Login: \texttt{openclaw models auth login --provider google-gemini-cli --set-default}
  \item Note: you do \textbf{not} paste a client id or secret into \texttt{openclaw.json}. The CLI login flow stores
\end{itemize}

tokens in auth profiles on the gateway host.

\subsubsection{Z.AI (GLM)}


\begin{itemize}
  \item Provider: \texttt{zai}
  \item Auth: \texttt{ZAI\_API\_KEY}
  \item Example model: \texttt{zai/glm-5}
  \item CLI: \texttt{openclaw onboard --auth-choice zai-api-key}
  \item Aliases: \texttt{z.ai/\textit{} and \texttt{z-ai/}} normalize to \texttt{zai/*}
\end{itemize}


\subsubsection{Vercel AI Gateway}


\begin{itemize}
  \item Provider: \texttt{vercel-ai-gateway}
  \item Auth: \texttt{AI\_GATEWAY\_API\_KEY}
  \item Example model: \texttt{vercel-ai-gateway/anthropic/claude-opus-4.6}
  \item CLI: \texttt{openclaw onboard --auth-choice ai-gateway-api-key}
\end{itemize}


\subsubsection{Kilo Gateway}


\begin{itemize}
  \item Provider: \texttt{kilocode}
  \item Auth: \texttt{KILOCODE\_API\_KEY}
  \item Example model: \texttt{kilocode/anthropic/claude-opus-4.6}
  \item CLI: \texttt{openclaw onboard --kilocode-api-key }
  \item Base URL: \texttt{https://api.kilo.ai/api/gateway/}
  \item Expanded built-in catalog includes GLM-5 Free, MiniMax M2.5 Free, GPT-5.2, Gemini 3 Pro Preview, Gemini 3 Flash Preview, Grok Code Fast 1, and Kimi K2.5.
\end{itemize}


See \href{/providers/kilocode}{/providers/kilocode} for setup details.

\subsubsection{Other built-in providers}


\begin{itemize}
  \item OpenRouter: \texttt{openrouter} (\texttt{OPENROUTER\_API\_KEY})
  \item Example model: \texttt{openrouter/anthropic/claude-sonnet-4-5}
  \item Kilo Gateway: \texttt{kilocode} (\texttt{KILOCODE\_API\_KEY})
  \item Example model: \texttt{kilocode/anthropic/claude-opus-4.6}
  \item xAI: \texttt{xai} (\texttt{XAI\_API\_KEY})
  \item Mistral: \texttt{mistral} (\texttt{MISTRAL\_API\_KEY})
  \item Example model: \texttt{mistral/mistral-large-latest}
  \item CLI: \texttt{openclaw onboard --auth-choice mistral-api-key}
  \item Groq: \texttt{groq} (\texttt{GROQ\_API\_KEY})
  \item Cerebras: \texttt{cerebras} (\texttt{CEREBRAS\_API\_KEY})
  \item GLM models on Cerebras use ids \texttt{zai-glm-4.7} and \texttt{zai-glm-4.6}.
  \item OpenAI-compatible base URL: \texttt{https://api.cerebras.ai/v1}.
  \item GitHub Copilot: \texttt{github-copilot} (\texttt{COPILOT\_GITHUB\_TOKEN} / \texttt{GH\_TOKEN} / \texttt{GITHUB\_TOKEN})
  \item Hugging Face Inference: \texttt{huggingface} (\texttt{HUGGINGFACE\_HUB\_TOKEN} or \texttt{HF\_TOKEN}) — OpenAI-compatible router; example model: \texttt{huggingface/deepseek-ai/DeepSeek-R1}; CLI: \texttt{openclaw onboard --auth-choice huggingface-api-key}. See \href{/providers/huggingface}{Hugging Face (Inference)}.
\end{itemize}


\subsection{Providers via models.providers (custom/base URL)}


Use \texttt{models.providers} (or \texttt{models.json}) to add \textbf{custom} providers or
OpenAI/Anthropic‑compatible proxies.

\subsubsection{Moonshot AI (Kimi)}


Moonshot uses OpenAI-compatible endpoints, so configure it as a custom provider:

\begin{itemize}
  \item Provider: \texttt{moonshot}
  \item Auth: \texttt{MOONSHOT\_API\_KEY}
  \item Example model: \texttt{moonshot/kimi-k2.5}
\end{itemize}


Kimi K2 model IDs:


\{/\_ moonshot-kimi-k2-model-refs:start \_/ \&\& null\}


\begin{itemize}
  \item \texttt{moonshot/kimi-k2.5}
  \item \texttt{moonshot/kimi-k2-0905-preview}
  \item \texttt{moonshot/kimi-k2-turbo-preview}
  \item \texttt{moonshot/kimi-k2-thinking}
  \item \texttt{moonshot/kimi-k2-thinking-turbo}
\end{itemize}

\{/\_ moonshot-kimi-k2-model-refs:end \_/ \&\& null\}

\begin{lstlisting}[language=json]
{
  agents: {
    defaults: { model: { primary: "moonshot/kimi-k2.5" } },
  },
  models: {
    mode: "merge",
    providers: {
      moonshot: {
        baseUrl: "https://api.moonshot.ai/v1",
        apiKey: "${MOONSHOT_API_KEY}",
        api: "openai-completions",
        models: [{ id: "kimi-k2.5", name: "Kimi K2.5" }],
      },
    },
  },
}
\end{lstlisting}


\subsubsection{Kimi Coding}


Kimi Coding uses Moonshot AI's Anthropic-compatible endpoint:

\begin{itemize}
  \item Provider: \texttt{kimi-coding}
  \item Auth: \texttt{KIMI\_API\_KEY}
  \item Example model: \texttt{kimi-coding/k2p5}
\end{itemize}


\begin{lstlisting}[language=json]
{
  env: { KIMI_API_KEY: "sk-..." },
  agents: {
    defaults: { model: { primary: "kimi-coding/k2p5" } },
  },
}
\end{lstlisting}


\subsubsection{Qwen OAuth (free tier)}


Qwen provides OAuth access to Qwen Coder + Vision via a device-code flow.
Enable the bundled plugin, then log in:

\begin{lstlisting}[language=bash]
openclaw plugins enable qwen-portal-auth
openclaw models auth login --provider qwen-portal --set-default
\end{lstlisting}


Model refs:

\begin{itemize}
  \item \texttt{qwen-portal/coder-model}
  \item \texttt{qwen-portal/vision-model}
\end{itemize}


See \href{/providers/qwen}{/providers/qwen} for setup details and notes.

\subsubsection{Volcano Engine (Doubao)}


Volcano Engine (火山引擎) provides access to Doubao and other models in China.

\begin{itemize}
  \item Provider: \texttt{volcengine} (coding: \texttt{volcengine-plan})
  \item Auth: \texttt{VOLCANO\_ENGINE\_API\_KEY}
  \item Example model: \texttt{volcengine/doubao-seed-1-8-251228}
  \item CLI: \texttt{openclaw onboard --auth-choice volcengine-api-key}
\end{itemize}


\begin{lstlisting}[language=json]
{
  agents: {
    defaults: { model: { primary: "volcengine/doubao-seed-1-8-251228" } },
  },
}
\end{lstlisting}


Available models:

\begin{itemize}
  \item \texttt{volcengine/doubao-seed-1-8-251228} (Doubao Seed 1.8)
  \item \texttt{volcengine/doubao-seed-code-preview-251028}
  \item \texttt{volcengine/kimi-k2-5-260127} (Kimi K2.5)
  \item \texttt{volcengine/glm-4-7-251222} (GLM 4.7)
  \item \texttt{volcengine/deepseek-v3-2-251201} (DeepSeek V3.2 128K)
\end{itemize}


Coding models (\texttt{volcengine-plan}):

\begin{itemize}
  \item \texttt{volcengine-plan/ark-code-latest}
  \item \texttt{volcengine-plan/doubao-seed-code}
  \item \texttt{volcengine-plan/kimi-k2.5}
  \item \texttt{volcengine-plan/kimi-k2-thinking}
  \item \texttt{volcengine-plan/glm-4.7}
\end{itemize}


\subsubsection{BytePlus (International)}


BytePlus ARK provides access to the same models as Volcano Engine for international users.

\begin{itemize}
  \item Provider: \texttt{byteplus} (coding: \texttt{byteplus-plan})
  \item Auth: \texttt{BYTEPLUS\_API\_KEY}
  \item Example model: \texttt{byteplus/seed-1-8-251228}
  \item CLI: \texttt{openclaw onboard --auth-choice byteplus-api-key}
\end{itemize}


\begin{lstlisting}[language=json]
{
  agents: {
    defaults: { model: { primary: "byteplus/seed-1-8-251228" } },
  },
}
\end{lstlisting}


Available models:

\begin{itemize}
  \item \texttt{byteplus/seed-1-8-251228} (Seed 1.8)
  \item \texttt{byteplus/kimi-k2-5-260127} (Kimi K2.5)
  \item \texttt{byteplus/glm-4-7-251222} (GLM 4.7)
\end{itemize}


Coding models (\texttt{byteplus-plan}):

\begin{itemize}
  \item \texttt{byteplus-plan/ark-code-latest}
  \item \texttt{byteplus-plan/doubao-seed-code}
  \item \texttt{byteplus-plan/kimi-k2.5}
  \item \texttt{byteplus-plan/kimi-k2-thinking}
  \item \texttt{byteplus-plan/glm-4.7}
\end{itemize}


\subsubsection{Synthetic}


Synthetic provides Anthropic-compatible models behind the \texttt{synthetic} provider:

\begin{itemize}
  \item Provider: \texttt{synthetic}
  \item Auth: \texttt{SYNTHETIC\_API\_KEY}
  \item Example model: \texttt{synthetic/hf:MiniMaxAI/MiniMax-M2.5}
  \item CLI: \texttt{openclaw onboard --auth-choice synthetic-api-key}
\end{itemize}


\begin{lstlisting}[language=json]
{
  agents: {
    defaults: { model: { primary: "synthetic/hf:MiniMaxAI/MiniMax-M2.5" } },
  },
  models: {
    mode: "merge",
    providers: {
      synthetic: {
        baseUrl: "https://api.synthetic.new/anthropic",
        apiKey: "${SYNTHETIC_API_KEY}",
        api: "anthropic-messages",
        models: [{ id: "hf:MiniMaxAI/MiniMax-M2.5", name: "MiniMax M2.5" }],
      },
    },
  },
}
\end{lstlisting}


\subsubsection{MiniMax}


MiniMax is configured via \texttt{models.providers} because it uses custom endpoints:

\begin{itemize}
  \item MiniMax (Anthropic‑compatible): \texttt{--auth-choice minimax-api}
  \item Auth: \texttt{MINIMAX\_API\_KEY}
\end{itemize}


See \href{/providers/minimax}{/providers/minimax} for setup details, model options, and config snippets.

\subsubsection{Ollama}


Ollama is a local LLM runtime that provides an OpenAI-compatible API:

\begin{itemize}
  \item Provider: \texttt{ollama}
  \item Auth: None required (local server)
  \item Example model: \texttt{ollama/llama3.3}
  \item Installation: \href{https://ollama.ai}{https://ollama.ai}
\end{itemize}


\begin{lstlisting}[language=bash]
# Install Ollama, then pull a model:
ollama pull llama3.3
\end{lstlisting}


\begin{lstlisting}[language=json]
{
  agents: {
    defaults: { model: { primary: "ollama/llama3.3" } },
  },
}
\end{lstlisting}


Ollama is automatically detected when running locally at \texttt{http://127.0.0.1:11434/v1}. See \href{/providers/ollama}{/providers/ollama} for model recommendations and custom configuration.

\subsubsection{vLLM}


vLLM is a local (or self-hosted) OpenAI-compatible server:

\begin{itemize}
  \item Provider: \texttt{vllm}
  \item Auth: Optional (depends on your server)
  \item Default base URL: \texttt{http://127.0.0.1:8000/v1}
\end{itemize}


To opt in to auto-discovery locally (any value works if your server doesn’t enforce auth):

\begin{lstlisting}[language=bash]
export VLLM_API_KEY="vllm-local"
\end{lstlisting}


Then set a model (replace with one of the IDs returned by \texttt{/v1/models}):

\begin{lstlisting}[language=json]
{
  agents: {
    defaults: { model: { primary: "vllm/your-model-id" } },
  },
}
\end{lstlisting}


See \href{/providers/vllm}{/providers/vllm} for details.

\subsubsection{Local proxies (LM Studio, vLLM, LiteLLM, etc.)}


Example (OpenAI‑compatible):

\begin{lstlisting}[language=json]
{
  agents: {
    defaults: {
      model: { primary: "lmstudio/minimax-m2.5-gs32" },
      models: { "lmstudio/minimax-m2.5-gs32": { alias: "Minimax" } },
    },
  },
  models: {
    providers: {
      lmstudio: {
        baseUrl: "http://localhost:1234/v1",
        apiKey: "LMSTUDIO_KEY",
        api: "openai-completions",
        models: [
          {
            id: "minimax-m2.5-gs32",
            name: "MiniMax M2.5",
            reasoning: false,
            input: ["text"],
            cost: { input: 0, output: 0, cacheRead: 0, cacheWrite: 0 },
            contextWindow: 200000,
            maxTokens: 8192,
          },
        ],
      },
    },
  },
}
\end{lstlisting}


Notes:

\begin{itemize}
  \item For custom providers, \texttt{reasoning}, \texttt{input}, \texttt{cost}, \texttt{contextWindow}, and \texttt{maxTokens} are optional.
\end{itemize}

When omitted, OpenClaw defaults to:
\begin{itemize}
  \item \texttt{reasoning: false}
  \item \texttt{input: ["text"]}
  \item \texttt{cost: \{ input: 0, output: 0, cacheRead: 0, cacheWrite: 0 \}}
  \item \texttt{contextWindow: 200000}
  \item \texttt{maxTokens: 8192}
  \item Recommended: set explicit values that match your proxy/model limits.
  \item For \texttt{api: "openai-completions"} on non-native endpoints (any non-empty \texttt{baseUrl} whose host is not \texttt{api.openai.com}), OpenClaw forces \texttt{compat.supportsDeveloperRole: false} to avoid provider 400 errors for unsupported \texttt{developer} roles.
  \item If \texttt{baseUrl} is empty/omitted, OpenClaw keeps the default OpenAI behavior (which resolves to \texttt{api.openai.com}).
  \item For safety, an explicit \texttt{compat.supportsDeveloperRole: true} is still overridden on non-native \texttt{openai-completions} endpoints.
\end{itemize}


\subsection{CLI examples}


\begin{lstlisting}[language=bash]
openclaw onboard --auth-choice opencode-zen
openclaw models set opencode/claude-opus-4-6
openclaw models list
\end{lstlisting}


See also: \href{/gateway/configuration}{/gateway/configuration} for full configuration examples.


\clearpage

% === Source file: concepts/models.md ===
\section{Models CLI}


See \href{/concepts/model-failover}{/concepts/model-failover} for auth profile
rotation, cooldowns, and how that interacts with fallbacks.
Quick provider overview + examples: \href{/concepts/model-providers}{/concepts/model-providers}.

\subsection{How model selection works}


OpenClaw selects models in this order:

\begin{enumerate}
  \item \textbf{Primary} model (\texttt{agents.defaults.model.primary} or \texttt{agents.defaults.model}).
  \item \textbf{Fallbacks} in \texttt{agents.defaults.model.fallbacks} (in order).
  \item \textbf{Provider auth failover} happens inside a provider before moving to the
\end{enumerate}

next model.

Related:

\begin{itemize}
  \item \texttt{agents.defaults.models} is the allowlist/catalog of models OpenClaw can use (plus aliases).
  \item \texttt{agents.defaults.imageModel} is used \textbf{only when} the primary model can’t accept images.
  \item Per-agent defaults can override \texttt{agents.defaults.model} via \texttt{agents.list[].model} plus bindings (see \href{/concepts/multi-agent}{/concepts/multi-agent}).
\end{itemize}


\subsection{Quick model policy}


\begin{itemize}
  \item Set your primary to the strongest latest-generation model available to you.
  \item Use fallbacks for cost/latency-sensitive tasks and lower-stakes chat.
  \item For tool-enabled agents or untrusted inputs, avoid older/weaker model tiers.
\end{itemize}


\subsection{Setup wizard (recommended)}


If you don’t want to hand-edit config, run the onboarding wizard:

\begin{lstlisting}[language=bash]
openclaw onboard
\end{lstlisting}


It can set up model + auth for common providers, including **OpenAI Code (Codex)
subscription\textbf{ (OAuth) and }Anthropic** (API key or \texttt{claude setup-token}).

\subsection{Config keys (overview)}


\begin{itemize}
  \item \texttt{agents.defaults.model.primary} and \texttt{agents.defaults.model.fallbacks}
  \item \texttt{agents.defaults.imageModel.primary} and \texttt{agents.defaults.imageModel.fallbacks}
  \item \texttt{agents.defaults.models} (allowlist + aliases + provider params)
  \item \texttt{models.providers} (custom providers written into \texttt{models.json})
\end{itemize}


Model refs are normalized to lowercase. Provider aliases like \texttt{z.ai/*} normalize
to \texttt{zai/*}.

Provider configuration examples (including OpenCode Zen) live in
\href{/gateway/configuration\#opencode-zen-multi-model-proxy}{/gateway/configuration}.

\subsection{“Model is not allowed” (and why replies stop)}


If \texttt{agents.defaults.models} is set, it becomes the \textbf{allowlist} for \texttt{/model} and for
session overrides. When a user selects a model that isn’t in that allowlist,
OpenClaw returns:

\begin{lstlisting}
Model "provider/model" is not allowed. Use /model to list available models.
\end{lstlisting}


This happens \textbf{before} a normal reply is generated, so the message can feel
like it “didn’t respond.” The fix is to either:

\begin{itemize}
  \item Add the model to \texttt{agents.defaults.models}, or
  \item Clear the allowlist (remove \texttt{agents.defaults.models}), or
  \item Pick a model from \texttt{/model list}.
\end{itemize}


Example allowlist config:

\begin{lstlisting}[language=json]
{
  agent: {
    model: { primary: "anthropic/claude-sonnet-4-5" },
    models: {
      "anthropic/claude-sonnet-4-5": { alias: "Sonnet" },
      "anthropic/claude-opus-4-6": { alias: "Opus" },
    },
  },
}
\end{lstlisting}


\subsection{Switching models in chat (/model)}


You can switch models for the current session without restarting:

\begin{lstlisting}
/model
/model list
/model 3
/model openai/gpt-5.2
/model status
\end{lstlisting}


Notes:

\begin{itemize}
  \item \texttt{/model} (and \texttt{/model list}) is a compact, numbered picker (model family + available providers).
  \item On Discord, \texttt{/model} and \texttt{/models} open an interactive picker with provider and model dropdowns plus a Submit step.
  \item \texttt{/model <\#>} selects from that picker.
  \item \texttt{/model status} is the detailed view (auth candidates and, when configured, provider endpoint \texttt{baseUrl} + \texttt{api} mode).
  \item Model refs are parsed by splitting on the \textbf{first} \texttt{/}. Use \texttt{provider/model} when typing \texttt{/model }.
  \item If the model ID itself contains \texttt{/} (OpenRouter-style), you must include the provider prefix (example: \texttt{/model openrouter/moonshotai/kimi-k2}).
  \item If you omit the provider, OpenClaw treats the input as an alias or a model for the \textbf{default provider} (only works when there is no \texttt{/} in the model ID).
\end{itemize}


Full command behavior/config: \href{/tools/slash-commands}{Slash commands}.

\subsection{CLI commands}


\begin{lstlisting}[language=bash]
openclaw models list
openclaw models status
openclaw models set <provider/model>
openclaw models set-image <provider/model>

openclaw models aliases list
openclaw models aliases add <alias> <provider/model>
openclaw models aliases remove <alias>

openclaw models fallbacks list
openclaw models fallbacks add <provider/model>
openclaw models fallbacks remove <provider/model>
openclaw models fallbacks clear

openclaw models image-fallbacks list
openclaw models image-fallbacks add <provider/model>
openclaw models image-fallbacks remove <provider/model>
openclaw models image-fallbacks clear
\end{lstlisting}


\texttt{openclaw models} (no subcommand) is a shortcut for \texttt{models status}.

\subsubsection{models list}


Shows configured models by default. Useful flags:

\begin{itemize}
  \item \texttt{--all}: full catalog
  \item \texttt{--local}: local providers only
  \item \texttt{--provider }: filter by provider
  \item \texttt{--plain}: one model per line
  \item \texttt{--json}: machine‑readable output
\end{itemize}


\subsubsection{models status}


Shows the resolved primary model, fallbacks, image model, and an auth overview
of configured providers. It also surfaces OAuth expiry status for profiles found
in the auth store (warns within 24h by default). \texttt{--plain} prints only the
resolved primary model.
OAuth status is always shown (and included in \texttt{--json} output). If a configured
provider has no credentials, \texttt{models status} prints a \textbf{Missing auth} section.
JSON includes \texttt{auth.oauth} (warn window + profiles) and \texttt{auth.providers}
(effective auth per provider).
Use \texttt{--check} for automation (exit \texttt{1} when missing/expired, \texttt{2} when expiring).

Auth choice is provider/account dependent. For always-on gateway hosts, API keys are usually the most predictable; subscription token flows are also supported.

Example (Anthropic setup-token):

\begin{lstlisting}[language=bash]
claude setup-token
openclaw models status
\end{lstlisting}


\subsection{Scanning (OpenRouter free models)}


\texttt{openclaw models scan} inspects OpenRouter’s \textbf{free model catalog} and can
optionally probe models for tool and image support.

Key flags:

\begin{itemize}
  \item \texttt{--no-probe}: skip live probes (metadata only)
  \item \texttt{--min-params }: minimum parameter size (billions)
  \item \texttt{--max-age-days }: skip older models
  \item \texttt{--provider }: provider prefix filter
  \item \texttt{--max-candidates }: fallback list size
  \item \texttt{--set-default}: set \texttt{agents.defaults.model.primary} to the first selection
  \item \texttt{--set-image}: set \texttt{agents.defaults.imageModel.primary} to the first image selection
\end{itemize}


Probing requires an OpenRouter API key (from auth profiles or
\texttt{OPENROUTER\_API\_KEY}). Without a key, use \texttt{--no-probe} to list candidates only.

Scan results are ranked by:

\begin{enumerate}
  \item Image support
  \item Tool latency
  \item Context size
  \item Parameter count
\end{enumerate}


Input

\begin{itemize}
  \item OpenRouter \texttt{/models} list (filter \texttt{:free})
  \item Requires OpenRouter API key from auth profiles or \texttt{OPENROUTER\_API\_KEY} (see \href{/help/environment}{/environment})
  \item Optional filters: \texttt{--max-age-days}, \texttt{--min-params}, \texttt{--provider}, \texttt{--max-candidates}
  \item Probe controls: \texttt{--timeout}, \texttt{--concurrency}
\end{itemize}


When run in a TTY, you can select fallbacks interactively. In non‑interactive
mode, pass \texttt{--yes} to accept defaults.

\subsection{Models registry (models.json)}


Custom providers in \texttt{models.providers} are written into \texttt{models.json} under the
agent directory (default \texttt{\textasciitilde{}/.openclaw/agents//models.json}). This file
is merged by default unless \texttt{models.mode} is set to \texttt{replace}.

Merge mode precedence for matching provider IDs:

\begin{itemize}
  \item Non-empty \texttt{apiKey}/\texttt{baseUrl} already present in the agent \texttt{models.json} win.
  \item Empty or missing agent \texttt{apiKey}/\texttt{baseUrl} fall back to config \texttt{models.providers}.
  \item Other provider fields are refreshed from config and normalized catalog data.
\end{itemize}



\clearpage

% === Source file: concepts/multi-agent.md ===
\section{Multi-Agent Routing}


Goal: multiple \_isolated\_ agents (separate workspace + \texttt{agentDir} + sessions), plus multiple channel accounts (e.g. two WhatsApps) in one running Gateway. Inbound is routed to an agent via bindings.

\subsection{What is “one agent”?}


An \textbf{agent} is a fully scoped brain with its own:

\begin{itemize}
  \item \textbf{Workspace} (files, AGENTS.md/SOUL.md/USER.md, local notes, persona rules).
  \item \textbf{State directory} (\texttt{agentDir}) for auth profiles, model registry, and per-agent config.
  \item \textbf{Session store} (chat history + routing state) under \texttt{\textasciitilde{}/.openclaw/agents//sessions}.
\end{itemize}


Auth profiles are \textbf{per-agent}. Each agent reads from its own:

\begin{lstlisting}
~/.openclaw/agents/<agentId>/agent/auth-profiles.json
\end{lstlisting}


Main agent credentials are \textbf{not} shared automatically. Never reuse \texttt{agentDir}
across agents (it causes auth/session collisions). If you want to share creds,
copy \texttt{auth-profiles.json} into the other agent's \texttt{agentDir}.

Skills are per-agent via each workspace’s \texttt{skills/} folder, with shared skills
available from \texttt{\textasciitilde{}/.openclaw/skills}. See \href{/tools/skills\#per-agent-vs-shared-skills}{Skills: per-agent vs shared}.

The Gateway can host \textbf{one agent} (default) or \textbf{many agents} side-by-side.

\textbf{Workspace note:} each agent’s workspace is the \textbf{default cwd}, not a hard
sandbox. Relative paths resolve inside the workspace, but absolute paths can
reach other host locations unless sandboxing is enabled. See
\href{/gateway/sandboxing}{Sandboxing}.

\subsection{Paths (quick map)}


\begin{itemize}
  \item Config: \texttt{\textasciitilde{}/.openclaw/openclaw.json} (or \texttt{OPENCLAW\_CONFIG\_PATH})
  \item State dir: \texttt{\textasciitilde{}/.openclaw} (or \texttt{OPENCLAW\_STATE\_DIR})
  \item Workspace: \texttt{\textasciitilde{}/.openclaw/workspace} (or \texttt{\textasciitilde{}/.openclaw/workspace-})
  \item Agent dir: \texttt{\textasciitilde{}/.openclaw/agents//agent} (or \texttt{agents.list[].agentDir})
  \item Sessions: \texttt{\textasciitilde{}/.openclaw/agents//sessions}
\end{itemize}


\subsubsection{Single-agent mode (default)}


If you do nothing, OpenClaw runs a single agent:

\begin{itemize}
  \item \texttt{agentId} defaults to \textbf{\texttt{main}}.
  \item Sessions are keyed as \texttt{agent:main:}.
  \item Workspace defaults to \texttt{\textasciitilde{}/.openclaw/workspace} (or \texttt{\textasciitilde{}/.openclaw/workspace-} when \texttt{OPENCLAW\_PROFILE} is set).
  \item State defaults to \texttt{\textasciitilde{}/.openclaw/agents/main/agent}.
\end{itemize}


\subsection{Agent helper}


Use the agent wizard to add a new isolated agent:

\begin{lstlisting}[language=bash]
openclaw agents add work
\end{lstlisting}


Then add \texttt{bindings} (or let the wizard do it) to route inbound messages.

Verify with:

\begin{lstlisting}[language=bash]
openclaw agents list --bindings
\end{lstlisting}


\subsection{Quick start}



Use the wizard or create workspaces manually:

\begin{lstlisting}[language=bash]
openclaw agents add coding
openclaw agents add social
\end{lstlisting}


Each agent gets its own workspace with \texttt{SOUL.md}, \texttt{AGENTS.md}, and optional \texttt{USER.md}, plus a dedicated \texttt{agentDir} and session store under \texttt{\textasciitilde{}/.openclaw/agents/}.



Create one account per agent on your preferred channels:

\begin{itemize}
  \item Discord: one bot per agent, enable Message Content Intent, copy each token.
  \item Telegram: one bot per agent via BotFather, copy each token.
  \item WhatsApp: link each phone number per account.
\end{itemize}


\begin{lstlisting}[language=bash]
openclaw channels login --channel whatsapp --account work
\end{lstlisting}


See channel guides: \href{/channels/discord}{Discord}, \href{/channels/telegram}{Telegram}, \href{/channels/whatsapp}{WhatsApp}.



Add agents under \texttt{agents.list}, channel accounts under \texttt{channels..accounts}, and connect them with \texttt{bindings} (examples below).



\begin{lstlisting}[language=bash]
openclaw gateway restart
openclaw agents list --bindings
openclaw channels status --probe
\end{lstlisting}



\subsection{Multiple agents = multiple people, multiple personalities}


With \textbf{multiple agents}, each \texttt{agentId} becomes a \textbf{fully isolated persona}:

\begin{itemize}
  \item \textbf{Different phone numbers/accounts} (per channel \texttt{accountId}).
  \item \textbf{Different personalities} (per-agent workspace files like \texttt{AGENTS.md} and \texttt{SOUL.md}).
  \item \textbf{Separate auth + sessions} (no cross-talk unless explicitly enabled).
\end{itemize}


This lets \textbf{multiple people} share one Gateway server while keeping their AI “brains” and data isolated.

\subsection{One WhatsApp number, multiple people (DM split)}


You can route \textbf{different WhatsApp DMs} to different agents while staying on \textbf{one WhatsApp account}. Match on sender E.164 (like \texttt{+15551234567}) with \texttt{peer.kind: "direct"}. Replies still come from the same WhatsApp number (no per‑agent sender identity).

Important detail: direct chats collapse to the agent’s \textbf{main session key}, so true isolation requires \textbf{one agent per person}.

Example:

\begin{lstlisting}[language=json]
{
  agents: {
    list: [
      { id: "alex", workspace: "~/.openclaw/workspace-alex" },
      { id: "mia", workspace: "~/.openclaw/workspace-mia" },
    ],
  },
  bindings: [
    {
      agentId: "alex",
      match: { channel: "whatsapp", peer: { kind: "direct", id: "+15551230001" } },
    },
    {
      agentId: "mia",
      match: { channel: "whatsapp", peer: { kind: "direct", id: "+15551230002" } },
    },
  ],
  channels: {
    whatsapp: {
      dmPolicy: "allowlist",
      allowFrom: ["+15551230001", "+15551230002"],
    },
  },
}
\end{lstlisting}


Notes:

\begin{itemize}
  \item DM access control is \textbf{global per WhatsApp account} (pairing/allowlist), not per agent.
  \item For shared groups, bind the group to one agent or use \href{/channels/broadcast-groups}{Broadcast groups}.
\end{itemize}


\subsection{Routing rules (how messages pick an agent)}


Bindings are \textbf{deterministic} and \textbf{most-specific wins}:

\begin{enumerate}
  \item \texttt{peer} match (exact DM/group/channel id)
  \item \texttt{parentPeer} match (thread inheritance)
  \item \texttt{guildId + roles} (Discord role routing)
  \item \texttt{guildId} (Discord)
  \item \texttt{teamId} (Slack)
  \item \texttt{accountId} match for a channel
  \item channel-level match (\texttt{accountId: "*"})
  \item fallback to default agent (\texttt{agents.list[].default}, else first list entry, default: \texttt{main})
\end{enumerate}


If multiple bindings match in the same tier, the first one in config order wins.
If a binding sets multiple match fields (for example \texttt{peer} + \texttt{guildId}), all specified fields are required (\texttt{AND} semantics).

Important account-scope detail:

\begin{itemize}
  \item A binding that omits \texttt{accountId} matches the default account only.
  \item Use \texttt{accountId: "*"} for a channel-wide fallback across all accounts.
  \item If you later add the same binding for the same agent with an explicit account id, OpenClaw upgrades the existing channel-only binding to account-scoped instead of duplicating it.
\end{itemize}


\subsection{Multiple accounts / phone numbers}


Channels that support \textbf{multiple accounts} (e.g. WhatsApp) use \texttt{accountId} to identify
each login. Each \texttt{accountId} can be routed to a different agent, so one server can host
multiple phone numbers without mixing sessions.

If you want a channel-wide default account when \texttt{accountId} is omitted, set
\texttt{channels..defaultAccount} (optional). When unset, OpenClaw falls back
to \texttt{default} if present, otherwise the first configured account id (sorted).

Common channels supporting this pattern include:

\begin{itemize}
  \item \texttt{whatsapp}, \texttt{telegram}, \texttt{discord}, \texttt{slack}, \texttt{signal}, \texttt{imessage}
  \item \texttt{irc}, \texttt{line}, \texttt{googlechat}, \texttt{mattermost}, \texttt{matrix}, \texttt{nextcloud-talk}
  \item \texttt{bluebubbles}, \texttt{zalo}, \texttt{zalouser}, \texttt{nostr}, \texttt{feishu}
\end{itemize}


\subsection{Concepts}


\begin{itemize}
  \item \texttt{agentId}: one “brain” (workspace, per-agent auth, per-agent session store).
  \item \texttt{accountId}: one channel account instance (e.g. WhatsApp account \texttt{"personal"} vs \texttt{"biz"}).
  \item \texttt{binding}: routes inbound messages to an \texttt{agentId} by \texttt{(channel, accountId, peer)} and optionally guild/team ids.
  \item Direct chats collapse to \texttt{agent::} (per-agent “main”; \texttt{session.mainKey}).
\end{itemize}


\subsection{Platform examples}


\subsubsection{Discord bots per agent}


Each Discord bot account maps to a unique \texttt{accountId}. Bind each account to an agent and keep allowlists per bot.

\begin{lstlisting}[language=json]
{
  agents: {
    list: [
      { id: "main", workspace: "~/.openclaw/workspace-main" },
      { id: "coding", workspace: "~/.openclaw/workspace-coding" },
    ],
  },
  bindings: [
    { agentId: "main", match: { channel: "discord", accountId: "default" } },
    { agentId: "coding", match: { channel: "discord", accountId: "coding" } },
  ],
  channels: {
    discord: {
      groupPolicy: "allowlist",
      accounts: {
        default: {
          token: "DISCORD_BOT_TOKEN_MAIN",
          guilds: {
            "123456789012345678": {
              channels: {
                "222222222222222222": { allow: true, requireMention: false },
              },
            },
          },
        },
        coding: {
          token: "DISCORD_BOT_TOKEN_CODING",
          guilds: {
            "123456789012345678": {
              channels: {
                "333333333333333333": { allow: true, requireMention: false },
              },
            },
          },
        },
      },
    },
  },
}
\end{lstlisting}


Notes:

\begin{itemize}
  \item Invite each bot to the guild and enable Message Content Intent.
  \item Tokens live in \texttt{channels.discord.accounts..token} (default account can use \texttt{DISCORD\_BOT\_TOKEN}).
\end{itemize}


\subsubsection{Telegram bots per agent}


\begin{lstlisting}[language=json]
{
  agents: {
    list: [
      { id: "main", workspace: "~/.openclaw/workspace-main" },
      { id: "alerts", workspace: "~/.openclaw/workspace-alerts" },
    ],
  },
  bindings: [
    { agentId: "main", match: { channel: "telegram", accountId: "default" } },
    { agentId: "alerts", match: { channel: "telegram", accountId: "alerts" } },
  ],
  channels: {
    telegram: {
      accounts: {
        default: {
          botToken: "123456:ABC...",
          dmPolicy: "pairing",
        },
        alerts: {
          botToken: "987654:XYZ...",
          dmPolicy: "allowlist",
          allowFrom: ["tg:123456789"],
        },
      },
    },
  },
}
\end{lstlisting}


Notes:

\begin{itemize}
  \item Create one bot per agent with BotFather and copy each token.
  \item Tokens live in \texttt{channels.telegram.accounts..botToken} (default account can use \texttt{TELEGRAM\_BOT\_TOKEN}).
\end{itemize}


\subsubsection{WhatsApp numbers per agent}


Link each account before starting the gateway:

\begin{lstlisting}[language=bash]
openclaw channels login --channel whatsapp --account personal
openclaw channels login --channel whatsapp --account biz
\end{lstlisting}


\texttt{\textasciitilde{}/.openclaw/openclaw.json} (JSON5):

\begin{lstlisting}[language=javascript]
{
  agents: {
    list: [
      {
        id: "home",
        default: true,
        name: "Home",
        workspace: "~/.openclaw/workspace-home",
        agentDir: "~/.openclaw/agents/home/agent",
      },
      {
        id: "work",
        name: "Work",
        workspace: "~/.openclaw/workspace-work",
        agentDir: "~/.openclaw/agents/work/agent",
      },
    ],
  },

  // Deterministic routing: first match wins (most-specific first).
  bindings: [
    { agentId: "home", match: { channel: "whatsapp", accountId: "personal" } },
    { agentId: "work", match: { channel: "whatsapp", accountId: "biz" } },

    // Optional per-peer override (example: send a specific group to work agent).
    {
      agentId: "work",
      match: {
        channel: "whatsapp",
        accountId: "personal",
        peer: { kind: "group", id: "1203630...@g.us" },
      },
    },
  ],

  // Off by default: agent-to-agent messaging must be explicitly enabled + allowlisted.
  tools: {
    agentToAgent: {
      enabled: false,
      allow: ["home", "work"],
    },
  },

  channels: {
    whatsapp: {
      accounts: {
        personal: {
          // Optional override. Default: ~/.openclaw/credentials/whatsapp/personal
          // authDir: "~/.openclaw/credentials/whatsapp/personal",
        },
        biz: {
          // Optional override. Default: ~/.openclaw/credentials/whatsapp/biz
          // authDir: "~/.openclaw/credentials/whatsapp/biz",
        },
      },
    },
  },
}
\end{lstlisting}


\subsection{Example: WhatsApp daily chat + Telegram deep work}


Split by channel: route WhatsApp to a fast everyday agent and Telegram to an Opus agent.

\begin{lstlisting}[language=json]
{
  agents: {
    list: [
      {
        id: "chat",
        name: "Everyday",
        workspace: "~/.openclaw/workspace-chat",
        model: "anthropic/claude-sonnet-4-5",
      },
      {
        id: "opus",
        name: "Deep Work",
        workspace: "~/.openclaw/workspace-opus",
        model: "anthropic/claude-opus-4-6",
      },
    ],
  },
  bindings: [
    { agentId: "chat", match: { channel: "whatsapp" } },
    { agentId: "opus", match: { channel: "telegram" } },
  ],
}
\end{lstlisting}


Notes:

\begin{itemize}
  \item If you have multiple accounts for a channel, add \texttt{accountId} to the binding (for example \texttt{\{ channel: "whatsapp", accountId: "personal" \}}).
  \item To route a single DM/group to Opus while keeping the rest on chat, add a \texttt{match.peer} binding for that peer; peer matches always win over channel-wide rules.
\end{itemize}


\subsection{Example: same channel, one peer to Opus}


Keep WhatsApp on the fast agent, but route one DM to Opus:

\begin{lstlisting}[language=json]
{
  agents: {
    list: [
      {
        id: "chat",
        name: "Everyday",
        workspace: "~/.openclaw/workspace-chat",
        model: "anthropic/claude-sonnet-4-5",
      },
      {
        id: "opus",
        name: "Deep Work",
        workspace: "~/.openclaw/workspace-opus",
        model: "anthropic/claude-opus-4-6",
      },
    ],
  },
  bindings: [
    {
      agentId: "opus",
      match: { channel: "whatsapp", peer: { kind: "direct", id: "+15551234567" } },
    },
    { agentId: "chat", match: { channel: "whatsapp" } },
  ],
}
\end{lstlisting}


Peer bindings always win, so keep them above the channel-wide rule.

\subsection{Family agent bound to a WhatsApp group}


Bind a dedicated family agent to a single WhatsApp group, with mention gating
and a tighter tool policy:

\begin{lstlisting}[language=json]
{
  agents: {
    list: [
      {
        id: "family",
        name: "Family",
        workspace: "~/.openclaw/workspace-family",
        identity: { name: "Family Bot" },
        groupChat: {
          mentionPatterns: ["@family", "@familybot", "@Family Bot"],
        },
        sandbox: {
          mode: "all",
          scope: "agent",
        },
        tools: {
          allow: [
            "exec",
            "read",
            "sessions_list",
            "sessions_history",
            "sessions_send",
            "sessions_spawn",
            "session_status",
          ],
          deny: ["write", "edit", "apply_patch", "browser", "canvas", "nodes", "cron"],
        },
      },
    ],
  },
  bindings: [
    {
      agentId: "family",
      match: {
        channel: "whatsapp",
        peer: { kind: "group", id: "120363999999999999@g.us" },
      },
    },
  ],
}
\end{lstlisting}


Notes:

\begin{itemize}
  \item Tool allow/deny lists are \textbf{tools}, not skills. If a skill needs to run a
\end{itemize}

binary, ensure \texttt{exec} is allowed and the binary exists in the sandbox.
\begin{itemize}
  \item For stricter gating, set \texttt{agents.list[].groupChat.mentionPatterns} and keep
\end{itemize}

group allowlists enabled for the channel.

\subsection{Per-Agent Sandbox and Tool Configuration}


Starting with v2026.1.6, each agent can have its own sandbox and tool restrictions:

\begin{lstlisting}[language=javascript]
{
  agents: {
    list: [
      {
        id: "personal",
        workspace: "~/.openclaw/workspace-personal",
        sandbox: {
          mode: "off",  // No sandbox for personal agent
        },
        // No tool restrictions - all tools available
      },
      {
        id: "family",
        workspace: "~/.openclaw/workspace-family",
        sandbox: {
          mode: "all",     // Always sandboxed
          scope: "agent",  // One container per agent
          docker: {
            // Optional one-time setup after container creation
            setupCommand: "apt-get update && apt-get install -y git curl",
          },
        },
        tools: {
          allow: ["read"],                    // Only read tool
          deny: ["exec", "write", "edit", "apply_patch"],    // Deny others
        },
      },
    ],
  },
}
\end{lstlisting}


Note: \texttt{setupCommand} lives under \texttt{sandbox.docker} and runs once on container creation.
Per-agent \texttt{sandbox.docker.*} overrides are ignored when the resolved scope is \texttt{"shared"}.

\textbf{Benefits:}

\begin{itemize}
  \item \textbf{Security isolation}: Restrict tools for untrusted agents
  \item \textbf{Resource control}: Sandbox specific agents while keeping others on host
  \item \textbf{Flexible policies}: Different permissions per agent
\end{itemize}


Note: \texttt{tools.elevated} is \textbf{global} and sender-based; it is not configurable per agent.
If you need per-agent boundaries, use \texttt{agents.list[].tools} to deny \texttt{exec}.
For group targeting, use \texttt{agents.list[].groupChat.mentionPatterns} so @mentions map cleanly to the intended agent.

See \href{/tools/multi-agent-sandbox-tools}{Multi-Agent Sandbox \& Tools} for detailed examples.


\clearpage

% === Source file: concepts/oauth.md ===
\section{OAuth}


OpenClaw supports “subscription auth” via OAuth for providers that offer it (notably \textbf{OpenAI Codex (ChatGPT OAuth)}). For Anthropic subscriptions, use the \textbf{setup-token} flow. Anthropic subscription use outside Claude Code has been restricted for some users in the past, so treat it as a user-choice risk and verify current Anthropic policy yourself. OpenAI Codex OAuth is explicitly supported for use in external tools like OpenClaw. This page explains:

For Anthropic in production, API key auth is the safer recommended path over subscription setup-token auth.

\begin{itemize}
  \item how the OAuth \textbf{token exchange} works (PKCE)
  \item where tokens are \textbf{stored} (and why)
  \item how to handle \textbf{multiple accounts} (profiles + per-session overrides)
\end{itemize}


OpenClaw also supports \textbf{provider plugins} that ship their own OAuth or API‑key
flows. Run them via:

\begin{lstlisting}[language=bash]
openclaw models auth login --provider <id>
\end{lstlisting}


\subsection{The token sink (why it exists)}


OAuth providers commonly mint a \textbf{new refresh token} during login/refresh flows. Some providers (or OAuth clients) can invalidate older refresh tokens when a new one is issued for the same user/app.

Practical symptom:

\begin{itemize}
  \item you log in via OpenClaw \_and\_ via Claude Code / Codex CLI → one of them randomly gets “logged out” later
\end{itemize}


To reduce that, OpenClaw treats \texttt{auth-profiles.json} as a \textbf{token sink}:

\begin{itemize}
  \item the runtime reads credentials from \textbf{one place}
  \item we can keep multiple profiles and route them deterministically
\end{itemize}


\subsection{Storage (where tokens live)}


Secrets are stored \textbf{per-agent}:

\begin{itemize}
  \item Auth profiles (OAuth + API keys + optional value-level refs): \texttt{\textasciitilde{}/.openclaw/agents//agent/auth-profiles.json}
  \item Legacy compatibility file: \texttt{\textasciitilde{}/.openclaw/agents//agent/auth.json}
\end{itemize}

(static \texttt{api\_key} entries are scrubbed when discovered)

Legacy import-only file (still supported, but not the main store):

\begin{itemize}
  \item \texttt{\textasciitilde{}/.openclaw/credentials/oauth.json} (imported into \texttt{auth-profiles.json} on first use)
\end{itemize}


All of the above also respect \texttt{\$OPENCLAW\_STATE\_DIR} (state dir override). Full reference: \href{/gateway/configuration\#auth-storage-oauth--api-keys}{/gateway/configuration}

For static secret refs and runtime snapshot activation behavior, see \href{/gateway/secrets}{Secrets Management}.

\subsection{Anthropic setup-token (subscription auth)}


Anthropic setup-token support is technical compatibility, not a policy guarantee.
Anthropic has blocked some subscription usage outside Claude Code in the past.
Decide for yourself whether to use subscription auth, and verify Anthropic's current terms.

Run \texttt{claude setup-token} on any machine, then paste it into OpenClaw:

\begin{lstlisting}[language=bash]
openclaw models auth setup-token --provider anthropic
\end{lstlisting}


If you generated the token elsewhere, paste it manually:

\begin{lstlisting}[language=bash]
openclaw models auth paste-token --provider anthropic
\end{lstlisting}


Verify:

\begin{lstlisting}[language=bash]
openclaw models status
\end{lstlisting}


\subsection{OAuth exchange (how login works)}


OpenClaw’s interactive login flows are implemented in \texttt{@mariozechner/pi-ai} and wired into the wizards/commands.

\subsubsection{Anthropic setup-token}


Flow shape:

\begin{enumerate}
  \item run \texttt{claude setup-token}
  \item paste the token into OpenClaw
  \item store as a token auth profile (no refresh)
\end{enumerate}


The wizard path is \texttt{openclaw onboard} → auth choice \texttt{setup-token} (Anthropic).

\subsubsection{OpenAI Codex (ChatGPT OAuth)}


OpenAI Codex OAuth is explicitly supported for use outside the Codex CLI, including OpenClaw workflows.

Flow shape (PKCE):

\begin{enumerate}
  \item generate PKCE verifier/challenge + random \texttt{state}
  \item open \texttt{https://auth.openai.com/oauth/authorize?...}
  \item try to capture callback on \texttt{http://127.0.0.1:1455/auth/callback}
  \item if callback can’t bind (or you’re remote/headless), paste the redirect URL/code
  \item exchange at \texttt{https://auth.openai.com/oauth/token}
  \item extract \texttt{accountId} from the access token and store \texttt{\{ access, refresh, expires, accountId \}}
\end{enumerate}


Wizard path is \texttt{openclaw onboard} → auth choice \texttt{openai-codex}.

\subsection{Refresh + expiry}


Profiles store an \texttt{expires} timestamp.

At runtime:

\begin{itemize}
  \item if \texttt{expires} is in the future → use the stored access token
  \item if expired → refresh (under a file lock) and overwrite the stored credentials
\end{itemize}


The refresh flow is automatic; you generally don't need to manage tokens manually.

\subsection{Multiple accounts (profiles) + routing}


Two patterns:

\subsubsection{1) Preferred: separate agents}


If you want “personal” and “work” to never interact, use isolated agents (separate sessions + credentials + workspace):

\begin{lstlisting}[language=bash]
openclaw agents add work
openclaw agents add personal
\end{lstlisting}


Then configure auth per-agent (wizard) and route chats to the right agent.

\subsubsection{2) Advanced: multiple profiles in one agent}


\texttt{auth-profiles.json} supports multiple profile IDs for the same provider.

Pick which profile is used:

\begin{itemize}
  \item globally via config ordering (\texttt{auth.order})
  \item per-session via \texttt{/model ...@}
\end{itemize}


Example (session override):

\begin{itemize}
  \item \texttt{/model Opus@anthropic:work}
\end{itemize}


How to see what profile IDs exist:

\begin{itemize}
  \item \texttt{openclaw channels list --json} (shows \texttt{auth[]})
\end{itemize}


Related docs:

\begin{itemize}
  \item \href{/concepts/model-failover}{/concepts/model-failover} (rotation + cooldown rules)
  \item \href{/tools/slash-commands}{/tools/slash-commands} (command surface)
\end{itemize}



\clearpage

% === Source file: concepts/presence.md ===
\section{Presence}


OpenClaw “presence” is a lightweight, best‑effort view of:

\begin{itemize}
  \item the \textbf{Gateway} itself, and
  \item \textbf{clients connected to the Gateway} (mac app, WebChat, CLI, etc.)
\end{itemize}


Presence is used primarily to render the macOS app’s \textbf{Instances} tab and to
provide quick operator visibility.

\subsection{Presence fields (what shows up)}


Presence entries are structured objects with fields like:

\begin{itemize}
  \item \texttt{instanceId} (optional but strongly recommended): stable client identity (usually \texttt{connect.client.instanceId})
  \item \texttt{host}: human‑friendly host name
  \item \texttt{ip}: best‑effort IP address
  \item \texttt{version}: client version string
  \item \texttt{deviceFamily} / \texttt{modelIdentifier}: hardware hints
  \item \texttt{mode}: \texttt{ui}, \texttt{webchat}, \texttt{cli}, \texttt{backend}, \texttt{probe}, \texttt{test}, \texttt{node}, ...
  \item \texttt{lastInputSeconds}: “seconds since last user input” (if known)
  \item \texttt{reason}: \texttt{self}, \texttt{connect}, \texttt{node-connected}, \texttt{periodic}, ...
  \item \texttt{ts}: last update timestamp (ms since epoch)
\end{itemize}


\subsection{Producers (where presence comes from)}


Presence entries are produced by multiple sources and \textbf{merged}.

\subsubsection{1) Gateway self entry}


The Gateway always seeds a “self” entry at startup so UIs show the gateway host
even before any clients connect.

\subsubsection{2) WebSocket connect}


Every WS client begins with a \texttt{connect} request. On successful handshake the
Gateway upserts a presence entry for that connection.

\paragraph{Why one‑off CLI commands don’t show up}


The CLI often connects for short, one‑off commands. To avoid spamming the
Instances list, \texttt{client.mode === "cli"} is \textbf{not} turned into a presence entry.

\subsubsection{3) system-event beacons}


Clients can send richer periodic beacons via the \texttt{system-event} method. The mac
app uses this to report host name, IP, and \texttt{lastInputSeconds}.

\subsubsection{4) Node connects (role: node)}


When a node connects over the Gateway WebSocket with \texttt{role: node}, the Gateway
upserts a presence entry for that node (same flow as other WS clients).

\subsection{Merge + dedupe rules (why instanceId matters)}


Presence entries are stored in a single in‑memory map:

\begin{itemize}
  \item Entries are keyed by a \textbf{presence key}.
  \item The best key is a stable \texttt{instanceId} (from \texttt{connect.client.instanceId}) that survives restarts.
  \item Keys are case‑insensitive.
\end{itemize}


If a client reconnects without a stable \texttt{instanceId}, it may show up as a
\textbf{duplicate} row.

\subsection{TTL and bounded size}


Presence is intentionally ephemeral:

\begin{itemize}
  \item \textbf{TTL:} entries older than 5 minutes are pruned
  \item \textbf{Max entries:} 200 (oldest dropped first)
\end{itemize}


This keeps the list fresh and avoids unbounded memory growth.

\subsection{Remote/tunnel caveat (loopback IPs)}


When a client connects over an SSH tunnel / local port forward, the Gateway may
see the remote address as \texttt{127.0.0.1}. To avoid overwriting a good client‑reported
IP, loopback remote addresses are ignored.

\subsection{Consumers}


\subsubsection{macOS Instances tab}


The macOS app renders the output of \texttt{system-presence} and applies a small status
indicator (Active/Idle/Stale) based on the age of the last update.

\subsection{Debugging tips}


\begin{itemize}
  \item To see the raw list, call \texttt{system-presence} against the Gateway.
  \item If you see duplicates:
  \item confirm clients send a stable \texttt{client.instanceId} in the handshake
  \item confirm periodic beacons use the same \texttt{instanceId}
  \item check whether the connection‑derived entry is missing \texttt{instanceId} (duplicates are expected)
\end{itemize}



\clearpage

% === Source file: concepts/queue.md ===
\section{Command Queue (2026-01-16)}


We serialize inbound auto-reply runs (all channels) through a tiny in-process queue to prevent multiple agent runs from colliding, while still allowing safe parallelism across sessions.

\subsection{Why}


\begin{itemize}
  \item Auto-reply runs can be expensive (LLM calls) and can collide when multiple inbound messages arrive close together.
  \item Serializing avoids competing for shared resources (session files, logs, CLI stdin) and reduces the chance of upstream rate limits.
\end{itemize}


\subsection{How it works}


\begin{itemize}
  \item A lane-aware FIFO queue drains each lane with a configurable concurrency cap (default 1 for unconfigured lanes; main defaults to 4, subagent to 8).
  \item \texttt{runEmbeddedPiAgent} enqueues by \textbf{session key} (lane \texttt{session:}) to guarantee only one active run per session.
  \item Each session run is then queued into a \textbf{global lane} (\texttt{main} by default) so overall parallelism is capped by \texttt{agents.defaults.maxConcurrent}.
  \item When verbose logging is enabled, queued runs emit a short notice if they waited more than \textasciitilde{}2s before starting.
  \item Typing indicators still fire immediately on enqueue (when supported by the channel) so user experience is unchanged while we wait our turn.
\end{itemize}


\subsection{Queue modes (per channel)}


Inbound messages can steer the current run, wait for a followup turn, or do both:

\begin{itemize}
  \item \texttt{steer}: inject immediately into the current run (cancels pending tool calls after the next tool boundary). If not streaming, falls back to followup.
  \item \texttt{followup}: enqueue for the next agent turn after the current run ends.
  \item \texttt{collect}: coalesce all queued messages into a \textbf{single} followup turn (default). If messages target different channels/threads, they drain individually to preserve routing.
  \item \texttt{steer-backlog} (aka \texttt{steer+backlog}): steer now \textbf{and} preserve the message for a followup turn.
  \item \texttt{interrupt} (legacy): abort the active run for that session, then run the newest message.
  \item \texttt{queue} (legacy alias): same as \texttt{steer}.
\end{itemize}


Steer-backlog means you can get a followup response after the steered run, so
streaming surfaces can look like duplicates. Prefer \texttt{collect}/\texttt{steer} if you want
one response per inbound message.
Send \texttt{/queue collect} as a standalone command (per-session) or set \texttt{messages.queue.byChannel.discord: "collect"}.

Defaults (when unset in config):

\begin{itemize}
  \item All surfaces → \texttt{collect}
\end{itemize}


Configure globally or per channel via \texttt{messages.queue}:

\begin{lstlisting}[language=json]
{
  messages: {
    queue: {
      mode: "collect",
      debounceMs: 1000,
      cap: 20,
      drop: "summarize",
      byChannel: { discord: "collect" },
    },
  },
}
\end{lstlisting}


\subsection{Queue options}


Options apply to \texttt{followup}, \texttt{collect}, and \texttt{steer-backlog} (and to \texttt{steer} when it falls back to followup):

\begin{itemize}
  \item \texttt{debounceMs}: wait for quiet before starting a followup turn (prevents “continue, continue”).
  \item \texttt{cap}: max queued messages per session.
  \item \texttt{drop}: overflow policy (\texttt{old}, \texttt{new}, \texttt{summarize}).
\end{itemize}


Summarize keeps a short bullet list of dropped messages and injects it as a synthetic followup prompt.
Defaults: \texttt{debounceMs: 1000}, \texttt{cap: 20}, \texttt{drop: summarize}.

\subsection{Per-session overrides}


\begin{itemize}
  \item Send \texttt{/queue } as a standalone command to store the mode for the current session.
  \item Options can be combined: \texttt{/queue collect debounce:2s cap:25 drop:summarize}
  \item \texttt{/queue default} or \texttt{/queue reset} clears the session override.
\end{itemize}


\subsection{Scope and guarantees}


\begin{itemize}
  \item Applies to auto-reply agent runs across all inbound channels that use the gateway reply pipeline (WhatsApp web, Telegram, Slack, Discord, Signal, iMessage, webchat, etc.).
  \item Default lane (\texttt{main}) is process-wide for inbound + main heartbeats; set \texttt{agents.defaults.maxConcurrent} to allow multiple sessions in parallel.
  \item Additional lanes may exist (e.g. \texttt{cron}, \texttt{subagent}) so background jobs can run in parallel without blocking inbound replies.
  \item Per-session lanes guarantee that only one agent run touches a given session at a time.
  \item No external dependencies or background worker threads; pure TypeScript + promises.
\end{itemize}


\subsection{Troubleshooting}


\begin{itemize}
  \item If commands seem stuck, enable verbose logs and look for “queued for …ms” lines to confirm the queue is draining.
  \item If you need queue depth, enable verbose logs and watch for queue timing lines.
\end{itemize}



\clearpage

% === Source file: concepts/retry.md ===
\section{Retry policy}


\subsection{Goals}


\begin{itemize}
  \item Retry per HTTP request, not per multi-step flow.
  \item Preserve ordering by retrying only the current step.
  \item Avoid duplicating non-idempotent operations.
\end{itemize}


\subsection{Defaults}


\begin{itemize}
  \item Attempts: 3
  \item Max delay cap: 30000 ms
  \item Jitter: 0.1 (10 percent)
  \item Provider defaults:
  \item Telegram min delay: 400 ms
  \item Discord min delay: 500 ms
\end{itemize}


\subsection{Behavior}


\subsubsection{Discord}


\begin{itemize}
  \item Retries only on rate-limit errors (HTTP 429).
  \item Uses Discord \texttt{retry\_after} when available, otherwise exponential backoff.
\end{itemize}


\subsubsection{Telegram}


\begin{itemize}
  \item Retries on transient errors (429, timeout, connect/reset/closed, temporarily unavailable).
  \item Uses \texttt{retry\_after} when available, otherwise exponential backoff.
  \item Markdown parse errors are not retried; they fall back to plain text.
\end{itemize}


\subsection{Configuration}


Set retry policy per provider in \texttt{\textasciitilde{}/.openclaw/openclaw.json}:

\begin{lstlisting}[language=json]
{
  channels: {
    telegram: {
      retry: {
        attempts: 3,
        minDelayMs: 400,
        maxDelayMs: 30000,
        jitter: 0.1,
      },
    },
    discord: {
      retry: {
        attempts: 3,
        minDelayMs: 500,
        maxDelayMs: 30000,
        jitter: 0.1,
      },
    },
  },
}
\end{lstlisting}


\subsection{Notes}


\begin{itemize}
  \item Retries apply per request (message send, media upload, reaction, poll, sticker).
  \item Composite flows do not retry completed steps.
\end{itemize}



\clearpage

% === Source file: concepts/session-pruning.md ===
\section{Session Pruning}


Session pruning trims \textbf{old tool results} from the in-memory context right before each LLM call. It does \textbf{not} rewrite the on-disk session history (\texttt{*.jsonl}).

\subsection{When it runs}


\begin{itemize}
  \item When \texttt{mode: "cache-ttl"} is enabled and the last Anthropic call for the session is older than \texttt{ttl}.
  \item Only affects the messages sent to the model for that request.
  \item Only active for Anthropic API calls (and OpenRouter Anthropic models).
  \item For best results, match \texttt{ttl} to your model \texttt{cacheRetention} policy (\texttt{short} = 5m, \texttt{long} = 1h).
  \item After a prune, the TTL window resets so subsequent requests keep cache until \texttt{ttl} expires again.
\end{itemize}


\subsection{Smart defaults (Anthropic)}


\begin{itemize}
  \item \textbf{OAuth or setup-token} profiles: enable \texttt{cache-ttl} pruning and set heartbeat to \texttt{1h}.
  \item \textbf{API key} profiles: enable \texttt{cache-ttl} pruning, set heartbeat to \texttt{30m}, and default \texttt{cacheRetention: "short"} on Anthropic models.
  \item If you set any of these values explicitly, OpenClaw does \textbf{not} override them.
\end{itemize}


\subsection{What this improves (cost + cache behavior)}


\begin{itemize}
  \item \textbf{Why prune:} Anthropic prompt caching only applies within the TTL. If a session goes idle past the TTL, the next request re-caches the full prompt unless you trim it first.
  \item \textbf{What gets cheaper:} pruning reduces the \textbf{cacheWrite} size for that first request after the TTL expires.
  \item \textbf{Why the TTL reset matters:} once pruning runs, the cache window resets, so follow‑up requests can reuse the freshly cached prompt instead of re-caching the full history again.
  \item \textbf{What it does not do:} pruning doesn’t add tokens or “double” costs; it only changes what gets cached on that first post‑TTL request.
\end{itemize}


\subsection{What can be pruned}


\begin{itemize}
  \item Only \texttt{toolResult} messages.
  \item User + assistant messages are \textbf{never} modified.
  \item The last \texttt{keepLastAssistants} assistant messages are protected; tool results after that cutoff are not pruned.
  \item If there aren’t enough assistant messages to establish the cutoff, pruning is skipped.
  \item Tool results containing \textbf{image blocks} are skipped (never trimmed/cleared).
\end{itemize}


\subsection{Context window estimation}


Pruning uses an estimated context window (chars ≈ tokens × 4). The base window is resolved in this order:

\begin{enumerate}
  \item \texttt{models.providers.*.models[].contextWindow} override.
  \item Model definition \texttt{contextWindow} (from the model registry).
  \item Default \texttt{200000} tokens.
\end{enumerate}


If \texttt{agents.defaults.contextTokens} is set, it is treated as a cap (min) on the resolved window.

\subsection{Mode}


\subsubsection{cache-ttl}


\begin{itemize}
  \item Pruning only runs if the last Anthropic call is older than \texttt{ttl} (default \texttt{5m}).
  \item When it runs: same soft-trim + hard-clear behavior as before.
\end{itemize}


\subsection{Soft vs hard pruning}


\begin{itemize}
  \item \textbf{Soft-trim}: only for oversized tool results.
  \item Keeps head + tail, inserts \texttt{...}, and appends a note with the original size.
  \item Skips results with image blocks.
  \item \textbf{Hard-clear}: replaces the entire tool result with \texttt{hardClear.placeholder}.
\end{itemize}


\subsection{Tool selection}


\begin{itemize}
  \item \texttt{tools.allow} / \texttt{tools.deny} support \texttt{*} wildcards.
  \item Deny wins.
  \item Matching is case-insensitive.
  \item Empty allow list => all tools allowed.
\end{itemize}


\subsection{Interaction with other limits}


\begin{itemize}
  \item Built-in tools already truncate their own output; session pruning is an extra layer that prevents long-running chats from accumulating too much tool output in the model context.
  \item Compaction is separate: compaction summarizes and persists, pruning is transient per request. See \href{/concepts/compaction}{/concepts/compaction}.
\end{itemize}


\subsection{Defaults (when enabled)}


\begin{itemize}
  \item \texttt{ttl}: \texttt{"5m"}
  \item \texttt{keepLastAssistants}: \texttt{3}
  \item \texttt{softTrimRatio}: \texttt{0.3}
  \item \texttt{hardClearRatio}: \texttt{0.5}
  \item \texttt{minPrunableToolChars}: \texttt{50000}
  \item \texttt{softTrim}: \texttt{\{ maxChars: 4000, headChars: 1500, tailChars: 1500 \}}
  \item \texttt{hardClear}: \texttt{\{ enabled: true, placeholder: "[Old tool result content cleared]" \}}
\end{itemize}


\subsection{Examples}


Default (off):

\begin{lstlisting}[language=json]
{
  agents: { defaults: { contextPruning: { mode: "off" } } },
}
\end{lstlisting}


Enable TTL-aware pruning:

\begin{lstlisting}[language=json]
{
  agents: { defaults: { contextPruning: { mode: "cache-ttl", ttl: "5m" } } },
}
\end{lstlisting}


Restrict pruning to specific tools:

\begin{lstlisting}[language=json]
{
  agents: {
    defaults: {
      contextPruning: {
        mode: "cache-ttl",
        tools: { allow: ["exec", "read"], deny: ["*image*"] },
      },
    },
  },
}
\end{lstlisting}


See config reference: \href{/gateway/configuration}{Gateway Configuration}


\clearpage

% === Source file: concepts/session-tool.md ===
\section{Session Tools}


Goal: small, hard-to-misuse tool set so agents can list sessions, fetch history, and send to another session.

\subsection{Tool Names}


\begin{itemize}
  \item \texttt{sessions\_list}
  \item \texttt{sessions\_history}
  \item \texttt{sessions\_send}
  \item \texttt{sessions\_spawn}
\end{itemize}


\subsection{Key Model}


\begin{itemize}
  \item Main direct chat bucket is always the literal key \texttt{"main"} (resolved to the current agent’s main key).
  \item Group chats use \texttt{agent:::group:} or \texttt{agent:::channel:} (pass the full key).
  \item Cron jobs use \texttt{cron:}.
  \item Hooks use \texttt{hook:} unless explicitly set.
  \item Node sessions use \texttt{node-} unless explicitly set.
\end{itemize}


\texttt{global} and \texttt{unknown} are reserved values and are never listed. If \texttt{session.scope = "global"}, we alias it to \texttt{main} for all tools so callers never see \texttt{global}.

\subsection{sessions\_list}


List sessions as an array of rows.

Parameters:

\begin{itemize}
  \item \texttt{kinds?: string[]} filter: any of \texttt{"main" | "group" | "cron" | "hook" | "node" | "other"}
  \item \texttt{limit?: number} max rows (default: server default, clamp e.g. 200)
  \item \texttt{activeMinutes?: number} only sessions updated within N minutes
  \item \texttt{messageLimit?: number} 0 = no messages (default 0); >0 = include last N messages
\end{itemize}


Behavior:

\begin{itemize}
  \item \texttt{messageLimit > 0} fetches \texttt{chat.history} per session and includes the last N messages.
  \item Tool results are filtered out in list output; use \texttt{sessions\_history} for tool messages.
  \item When running in a \textbf{sandboxed} agent session, session tools default to \textbf{spawned-only visibility} (see below).
\end{itemize}


Row shape (JSON):

\begin{itemize}
  \item \texttt{key}: session key (string)
  \item \texttt{kind}: \texttt{main | group | cron | hook | node | other}
  \item \texttt{channel}: \texttt{whatsapp | telegram | discord | signal | imessage | webchat | internal | unknown}
  \item \texttt{displayName} (group display label if available)
  \item \texttt{updatedAt} (ms)
  \item \texttt{sessionId}
  \item \texttt{model}, \texttt{contextTokens}, \texttt{totalTokens}
  \item \texttt{thinkingLevel}, \texttt{verboseLevel}, \texttt{systemSent}, \texttt{abortedLastRun}
  \item \texttt{sendPolicy} (session override if set)
  \item \texttt{lastChannel}, \texttt{lastTo}
  \item \texttt{deliveryContext} (normalized \texttt{\{ channel, to, accountId \}} when available)
  \item \texttt{transcriptPath} (best-effort path derived from store dir + sessionId)
  \item \texttt{messages?} (only when \texttt{messageLimit > 0})
\end{itemize}


\subsection{sessions\_history}


Fetch transcript for one session.

Parameters:

\begin{itemize}
  \item \texttt{sessionKey} (required; accepts session key or \texttt{sessionId} from \texttt{sessions\_list})
  \item \texttt{limit?: number} max messages (server clamps)
  \item \texttt{includeTools?: boolean} (default false)
\end{itemize}


Behavior:

\begin{itemize}
  \item \texttt{includeTools=false} filters \texttt{role: "toolResult"} messages.
  \item Returns messages array in the raw transcript format.
  \item When given a \texttt{sessionId}, OpenClaw resolves it to the corresponding session key (missing ids error).
\end{itemize}


\subsection{sessions\_send}


Send a message into another session.

Parameters:

\begin{itemize}
  \item \texttt{sessionKey} (required; accepts session key or \texttt{sessionId} from \texttt{sessions\_list})
  \item \texttt{message} (required)
  \item \texttt{timeoutSeconds?: number} (default >0; 0 = fire-and-forget)
\end{itemize}


Behavior:

\begin{itemize}
  \item \texttt{timeoutSeconds = 0}: enqueue and return \texttt{\{ runId, status: "accepted" \}}.
  \item \texttt{timeoutSeconds > 0}: wait up to N seconds for completion, then return \texttt{\{ runId, status: "ok", reply \}}.
  \item If wait times out: \texttt{\{ runId, status: "timeout", error \}}. Run continues; call \texttt{sessions\_history} later.
  \item If the run fails: \texttt{\{ runId, status: "error", error \}}.
  \item Announce delivery runs after the primary run completes and is best-effort; \texttt{status: "ok"} does not guarantee the announce was delivered.
  \item Waits via gateway \texttt{agent.wait} (server-side) so reconnects don't drop the wait.
  \item Agent-to-agent message context is injected for the primary run.
  \item Inter-session messages are persisted with \texttt{message.provenance.kind = "inter\_session"} so transcript readers can distinguish routed agent instructions from external user input.
  \item After the primary run completes, OpenClaw runs a \textbf{reply-back loop}:
  \item Round 2+ alternates between requester and target agents.
  \item Reply exactly \texttt{REPLY\_SKIP} to stop the ping‑pong.
  \item Max turns is \texttt{session.agentToAgent.maxPingPongTurns} (0–5, default 5).
  \item Once the loop ends, OpenClaw runs the \textbf{agent‑to‑agent announce step} (target agent only):
  \item Reply exactly \texttt{ANNOUNCE\_SKIP} to stay silent.
  \item Any other reply is sent to the target channel.
  \item Announce step includes the original request + round‑1 reply + latest ping‑pong reply.
\end{itemize}


\subsection{Channel Field}


\begin{itemize}
  \item For groups, \texttt{channel} is the channel recorded on the session entry.
  \item For direct chats, \texttt{channel} maps from \texttt{lastChannel}.
  \item For cron/hook/node, \texttt{channel} is \texttt{internal}.
  \item If missing, \texttt{channel} is \texttt{unknown}.
\end{itemize}


\subsection{Security / Send Policy}


Policy-based blocking by channel/chat type (not per session id).

\begin{lstlisting}[language=json]
{
  "session": {
    "sendPolicy": {
      "rules": [
        {
          "match": { "channel": "discord", "chatType": "group" },
          "action": "deny"
        }
      ],
      "default": "allow"
    }
  }
}
\end{lstlisting}


Runtime override (per session entry):

\begin{itemize}
  \item \texttt{sendPolicy: "allow" | "deny"} (unset = inherit config)
  \item Settable via \texttt{sessions.patch} or owner-only \texttt{/send on|off|inherit} (standalone message).
\end{itemize}


Enforcement points:

\begin{itemize}
  \item \texttt{chat.send} / \texttt{agent} (gateway)
  \item auto-reply delivery logic
\end{itemize}


\subsection{sessions\_spawn}


Spawn a sub-agent run in an isolated session and announce the result back to the requester chat channel.

Parameters:

\begin{itemize}
  \item \texttt{task} (required)
  \item \texttt{label?} (optional; used for logs/UI)
  \item \texttt{agentId?} (optional; spawn under another agent id if allowed)
  \item \texttt{model?} (optional; overrides the sub-agent model; invalid values error)
  \item \texttt{thinking?} (optional; overrides thinking level for the sub-agent run)
  \item \texttt{runTimeoutSeconds?} (defaults to \texttt{agents.defaults.subagents.runTimeoutSeconds} when set, otherwise \texttt{0}; when set, aborts the sub-agent run after N seconds)
  \item \texttt{thread?} (default false; request thread-bound routing for this spawn when supported by the channel/plugin)
  \item \texttt{mode?} (\texttt{run|session}; defaults to \texttt{run}, but defaults to \texttt{session} when \texttt{thread=true}; \texttt{mode="session"} requires \texttt{thread=true})
  \item \texttt{cleanup?} (\texttt{delete|keep}, default \texttt{keep})
  \item \texttt{sandbox?} (\texttt{inherit|require}, default \texttt{inherit}; \texttt{require} rejects spawn unless the target child runtime is sandboxed)
  \item \texttt{attachments?} (optional array of inline files; subagent runtime only, ACP rejects). Each entry: \texttt{\{ name, content, encoding?: "utf8" | "base64", mimeType? \}}. Files are materialized into the child workspace at \texttt{.openclaw/attachments//}. Returns a receipt with sha256 per file.
  \item \texttt{attachAs?} (optional; \texttt{\{ mountPath? \}} hint reserved for future mount implementations)
\end{itemize}


Allowlist:

\begin{itemize}
  \item \texttt{agents.list[].subagents.allowAgents}: list of agent ids allowed via \texttt{agentId} (\texttt{["*"]} to allow any). Default: only the requester agent.
  \item Sandbox inheritance guard: if the requester session is sandboxed, \texttt{sessions\_spawn} rejects targets that would run unsandboxed.
\end{itemize}


Discovery:

\begin{itemize}
  \item Use \texttt{agents\_list} to discover which agent ids are allowed for \texttt{sessions\_spawn}.
\end{itemize}


Behavior:

\begin{itemize}
  \item Starts a new \texttt{agent::subagent:} session with \texttt{deliver: false}.
  \item Sub-agents default to the full tool set \textbf{minus session tools} (configurable via \texttt{tools.subagents.tools}).
  \item Sub-agents are not allowed to call \texttt{sessions\_spawn} (no sub-agent → sub-agent spawning).
  \item Always non-blocking: returns \texttt{\{ status: "accepted", runId, childSessionKey \}} immediately.
  \item With \texttt{thread=true}, channel plugins can bind delivery/routing to a thread target (Discord support is controlled by \texttt{session.threadBindings.\textit{} and \texttt{channels.discord.threadBindings.}}).
  \item After completion, OpenClaw runs a sub-agent \textbf{announce step} and posts the result to the requester chat channel.
  \item If the assistant final reply is empty, the latest \texttt{toolResult} from sub-agent history is included as \texttt{Result}.
  \item Reply exactly \texttt{ANNOUNCE\_SKIP} during the announce step to stay silent.
  \item Announce replies are normalized to \texttt{Status}/\texttt{Result}/\texttt{Notes}; \texttt{Status} comes from runtime outcome (not model text).
  \item Sub-agent sessions are auto-archived after \texttt{agents.defaults.subagents.archiveAfterMinutes} (default: 60).
  \item Announce replies include a stats line (runtime, tokens, sessionKey/sessionId, transcript path, and optional cost).
\end{itemize}


\subsection{Sandbox Session Visibility}


Session tools can be scoped to reduce cross-session access.

Default behavior:

\begin{itemize}
  \item \texttt{tools.sessions.visibility} defaults to \texttt{tree} (current session + spawned subagent sessions).
  \item For sandboxed sessions, \texttt{agents.defaults.sandbox.sessionToolsVisibility} can hard-clamp visibility.
\end{itemize}


Config:

\begin{lstlisting}[language=json]
{
  tools: {
    sessions: {
      // "self" | "tree" | "agent" | "all"
      // default: "tree"
      visibility: "tree",
    },
  },
  agents: {
    defaults: {
      sandbox: {
        // default: "spawned"
        sessionToolsVisibility: "spawned", // or "all"
      },
    },
  },
}
\end{lstlisting}


Notes:

\begin{itemize}
  \item \texttt{self}: only the current session key.
  \item \texttt{tree}: current session + sessions spawned by the current session.
  \item \texttt{agent}: any session belonging to the current agent id.
  \item \texttt{all}: any session (cross-agent access still requires \texttt{tools.agentToAgent}).
  \item When a session is sandboxed and \texttt{sessionToolsVisibility="spawned"}, OpenClaw clamps visibility to \texttt{tree} even if you set \texttt{tools.sessions.visibility="all"}.
\end{itemize}



\clearpage

% === Source file: concepts/session.md ===
\section{Session Management}


OpenClaw treats \textbf{one direct-chat session per agent} as primary. Direct chats collapse to \texttt{agent::} (default \texttt{main}), while group/channel chats get their own keys. \texttt{session.mainKey} is honored.

Use \texttt{session.dmScope} to control how \textbf{direct messages} are grouped:

\begin{itemize}
  \item \texttt{main} (default): all DMs share the main session for continuity.
  \item \texttt{per-peer}: isolate by sender id across channels.
  \item \texttt{per-channel-peer}: isolate by channel + sender (recommended for multi-user inboxes).
  \item \texttt{per-account-channel-peer}: isolate by account + channel + sender (recommended for multi-account inboxes).
\end{itemize}

Use \texttt{session.identityLinks} to map provider-prefixed peer ids to a canonical identity so the same person shares a DM session across channels when using \texttt{per-peer}, \texttt{per-channel-peer}, or \texttt{per-account-channel-peer}.

\subsection{Secure DM mode (recommended for multi-user setups)}


\begin{quote}
\textbf{Security Warning:} If your agent can receive DMs from \textbf{multiple people}, you should strongly consider enabling secure DM mode. Without it, all users share the same conversation context, which can leak private information between users.
\end{quote}


\textbf{Example of the problem with default settings:}

\begin{itemize}
  \item Alice (``) messages your agent about a private topic (for example, a medical appointment)
  \item Bob (``) messages your agent asking "What were we talking about?"
  \item Because both DMs share the same session, the model may answer Bob using Alice's prior context.
\end{itemize}


\textbf{The fix:} Set \texttt{dmScope} to isolate sessions per user:

\begin{lstlisting}[language=json]
// ~/.openclaw/openclaw.json
{
  session: {
    // Secure DM mode: isolate DM context per channel + sender.
    dmScope: "per-channel-peer",
  },
}
\end{lstlisting}


\textbf{When to enable this:}

\begin{itemize}
  \item You have pairing approvals for more than one sender
  \item You use a DM allowlist with multiple entries
  \item You set \texttt{dmPolicy: "open"}
  \item Multiple phone numbers or accounts can message your agent
\end{itemize}


Notes:

\begin{itemize}
  \item Default is \texttt{dmScope: "main"} for continuity (all DMs share the main session). This is fine for single-user setups.
  \item Local CLI onboarding writes \texttt{session.dmScope: "per-channel-peer"} by default when unset (existing explicit values are preserved).
  \item For multi-account inboxes on the same channel, prefer \texttt{per-account-channel-peer}.
  \item If the same person contacts you on multiple channels, use \texttt{session.identityLinks} to collapse their DM sessions into one canonical identity.
  \item You can verify your DM settings with \texttt{openclaw security audit} (see \href{/cli/security}{security}).
\end{itemize}


\subsection{Gateway is the source of truth}


All session state is \textbf{owned by the gateway} (the “master” OpenClaw). UI clients (macOS app, WebChat, etc.) must query the gateway for session lists and token counts instead of reading local files.

\begin{itemize}
  \item In \textbf{remote mode}, the session store you care about lives on the remote gateway host, not your Mac.
  \item Token counts shown in UIs come from the gateway’s store fields (\texttt{inputTokens}, \texttt{outputTokens}, \texttt{totalTokens}, \texttt{contextTokens}). Clients do not parse JSONL transcripts to “fix up” totals.
\end{itemize}


\subsection{Where state lives}


\begin{itemize}
  \item On the \textbf{gateway host}:
  \item Store file: \texttt{\textasciitilde{}/.openclaw/agents//sessions/sessions.json} (per agent).
  \item Transcripts: \texttt{\textasciitilde{}/.openclaw/agents//sessions/.jsonl} (Telegram topic sessions use \texttt{.../-topic-.jsonl}).
  \item The store is a map \texttt{sessionKey -> \{ sessionId, updatedAt, ... \}}. Deleting entries is safe; they are recreated on demand.
  \item Group entries may include \texttt{displayName}, \texttt{channel}, \texttt{subject}, \texttt{room}, and \texttt{space} to label sessions in UIs.
  \item Session entries include \texttt{origin} metadata (label + routing hints) so UIs can explain where a session came from.
  \item OpenClaw does \textbf{not} read legacy Pi/Tau session folders.
\end{itemize}


\subsection{Maintenance}


OpenClaw applies session-store maintenance to keep \texttt{sessions.json} and transcript artifacts bounded over time.

\subsubsection{Defaults}


\begin{itemize}
  \item \texttt{session.maintenance.mode}: \texttt{warn}
  \item \texttt{session.maintenance.pruneAfter}: \texttt{30d}
  \item \texttt{session.maintenance.maxEntries}: \texttt{500}
  \item \texttt{session.maintenance.rotateBytes}: \texttt{10mb}
  \item \texttt{session.maintenance.resetArchiveRetention}: defaults to \texttt{pruneAfter} (\texttt{30d})
  \item \texttt{session.maintenance.maxDiskBytes}: unset (disabled)
  \item \texttt{session.maintenance.highWaterBytes}: defaults to \texttt{80\%} of \texttt{maxDiskBytes} when budgeting is enabled
\end{itemize}


\subsubsection{How it works}


Maintenance runs during session-store writes, and you can trigger it on demand with \texttt{openclaw sessions cleanup}.

\begin{itemize}
  \item \texttt{mode: "warn"}: reports what would be evicted but does not mutate entries/transcripts.
  \item \texttt{mode: "enforce"}: applies cleanup in this order:
\end{itemize}

\begin{enumerate}
  \item prune stale entries older than \texttt{pruneAfter}
  \item cap entry count to \texttt{maxEntries} (oldest first)
  \item archive transcript files for removed entries that are no longer referenced
  \item purge old \texttt{\textit{.deleted.} and \texttt{}.reset.} archives by retention policy
  \item rotate \texttt{sessions.json} when it exceeds \texttt{rotateBytes}
  \item if \texttt{maxDiskBytes} is set, enforce disk budget toward \texttt{highWaterBytes} (oldest artifacts first, then oldest sessions)
\end{enumerate}


\subsubsection{Performance caveat for large stores}


Large session stores are common in high-volume setups. Maintenance work is write-path work, so very large stores can increase write latency.

What increases cost most:

\begin{itemize}
  \item very high \texttt{session.maintenance.maxEntries} values
  \item long \texttt{pruneAfter} windows that keep stale entries around
  \item many transcript/archive artifacts in \texttt{\textasciitilde{}/.openclaw/agents//sessions/}
  \item enabling disk budgets (\texttt{maxDiskBytes}) without reasonable pruning/cap limits
\end{itemize}


What to do:

\begin{itemize}
  \item use \texttt{mode: "enforce"} in production so growth is bounded automatically
  \item set both time and count limits (\texttt{pruneAfter} + \texttt{maxEntries}), not just one
  \item set \texttt{maxDiskBytes} + \texttt{highWaterBytes} for hard upper bounds in large deployments
  \item keep \texttt{highWaterBytes} meaningfully below \texttt{maxDiskBytes} (default is 80\%)
  \item run \texttt{openclaw sessions cleanup --dry-run --json} after config changes to verify projected impact before enforcing
  \item for frequent active sessions, pass \texttt{--active-key} when running manual cleanup
\end{itemize}


\subsubsection{Customize examples}


Use a conservative enforce policy:

\begin{lstlisting}[language=json]
{
  session: {
    maintenance: {
      mode: "enforce",
      pruneAfter: "45d",
      maxEntries: 800,
      rotateBytes: "20mb",
      resetArchiveRetention: "14d",
    },
  },
}
\end{lstlisting}


Enable a hard disk budget for the sessions directory:

\begin{lstlisting}[language=json]
{
  session: {
    maintenance: {
      mode: "enforce",
      maxDiskBytes: "1gb",
      highWaterBytes: "800mb",
    },
  },
}
\end{lstlisting}


Tune for larger installs (example):

\begin{lstlisting}[language=json]
{
  session: {
    maintenance: {
      mode: "enforce",
      pruneAfter: "14d",
      maxEntries: 2000,
      rotateBytes: "25mb",
      maxDiskBytes: "2gb",
      highWaterBytes: "1.6gb",
    },
  },
}
\end{lstlisting}


Preview or force maintenance from CLI:

\begin{lstlisting}[language=bash]
openclaw sessions cleanup --dry-run
openclaw sessions cleanup --enforce
\end{lstlisting}


\subsection{Session pruning}


OpenClaw trims \textbf{old tool results} from the in-memory context right before LLM calls by default.
This does \textbf{not} rewrite JSONL history. See \href{/concepts/session-pruning}{/concepts/session-pruning}.

\subsection{Pre-compaction memory flush}


When a session nears auto-compaction, OpenClaw can run a \textbf{silent memory flush}
turn that reminds the model to write durable notes to disk. This only runs when
the workspace is writable. See \href{/concepts/memory}{Memory} and
\href{/concepts/compaction}{Compaction}.

\subsection{Mapping transports → session keys}


\begin{itemize}
  \item Direct chats follow \texttt{session.dmScope} (default \texttt{main}).
  \item \texttt{main}: \texttt{agent::} (continuity across devices/channels).
  \item Multiple phone numbers and channels can map to the same agent main key; they act as transports into one conversation.
  \item \texttt{per-peer}: \texttt{agent::dm:}.
  \item \texttt{per-channel-peer}: \texttt{agent:::dm:}.
  \item \texttt{per-account-channel-peer}: \texttt{agent::::dm:} (accountId defaults to \texttt{default}).
  \item If \texttt{session.identityLinks} matches a provider-prefixed peer id (for example \texttt{telegram:123}), the canonical key replaces `` so the same person shares a session across channels.
  \item Group chats isolate state: \texttt{agent:::group:} (rooms/channels use \texttt{agent:::channel:}).
  \item Telegram forum topics append \texttt{:topic:} to the group id for isolation.
  \item Legacy \texttt{group:} keys are still recognized for migration.
  \item Inbound contexts may still use \texttt{group:}; the channel is inferred from \texttt{Provider} and normalized to the canonical \texttt{agent:::group:} form.
  \item Other sources:
  \item Cron jobs: \texttt{cron:}
  \item Webhooks: \texttt{hook:} (unless explicitly set by the hook)
  \item Node runs: \texttt{node-}
\end{itemize}


\subsection{Lifecycle}


\begin{itemize}
  \item Reset policy: sessions are reused until they expire, and expiry is evaluated on the next inbound message.
  \item Daily reset: defaults to \textbf{4:00 AM local time on the gateway host}. A session is stale once its last update is earlier than the most recent daily reset time.
  \item Idle reset (optional): \texttt{idleMinutes} adds a sliding idle window. When both daily and idle resets are configured, \textbf{whichever expires first} forces a new session.
  \item Legacy idle-only: if you set \texttt{session.idleMinutes} without any \texttt{session.reset}/\texttt{resetByType} config, OpenClaw stays in idle-only mode for backward compatibility.
  \item Per-type overrides (optional): \texttt{resetByType} lets you override the policy for \texttt{direct}, \texttt{group}, and \texttt{thread} sessions (thread = Slack/Discord threads, Telegram topics, Matrix threads when provided by the connector).
  \item Per-channel overrides (optional): \texttt{resetByChannel} overrides the reset policy for a channel (applies to all session types for that channel and takes precedence over \texttt{reset}/\texttt{resetByType}).
  \item Reset triggers: exact \texttt{/new} or \texttt{/reset} (plus any extras in \texttt{resetTriggers}) start a fresh session id and pass the remainder of the message through. \texttt{/new } accepts a model alias, \texttt{provider/model}, or provider name (fuzzy match) to set the new session model. If \texttt{/new} or \texttt{/reset} is sent alone, OpenClaw runs a short “hello” greeting turn to confirm the reset.
  \item Manual reset: delete specific keys from the store or remove the JSONL transcript; the next message recreates them.
  \item Isolated cron jobs always mint a fresh \texttt{sessionId} per run (no idle reuse).
\end{itemize}


\subsection{Send policy (optional)}


Block delivery for specific session types without listing individual ids.

\begin{lstlisting}[language=json]
{
  session: {
    sendPolicy: {
      rules: [
        { action: "deny", match: { channel: "discord", chatType: "group" } },
        { action: "deny", match: { keyPrefix: "cron:" } },
        // Match the raw session key (including the `agent:<id>:` prefix).
        { action: "deny", match: { rawKeyPrefix: "agent:main:discord:" } },
      ],
      default: "allow",
    },
  },
}
\end{lstlisting}


Runtime override (owner only):

\begin{itemize}
  \item \texttt{/send on} → allow for this session
  \item \texttt{/send off} → deny for this session
  \item \texttt{/send inherit} → clear override and use config rules
\end{itemize}

Send these as standalone messages so they register.

\subsection{Configuration (optional rename example)}


\begin{lstlisting}[language=json]
// ~/.openclaw/openclaw.json
{
  session: {
    scope: "per-sender", // keep group keys separate
    dmScope: "main", // DM continuity (set per-channel-peer/per-account-channel-peer for shared inboxes)
    identityLinks: {
      alice: ["telegram:123456789", "discord:987654321012345678"],
    },
    reset: {
      // Defaults: mode=daily, atHour=4 (gateway host local time).
      // If you also set idleMinutes, whichever expires first wins.
      mode: "daily",
      atHour: 4,
      idleMinutes: 120,
    },
    resetByType: {
      thread: { mode: "daily", atHour: 4 },
      direct: { mode: "idle", idleMinutes: 240 },
      group: { mode: "idle", idleMinutes: 120 },
    },
    resetByChannel: {
      discord: { mode: "idle", idleMinutes: 10080 },
    },
    resetTriggers: ["/new", "/reset"],
    store: "~/.openclaw/agents/{agentId}/sessions/sessions.json",
    mainKey: "main",
  },
}
\end{lstlisting}


\subsection{Inspecting}


\begin{itemize}
  \item \texttt{openclaw status} — shows store path and recent sessions.
  \item \texttt{openclaw sessions --json} — dumps every entry (filter with \texttt{--active }).
  \item \texttt{openclaw gateway call sessions.list --params '\{\}'} — fetch sessions from the running gateway (use \texttt{--url}/\texttt{--token} for remote gateway access).
  \item Send \texttt{/status} as a standalone message in chat to see whether the agent is reachable, how much of the session context is used, current thinking/verbose toggles, and when your WhatsApp web creds were last refreshed (helps spot relink needs).
  \item Send \texttt{/context list} or \texttt{/context detail} to see what’s in the system prompt and injected workspace files (and the biggest context contributors).
  \item Send \texttt{/stop} (or standalone abort phrases like \texttt{stop}, \texttt{stop action}, \texttt{stop run}, \texttt{stop openclaw}) to abort the current run, clear queued followups for that session, and stop any sub-agent runs spawned from it (the reply includes the stopped count).
  \item Send \texttt{/compact} (optional instructions) as a standalone message to summarize older context and free up window space. See \href{/concepts/compaction}{/concepts/compaction}.
  \item JSONL transcripts can be opened directly to review full turns.
\end{itemize}


\subsection{Tips}


\begin{itemize}
  \item Keep the primary key dedicated to 1:1 traffic; let groups keep their own keys.
  \item When automating cleanup, delete individual keys instead of the whole store to preserve context elsewhere.
\end{itemize}


\subsection{Session origin metadata}


Each session entry records where it came from (best-effort) in \texttt{origin}:

\begin{itemize}
  \item \texttt{label}: human label (resolved from conversation label + group subject/channel)
  \item \texttt{provider}: normalized channel id (including extensions)
  \item \texttt{from}/\texttt{to}: raw routing ids from the inbound envelope
  \item \texttt{accountId}: provider account id (when multi-account)
  \item \texttt{threadId}: thread/topic id when the channel supports it
\end{itemize}

The origin fields are populated for direct messages, channels, and groups. If a
connector only updates delivery routing (for example, to keep a DM main session
fresh), it should still provide inbound context so the session keeps its
explainer metadata. Extensions can do this by sending \texttt{ConversationLabel},
\texttt{GroupSubject}, \texttt{GroupChannel}, \texttt{GroupSpace}, and \texttt{SenderName} in the inbound
context and calling \texttt{recordSessionMetaFromInbound} (or passing the same context
to \texttt{updateLastRoute}).


\clearpage

% === Source file: concepts/streaming.md ===
\section{Streaming + chunking}


OpenClaw has two separate streaming layers:

\begin{itemize}
  \item \textbf{Block streaming (channels):} emit completed \textbf{blocks} as the assistant writes. These are normal channel messages (not token deltas).
  \item \textbf{Preview streaming (Telegram/Discord/Slack):} update a temporary \textbf{preview message} while generating.
\end{itemize}


There is \textbf{no true token-delta streaming} to channel messages today. Preview streaming is message-based (send + edits/appends).

\subsection{Block streaming (channel messages)}


Block streaming sends assistant output in coarse chunks as it becomes available.

\begin{lstlisting}
Model output
  └─ text_delta/events
       ├─ (blockStreamingBreak=text_end)
       │    └─ chunker emits blocks as buffer grows
       └─ (blockStreamingBreak=message_end)
            └─ chunker flushes at message_end
                   └─ channel send (block replies)
\end{lstlisting}


Legend:

\begin{itemize}
  \item \texttt{text\_delta/events}: model stream events (may be sparse for non-streaming models).
  \item \texttt{chunker}: \texttt{EmbeddedBlockChunker} applying min/max bounds + break preference.
  \item \texttt{channel send}: actual outbound messages (block replies).
\end{itemize}


\textbf{Controls:}

\begin{itemize}
  \item \texttt{agents.defaults.blockStreamingDefault}: \texttt{"on"}/\texttt{"off"} (default off).
  \item Channel overrides: \texttt{*.blockStreaming} (and per-account variants) to force \texttt{"on"}/\texttt{"off"} per channel.
  \item \texttt{agents.defaults.blockStreamingBreak}: \texttt{"text\_end"} or \texttt{"message\_end"}.
  \item \texttt{agents.defaults.blockStreamingChunk}: \texttt{\{ minChars, maxChars, breakPreference? \}}.
  \item \texttt{agents.defaults.blockStreamingCoalesce}: \texttt{\{ minChars?, maxChars?, idleMs? \}} (merge streamed blocks before send).
  \item Channel hard cap: \texttt{*.textChunkLimit} (e.g., \texttt{channels.whatsapp.textChunkLimit}).
  \item Channel chunk mode: \texttt{*.chunkMode} (\texttt{length} default, \texttt{newline} splits on blank lines (paragraph boundaries) before length chunking).
  \item Discord soft cap: \texttt{channels.discord.maxLinesPerMessage} (default 17) splits tall replies to avoid UI clipping.
\end{itemize}


\textbf{Boundary semantics:}

\begin{itemize}
  \item \texttt{text\_end}: stream blocks as soon as chunker emits; flush on each \texttt{text\_end}.
  \item \texttt{message\_end}: wait until assistant message finishes, then flush buffered output.
\end{itemize}


\texttt{message\_end} still uses the chunker if the buffered text exceeds \texttt{maxChars}, so it can emit multiple chunks at the end.

\subsection{Chunking algorithm (low/high bounds)}


Block chunking is implemented by \texttt{EmbeddedBlockChunker}:

\begin{itemize}
  \item \textbf{Low bound:} don’t emit until buffer >= \texttt{minChars} (unless forced).
  \item \textbf{High bound:} prefer splits before \texttt{maxChars}; if forced, split at \texttt{maxChars}.
  \item \textbf{Break preference:} \texttt{paragraph} → \texttt{newline} → \texttt{sentence} → \texttt{whitespace} → hard break.
  \item \textbf{Code fences:} never split inside fences; when forced at \texttt{maxChars}, close + reopen the fence to keep Markdown valid.
\end{itemize}


\texttt{maxChars} is clamped to the channel \texttt{textChunkLimit}, so you can’t exceed per-channel caps.

\subsection{Coalescing (merge streamed blocks)}


When block streaming is enabled, OpenClaw can \textbf{merge consecutive block chunks}
before sending them out. This reduces “single-line spam” while still providing
progressive output.

\begin{itemize}
  \item Coalescing waits for \textbf{idle gaps} (\texttt{idleMs}) before flushing.
  \item Buffers are capped by \texttt{maxChars} and will flush if they exceed it.
  \item \texttt{minChars} prevents tiny fragments from sending until enough text accumulates
\end{itemize}

(final flush always sends remaining text).
\begin{itemize}
  \item Joiner is derived from \texttt{blockStreamingChunk.breakPreference}
\end{itemize}

(\texttt{paragraph} → \texttt{\textbackslash\{\}n\textbackslash\{\}n}, \texttt{newline} → \texttt{\textbackslash\{\}n}, \texttt{sentence} → space).
\begin{itemize}
  \item Channel overrides are available via \texttt{*.blockStreamingCoalesce} (including per-account configs).
  \item Default coalesce \texttt{minChars} is bumped to 1500 for Signal/Slack/Discord unless overridden.
\end{itemize}


\subsection{Human-like pacing between blocks}


When block streaming is enabled, you can add a \textbf{randomized pause} between
block replies (after the first block). This makes multi-bubble responses feel
more natural.

\begin{itemize}
  \item Config: \texttt{agents.defaults.humanDelay} (override per agent via \texttt{agents.list[].humanDelay}).
  \item Modes: \texttt{off} (default), \texttt{natural} (800–2500ms), \texttt{custom} (\texttt{minMs}/\texttt{maxMs}).
  \item Applies only to \textbf{block replies}, not final replies or tool summaries.
\end{itemize}


\subsection{“Stream chunks or everything”}


This maps to:

\begin{itemize}
  \item \textbf{Stream chunks:} \texttt{blockStreamingDefault: "on"} + \texttt{blockStreamingBreak: "text\_end"} (emit as you go). Non-Telegram channels also need \texttt{*.blockStreaming: true}.
  \item \textbf{Stream everything at end:} \texttt{blockStreamingBreak: "message\_end"} (flush once, possibly multiple chunks if very long).
  \item \textbf{No block streaming:} \texttt{blockStreamingDefault: "off"} (only final reply).
\end{itemize}


\textbf{Channel note:} Block streaming is \textbf{off unless}
\texttt{*.blockStreaming} is explicitly set to \texttt{true}. Channels can stream a live preview
(\texttt{channels..streaming}) without block replies.

Config location reminder: the \texttt{blockStreaming*} defaults live under
\texttt{agents.defaults}, not the root config.

\subsection{Preview streaming modes}


Canonical key: \texttt{channels..streaming}

Modes:

\begin{itemize}
  \item \texttt{off}: disable preview streaming.
  \item \texttt{partial}: single preview that is replaced with latest text.
  \item \texttt{block}: preview updates in chunked/appended steps.
  \item \texttt{progress}: progress/status preview during generation, final answer at completion.
\end{itemize}


\subsubsection{Channel mapping}



\begin{table}[h]
\centering
\small
\begin{tabular}{|l|l|l|l|l|}
\hline
Channel & \texttt{off} & \texttt{partial} & \texttt{block} & \texttt{progress} \\
\hline
Telegram & ✅ & ✅ & ✅ & maps to \texttt{partial} \\
\hline
Discord & ✅ & ✅ & ✅ & maps to \texttt{partial} \\
\hline
Slack & ✅ & ✅ & ✅ & ✅ \\
\hline
\end{tabular}
\end{table}



Slack-only:

\begin{itemize}
  \item \texttt{channels.slack.nativeStreaming} toggles Slack native streaming API calls when \texttt{streaming=partial} (default: \texttt{true}).
\end{itemize}


Legacy key migration:

\begin{itemize}
  \item Telegram: \texttt{streamMode} + boolean \texttt{streaming} auto-migrate to \texttt{streaming} enum.
  \item Discord: \texttt{streamMode} + boolean \texttt{streaming} auto-migrate to \texttt{streaming} enum.
  \item Slack: \texttt{streamMode} auto-migrates to \texttt{streaming} enum; boolean \texttt{streaming} auto-migrates to \texttt{nativeStreaming}.
\end{itemize}


\subsubsection{Runtime behavior}


Telegram:

\begin{itemize}
  \item Uses Bot API \texttt{sendMessageDraft} in DMs when available, and \texttt{sendMessage} + \texttt{editMessageText} for group/topic preview updates.
  \item Preview streaming is skipped when Telegram block streaming is explicitly enabled (to avoid double-streaming).
  \item \texttt{/reasoning stream} can write reasoning to preview.
\end{itemize}


Discord:

\begin{itemize}
  \item Uses send + edit preview messages.
  \item \texttt{block} mode uses draft chunking (\texttt{draftChunk}).
  \item Preview streaming is skipped when Discord block streaming is explicitly enabled.
\end{itemize}


Slack:

\begin{itemize}
  \item \texttt{partial} can use Slack native streaming (\texttt{chat.startStream}/\texttt{append}/\texttt{stop}) when available.
  \item \texttt{block} uses append-style draft previews.
  \item \texttt{progress} uses status preview text, then final answer.
\end{itemize}



\clearpage

% === Source file: concepts/system-prompt.md ===
\section{System Prompt}


OpenClaw builds a custom system prompt for every agent run. The prompt is \textbf{OpenClaw-owned} and does not use the pi-coding-agent default prompt.

The prompt is assembled by OpenClaw and injected into each agent run.

\subsection{Structure}


The prompt is intentionally compact and uses fixed sections:

\begin{itemize}
  \item \textbf{Tooling}: current tool list + short descriptions.
  \item \textbf{Safety}: short guardrail reminder to avoid power-seeking behavior or bypassing oversight.
  \item \textbf{Skills} (when available): tells the model how to load skill instructions on demand.
  \item \textbf{OpenClaw Self-Update}: how to run \texttt{config.apply} and \texttt{update.run}.
  \item \textbf{Workspace}: working directory (\texttt{agents.defaults.workspace}).
  \item \textbf{Documentation}: local path to OpenClaw docs (repo or npm package) and when to read them.
  \item \textbf{Workspace Files (injected)}: indicates bootstrap files are included below.
  \item \textbf{Sandbox} (when enabled): indicates sandboxed runtime, sandbox paths, and whether elevated exec is available.
  \item \textbf{Current Date \& Time}: user-local time, timezone, and time format.
  \item \textbf{Reply Tags}: optional reply tag syntax for supported providers.
  \item \textbf{Heartbeats}: heartbeat prompt and ack behavior.
  \item \textbf{Runtime}: host, OS, node, model, repo root (when detected), thinking level (one line).
  \item \textbf{Reasoning}: current visibility level + /reasoning toggle hint.
\end{itemize}


Safety guardrails in the system prompt are advisory. They guide model behavior but do not enforce policy. Use tool policy, exec approvals, sandboxing, and channel allowlists for hard enforcement; operators can disable these by design.

\subsection{Prompt modes}


OpenClaw can render smaller system prompts for sub-agents. The runtime sets a
\texttt{promptMode} for each run (not a user-facing config):

\begin{itemize}
  \item \texttt{full} (default): includes all sections above.
  \item \texttt{minimal}: used for sub-agents; omits \textbf{Skills}, \textbf{Memory Recall}, **OpenClaw
\end{itemize}

Self-Update\textbf{, }Model Aliases\textbf{, }User Identity\textbf{, }Reply Tags**,
\textbf{Messaging}, \textbf{Silent Replies}, and \textbf{Heartbeats}. Tooling, \textbf{Safety},
Workspace, Sandbox, Current Date \& Time (when known), Runtime, and injected
context stay available.
\begin{itemize}
  \item \texttt{none}: returns only the base identity line.
\end{itemize}


When \texttt{promptMode=minimal}, extra injected prompts are labeled **Subagent
Context\textbf{ instead of }Group Chat Context**.

\subsection{Workspace bootstrap injection}


Bootstrap files are trimmed and appended under \textbf{Project Context} so the model sees identity and profile context without needing explicit reads:

\begin{itemize}
  \item \texttt{AGENTS.md}
  \item \texttt{SOUL.md}
  \item \texttt{TOOLS.md}
  \item \texttt{IDENTITY.md}
  \item \texttt{USER.md}
  \item \texttt{HEARTBEAT.md}
  \item \texttt{BOOTSTRAP.md} (only on brand-new workspaces)
  \item \texttt{MEMORY.md} and/or \texttt{memory.md} (when present in the workspace; either or both may be injected)
\end{itemize}


All of these files are \textbf{injected into the context window} on every turn, which
means they consume tokens. Keep them concise — especially \texttt{MEMORY.md}, which can
grow over time and lead to unexpectedly high context usage and more frequent
compaction.

\begin{quote}
\textbf{Note:} \texttt{memory/*.md} daily files are \textbf{not} injected automatically. They
are accessed on demand via the \texttt{memory\_search} and \texttt{memory\_get} tools, so they
do not count against the context window unless the model explicitly reads them.
\end{quote}


Large files are truncated with a marker. The max per-file size is controlled by
\texttt{agents.defaults.bootstrapMaxChars} (default: 20000). Total injected bootstrap
content across files is capped by \texttt{agents.defaults.bootstrapTotalMaxChars}
(default: 150000). Missing files inject a short missing-file marker.

Sub-agent sessions only inject \texttt{AGENTS.md} and \texttt{TOOLS.md} (other bootstrap files
are filtered out to keep the sub-agent context small).

Internal hooks can intercept this step via \texttt{agent:bootstrap} to mutate or replace
the injected bootstrap files (for example swapping \texttt{SOUL.md} for an alternate persona).

To inspect how much each injected file contributes (raw vs injected, truncation, plus tool schema overhead), use \texttt{/context list} or \texttt{/context detail}. See \href{/concepts/context}{Context}.

\subsection{Time handling}


The system prompt includes a dedicated \textbf{Current Date \& Time} section when the
user timezone is known. To keep the prompt cache-stable, it now only includes
the \textbf{time zone} (no dynamic clock or time format).

Use \texttt{session\_status} when the agent needs the current time; the status card
includes a timestamp line.

Configure with:

\begin{itemize}
  \item \texttt{agents.defaults.userTimezone}
  \item \texttt{agents.defaults.timeFormat} (\texttt{auto} | \texttt{12} | \texttt{24})
\end{itemize}


See \href{/date-time}{Date \& Time} for full behavior details.

\subsection{Skills}


When eligible skills exist, OpenClaw injects a compact \textbf{available skills list}
(\texttt{formatSkillsForPrompt}) that includes the \textbf{file path} for each skill. The
prompt instructs the model to use \texttt{read} to load the SKILL.md at the listed
location (workspace, managed, or bundled). If no skills are eligible, the
Skills section is omitted.

\begin{lstlisting}
<available_skills>
  <skill>
    <name>...</name>
    <description>...</description>
    <location>...</location>
  </skill>
</available_skills>
\end{lstlisting}


This keeps the base prompt small while still enabling targeted skill usage.

\subsection{Documentation}


When available, the system prompt includes a \textbf{Documentation} section that points to the
local OpenClaw docs directory (either \texttt{docs/} in the repo workspace or the bundled npm
package docs) and also notes the public mirror, source repo, community Discord, and
ClawHub (\href{https://clawhub.com}{https://clawhub.com}) for skills discovery. The prompt instructs the model to consult local docs first
for OpenClaw behavior, commands, configuration, or architecture, and to run
\texttt{openclaw status} itself when possible (asking the user only when it lacks access).


\clearpage

% === Source file: concepts/timezone.md ===
\section{Timezones}


OpenClaw standardizes timestamps so the model sees a \textbf{single reference time}.

\subsection{Message envelopes (local by default)}


Inbound messages are wrapped in an envelope like:

\begin{lstlisting}
[Provider ... 2026-01-05 16:26 PST] message text
\end{lstlisting}


The timestamp in the envelope is \textbf{host-local by default}, with minutes precision.

You can override this with:

\begin{lstlisting}[language=json]
{
  agents: {
    defaults: {
      envelopeTimezone: "local", // "utc" | "local" | "user" | IANA timezone
      envelopeTimestamp: "on", // "on" | "off"
      envelopeElapsed: "on", // "on" | "off"
    },
  },
}
\end{lstlisting}


\begin{itemize}
  \item \texttt{envelopeTimezone: "utc"} uses UTC.
  \item \texttt{envelopeTimezone: "user"} uses \texttt{agents.defaults.userTimezone} (falls back to host timezone).
  \item Use an explicit IANA timezone (e.g., \texttt{"Europe/Vienna"}) for a fixed offset.
  \item \texttt{envelopeTimestamp: "off"} removes absolute timestamps from envelope headers.
  \item \texttt{envelopeElapsed: "off"} removes elapsed time suffixes (the \texttt{+2m} style).
\end{itemize}


\subsubsection{Examples}


\textbf{Local (default):}

\begin{lstlisting}
[Signal Alice +1555 2026-01-18 00:19 PST] hello
\end{lstlisting}


\textbf{Fixed timezone:}

\begin{lstlisting}
[Signal Alice +1555 2026-01-18 06:19 GMT+1] hello
\end{lstlisting}


\textbf{Elapsed time:}

\begin{lstlisting}
[Signal Alice +1555 +2m 2026-01-18T05:19Z] follow-up
\end{lstlisting}


\subsection{Tool payloads (raw provider data + normalized fields)}


Tool calls (\texttt{channels.discord.readMessages}, \texttt{channels.slack.readMessages}, etc.) return \textbf{raw provider timestamps}.
We also attach normalized fields for consistency:

\begin{itemize}
  \item \texttt{timestampMs} (UTC epoch milliseconds)
  \item \texttt{timestampUtc} (ISO 8601 UTC string)
\end{itemize}


Raw provider fields are preserved.

\subsection{User timezone for the system prompt}


Set \texttt{agents.defaults.userTimezone} to tell the model the user's local time zone. If it is
unset, OpenClaw resolves the \textbf{host timezone at runtime} (no config write).

\begin{lstlisting}[language=json]
{
  agents: { defaults: { userTimezone: "America/Chicago" } },
}
\end{lstlisting}


The system prompt includes:

\begin{itemize}
  \item \texttt{Current Date \& Time} section with local time and timezone
  \item \texttt{Time format: 12-hour} or \texttt{24-hour}
\end{itemize}


You can control the prompt format with \texttt{agents.defaults.timeFormat} (\texttt{auto} | \texttt{12} | \texttt{24}).

See \href{/date-time}{Date \& Time} for the full behavior and examples.


\clearpage

% === Source file: concepts/typebox.md ===
\section{TypeBox as protocol source of truth}


Last updated: 2026-01-10

TypeBox is a TypeScript-first schema library. We use it to define the **Gateway
WebSocket protocol** (handshake, request/response, server events). Those schemas
drive \textbf{runtime validation}, \textbf{JSON Schema export}, and \textbf{Swift codegen} for
the macOS app. One source of truth; everything else is generated.

If you want the higher-level protocol context, start with
\href{/concepts/architecture}{Gateway architecture}.

\subsection{Mental model (30 seconds)}


Every Gateway WS message is one of three frames:

\begin{itemize}
  \item \textbf{Request}: \texttt{\{ type: "req", id, method, params \}}
  \item \textbf{Response}: \texttt{\{ type: "res", id, ok, payload | error \}}
  \item \textbf{Event}: \texttt{\{ type: "event", event, payload, seq?, stateVersion? \}}
\end{itemize}


The first frame \textbf{must} be a \texttt{connect} request. After that, clients can call
methods (e.g. \texttt{health}, \texttt{send}, \texttt{chat.send}) and subscribe to events (e.g.
\texttt{presence}, \texttt{tick}, \texttt{agent}).

Connection flow (minimal):

\begin{lstlisting}
Client                    Gateway
  |---- req:connect -------->|
  |<---- res:hello-ok --------|
  |<---- event:tick ----------|
  |---- req:health ---------->|
  |<---- res:health ----------|
\end{lstlisting}


Common methods + events:


\begin{table}[h]
\centering
\small
\begin{tabular}{|l|l|l|}
\hline
Category & Examples & Notes \\
\hline
Core & \texttt{connect}, \texttt{health}, \texttt{status} & \texttt{connect} must be first \\
\hline
Messaging & \texttt{send}, \texttt{poll}, \texttt{agent}, \texttt{agent.wait} & side-effects need \texttt{idempotencyKey} \\
\hline
Chat & \texttt{chat.history}, \texttt{chat.send}, \texttt{chat.abort}, \texttt{chat.inject} & WebChat uses these \\
\hline
Sessions & \texttt{sessions.list}, \texttt{sessions.patch}, \texttt{sessions.delete} & session admin \\
\hline
Nodes & \texttt{node.list}, \texttt{node.invoke}, \texttt{node.pair.*} & Gateway WS + node actions \\
\hline
Events & \texttt{tick}, \texttt{presence}, \texttt{agent}, \texttt{chat}, \texttt{health}, \texttt{shutdown} & server push \\
\hline
\end{tabular}
\end{table}



Authoritative list lives in \texttt{src/gateway/server.ts} (\texttt{METHODS}, \texttt{EVENTS}).

\subsection{Where the schemas live}


\begin{itemize}
  \item Source: \texttt{src/gateway/protocol/schema.ts}
  \item Runtime validators (AJV): \texttt{src/gateway/protocol/index.ts}
  \item Server handshake + method dispatch: \texttt{src/gateway/server.ts}
  \item Node client: \texttt{src/gateway/client.ts}
  \item Generated JSON Schema: \texttt{dist/protocol.schema.json}
  \item Generated Swift models: \texttt{apps/macos/Sources/OpenClawProtocol/GatewayModels.swift}
\end{itemize}


\subsection{Current pipeline}


\begin{itemize}
  \item \texttt{pnpm protocol:gen}
  \item writes JSON Schema (draft‑07) to \texttt{dist/protocol.schema.json}
  \item \texttt{pnpm protocol:gen:swift}
  \item generates Swift gateway models
  \item \texttt{pnpm protocol:check}
  \item runs both generators and verifies the output is committed
\end{itemize}


\subsection{How the schemas are used at runtime}


\begin{itemize}
  \item \textbf{Server side}: every inbound frame is validated with AJV. The handshake only
\end{itemize}

accepts a \texttt{connect} request whose params match \texttt{ConnectParams}.
\begin{itemize}
  \item \textbf{Client side}: the JS client validates event and response frames before
\end{itemize}

using them.
\begin{itemize}
  \item \textbf{Method surface}: the Gateway advertises the supported \texttt{methods} and
\end{itemize}

\texttt{events} in \texttt{hello-ok}.

\subsection{Example frames}


Connect (first message):

\begin{lstlisting}[language=json]
{
  "type": "req",
  "id": "c1",
  "method": "connect",
  "params": {
    "minProtocol": 2,
    "maxProtocol": 2,
    "client": {
      "id": "openclaw-macos",
      "displayName": "macos",
      "version": "1.0.0",
      "platform": "macos 15.1",
      "mode": "ui",
      "instanceId": "A1B2"
    }
  }
}
\end{lstlisting}


Hello-ok response:

\begin{lstlisting}[language=json]
{
  "type": "res",
  "id": "c1",
  "ok": true,
  "payload": {
    "type": "hello-ok",
    "protocol": 2,
    "server": { "version": "dev", "connId": "ws-1" },
    "features": { "methods": ["health"], "events": ["tick"] },
    "snapshot": {
      "presence": [],
      "health": {},
      "stateVersion": { "presence": 0, "health": 0 },
      "uptimeMs": 0
    },
    "policy": { "maxPayload": 1048576, "maxBufferedBytes": 1048576, "tickIntervalMs": 30000 }
  }
}
\end{lstlisting}


Request + response:

\begin{lstlisting}[language=json]
{ "type": "req", "id": "r1", "method": "health" }
\end{lstlisting}


\begin{lstlisting}[language=json]
{ "type": "res", "id": "r1", "ok": true, "payload": { "ok": true } }
\end{lstlisting}


Event:

\begin{lstlisting}[language=json]
{ "type": "event", "event": "tick", "payload": { "ts": 1730000000 }, "seq": 12 }
\end{lstlisting}


\subsection{Minimal client (Node.js)}


Smallest useful flow: connect + health.

\begin{lstlisting}[language=typescript]
import { WebSocket } from "ws";

const ws = new WebSocket("ws://127.0.0.1:18789");

ws.on("open", () => {
  ws.send(
    JSON.stringify({
      type: "req",
      id: "c1",
      method: "connect",
      params: {
        minProtocol: 3,
        maxProtocol: 3,
        client: {
          id: "cli",
          displayName: "example",
          version: "dev",
          platform: "node",
          mode: "cli",
        },
      },
    }),
  );
});

ws.on("message", (data) => {
  const msg = JSON.parse(String(data));
  if (msg.type === "res" && msg.id === "c1" && msg.ok) {
    ws.send(JSON.stringify({ type: "req", id: "h1", method: "health" }));
  }
  if (msg.type === "res" && msg.id === "h1") {
    console.log("health:", msg.payload);
    ws.close();
  }
});
\end{lstlisting}


\subsection{Worked example: add a method end‑to‑end}


Example: add a new \texttt{system.echo} request that returns \texttt{\{ ok: true, text \}}.

\begin{enumerate}
  \item \textbf{Schema (source of truth)}
\end{enumerate}


Add to \texttt{src/gateway/protocol/schema.ts}:

\begin{lstlisting}[language=typescript]
export const SystemEchoParamsSchema = Type.Object(
  { text: NonEmptyString },
  { additionalProperties: false },
);

export const SystemEchoResultSchema = Type.Object(
  { ok: Type.Boolean(), text: NonEmptyString },
  { additionalProperties: false },
);
\end{lstlisting}


Add both to \texttt{ProtocolSchemas} and export types:

\begin{lstlisting}[language=typescript]
  SystemEchoParams: SystemEchoParamsSchema,
  SystemEchoResult: SystemEchoResultSchema,
\end{lstlisting}


\begin{lstlisting}[language=typescript]
export type SystemEchoParams = Static<typeof SystemEchoParamsSchema>;
export type SystemEchoResult = Static<typeof SystemEchoResultSchema>;
\end{lstlisting}


\begin{enumerate}
  \item \textbf{Validation}
\end{enumerate}


In \texttt{src/gateway/protocol/index.ts}, export an AJV validator:

\begin{lstlisting}[language=typescript]
export const validateSystemEchoParams = ajv.compile<SystemEchoParams>(SystemEchoParamsSchema);
\end{lstlisting}


\begin{enumerate}
  \item \textbf{Server behavior}
\end{enumerate}


Add a handler in \texttt{src/gateway/server-methods/system.ts}:

\begin{lstlisting}[language=typescript]
export const systemHandlers: GatewayRequestHandlers = {
  "system.echo": ({ params, respond }) => {
    const text = String(params.text ?? "");
    respond(true, { ok: true, text });
  },
};
\end{lstlisting}


Register it in \texttt{src/gateway/server-methods.ts} (already merges \texttt{systemHandlers}),
then add \texttt{"system.echo"} to \texttt{METHODS} in \texttt{src/gateway/server.ts}.

\begin{enumerate}
  \item \textbf{Regenerate}
\end{enumerate}


\begin{lstlisting}[language=bash]
pnpm protocol:check
\end{lstlisting}


\begin{enumerate}
  \item \textbf{Tests + docs}
\end{enumerate}


Add a server test in \texttt{src/gateway/server.*.test.ts} and note the method in docs.

\subsection{Swift codegen behavior}


The Swift generator emits:

\begin{itemize}
  \item \texttt{GatewayFrame} enum with \texttt{req}, \texttt{res}, \texttt{event}, and \texttt{unknown} cases
  \item Strongly typed payload structs/enums
  \item \texttt{ErrorCode} values and \texttt{GATEWAY\_PROTOCOL\_VERSION}
\end{itemize}


Unknown frame types are preserved as raw payloads for forward compatibility.

\subsection{Versioning + compatibility}


\begin{itemize}
  \item \texttt{PROTOCOL\_VERSION} lives in \texttt{src/gateway/protocol/schema.ts}.
  \item Clients send \texttt{minProtocol} + \texttt{maxProtocol}; the server rejects mismatches.
  \item The Swift models keep unknown frame types to avoid breaking older clients.
\end{itemize}


\subsection{Schema patterns and conventions}


\begin{itemize}
  \item Most objects use \texttt{additionalProperties: false} for strict payloads.
  \item \texttt{NonEmptyString} is the default for IDs and method/event names.
  \item The top-level \texttt{GatewayFrame} uses a \textbf{discriminator} on \texttt{type}.
  \item Methods with side effects usually require an \texttt{idempotencyKey} in params
\end{itemize}

(example: \texttt{send}, \texttt{poll}, \texttt{agent}, \texttt{chat.send}).
\begin{itemize}
  \item \texttt{agent} accepts optional \texttt{internalEvents} for runtime-generated orchestration context
\end{itemize}

(for example subagent/cron task completion handoff); treat this as internal API surface.

\subsection{Live schema JSON}


Generated JSON Schema is in the repo at \texttt{dist/protocol.schema.json}. The
published raw file is typically available at:

\begin{itemize}
  \item \href{https://raw.githubusercontent.com/openclaw/openclaw/main/dist/protocol.schema.json}{https://raw.githubusercontent.com/openclaw/openclaw/main/dist/protocol.schema.json}
\end{itemize}


\subsection{When you change schemas}


\begin{enumerate}
  \item Update the TypeBox schemas.
  \item Run \texttt{pnpm protocol:check}.
  \item Commit the regenerated schema + Swift models.
\end{enumerate}



\clearpage

% === Source file: concepts/typing-indicators.md ===
\section{Typing indicators}


Typing indicators are sent to the chat channel while a run is active. Use
\texttt{agents.defaults.typingMode} to control \textbf{when} typing starts and \texttt{typingIntervalSeconds}
to control \textbf{how often} it refreshes.

\subsection{Defaults}


When \texttt{agents.defaults.typingMode} is \textbf{unset}, OpenClaw keeps the legacy behavior:

\begin{itemize}
  \item \textbf{Direct chats}: typing starts immediately once the model loop begins.
  \item \textbf{Group chats with a mention}: typing starts immediately.
  \item \textbf{Group chats without a mention}: typing starts only when message text begins streaming.
  \item \textbf{Heartbeat runs}: typing is disabled.
\end{itemize}


\subsection{Modes}


Set \texttt{agents.defaults.typingMode} to one of:

\begin{itemize}
  \item \texttt{never} — no typing indicator, ever.
  \item \texttt{instant} — start typing \textbf{as soon as the model loop begins}, even if the run
\end{itemize}

later returns only the silent reply token.
\begin{itemize}
  \item \texttt{thinking} — start typing on the \textbf{first reasoning delta} (requires
\end{itemize}

\texttt{reasoningLevel: "stream"} for the run).
\begin{itemize}
  \item \texttt{message} — start typing on the \textbf{first non-silent text delta} (ignores
\end{itemize}

the \texttt{NO\_REPLY} silent token).

Order of “how early it fires”:
\texttt{never} → \texttt{message} → \texttt{thinking} → \texttt{instant}

\subsection{Configuration}


\begin{lstlisting}[language=json]
{
  agent: {
    typingMode: "thinking",
    typingIntervalSeconds: 6,
  },
}
\end{lstlisting}


You can override mode or cadence per session:

\begin{lstlisting}[language=json]
{
  session: {
    typingMode: "message",
    typingIntervalSeconds: 4,
  },
}
\end{lstlisting}


\subsection{Notes}


\begin{itemize}
  \item \texttt{message} mode won’t show typing for silent-only replies (e.g. the \texttt{NO\_REPLY}
\end{itemize}

token used to suppress output).
\begin{itemize}
  \item \texttt{thinking} only fires if the run streams reasoning (\texttt{reasoningLevel: "stream"}).
\end{itemize}

If the model doesn’t emit reasoning deltas, typing won’t start.
\begin{itemize}
  \item Heartbeats never show typing, regardless of mode.
  \item \texttt{typingIntervalSeconds} controls the \textbf{refresh cadence}, not the start time.
\end{itemize}

The default is 6 seconds.


\clearpage

% === Source file: concepts/usage-tracking.md ===
\section{Usage tracking}


\subsection{What it is}


\begin{itemize}
  \item Pulls provider usage/quota directly from their usage endpoints.
  \item No estimated costs; only the provider-reported windows.
\end{itemize}


\subsection{Where it shows up}


\begin{itemize}
  \item \texttt{/status} in chats: emoji‑rich status card with session tokens + estimated cost (API key only). Provider usage shows for the \textbf{current model provider} when available.
  \item \texttt{/usage off|tokens|full} in chats: per-response usage footer (OAuth shows tokens only).
  \item \texttt{/usage cost} in chats: local cost summary aggregated from OpenClaw session logs.
  \item CLI: \texttt{openclaw status --usage} prints a full per-provider breakdown.
  \item CLI: \texttt{openclaw channels list} prints the same usage snapshot alongside provider config (use \texttt{--no-usage} to skip).
  \item macOS menu bar: “Usage” section under Context (only if available).
\end{itemize}


\subsection{Providers + credentials}


\begin{itemize}
  \item \textbf{Anthropic (Claude)}: OAuth tokens in auth profiles.
  \item \textbf{GitHub Copilot}: OAuth tokens in auth profiles.
  \item \textbf{Gemini CLI}: OAuth tokens in auth profiles.
  \item \textbf{Antigravity}: OAuth tokens in auth profiles.
  \item \textbf{OpenAI Codex}: OAuth tokens in auth profiles (accountId used when present).
  \item \textbf{MiniMax}: API key (coding plan key; \texttt{MINIMAX\_CODE\_PLAN\_KEY} or \texttt{MINIMAX\_API\_KEY}); uses the 5‑hour coding plan window.
  \item \textbf{z.ai}: API key via env/config/auth store.
\end{itemize}


Usage is hidden if no matching OAuth/API credentials exist.


\clearpage
